% Contraction de texte : reformuler les idées essentielles en vue d’un résumé — Français Tle Tech — Cours signé OnBuch+
% Programme officiel MINESEC. Compiler avec : tectonic N03-contraction-texte-reformuler-idees-essentielles-vue-res.tex
\newcommand{\DOCMATIERE}{Français}
\newcommand{\DOCNIVEAU}{Tle Tech}
\newcommand{\DOCMODULE}{Français – classe de Terminale technique}
\newcommand{\DOCLECON}{N03}
\newcommand{\DOCTITRE}{Contraction de texte : reformuler les idées essentielles en vue d’un résumé}
% OnBuch+ — préambule « cours signés » ENRICHI (compilé avec tectonic / XeLaTeX)
% Système visuel « Pop » d'OnBuch+ : fond crème (jamais blanc), bordures encre
% épaisses, ombres « dures » décalées, titres Archivo Black en capitales.
% Version MATHÉMATIQUES : graphes (pgfplots/tikz), boîtes pédagogiques
% (définition, propriété, méthode, attention, le savais-tu, exercice résolu,
% exercice, TP, bilan), page de garde et signature OnBuch+ ancrée en bas.
%
% Macros attendues dans main.tex AVANT \input{preamble} :
%   \DOCMATIERE  \DOCNIVEAU  \DOCMODULE  \DOCLECON (numéro)  \DOCTITRE
\documentclass[11pt]{article}

\usepackage{fontspec}
\usepackage[french]{babel}
\usepackage[a4paper,margin=2cm,top=3.4cm,bottom=2.8cm,headheight=1.4cm,footskip=1.3cm]{geometry}
\usepackage{amsmath,amssymb,amsfonts}
\usepackage{mathtools} % \xmapsto, \coloneqq... souvent utilisés par les modèles
\usepackage{cancel} % simplifications barrées (bilans de forces)
% \abs{z} (module, valeur absolue) et \norm{u} (norme), utilisés par les modèles
\providecommand{\abs}{}\let\abs\relax\DeclarePairedDelimiter{\abs}{\lvert}{\rvert}
\providecommand{\norm}{}\let\norm\relax\DeclarePairedDelimiter{\norm}{\lVert}{\rVert}
\usepackage[table]{xcolor} % [table] : fournit aussi \rowcolors (alternance de lignes)
\usepackage{pagecolor}
\usepackage{tikz}
\usetikzlibrary{arrows.meta,positioning,calc,shapes.geometric,shapes.misc,shapes.arrows,decorations.pathmorphing,decorations.pathreplacing,decorations.markings,patterns,shadows,mindmap,trees,fit,backgrounds,matrix,intersections,angles,quotes}
\usepackage{pgfplots}
\usepackage[version=4]{mhchem} % équations nucléaires \ce{^{235}_{92}U}
\usepackage[european,siunitx]{circuitikz} % schémas électriques
\pgfplotsset{compat=1.18}
\usepgfplotslibrary{fillbetween,groupplots} % aires entre courbes (calcul intégral)
\usepackage{enumitem}
\usepackage{fancyhdr}
% [most] charge toutes les bibliothèques tcolorbox (listings, minted, raster,
% documentation...) et gonfle inutilement la mémoire TeX sur un document long
% (~300 boîtes) : on ne charge que ce qui est réellement utilisé.
\usepackage[skins,breakable]{tcolorbox}
\usepackage{graphicx}
\usepackage{colortbl}
\usepackage{booktabs}
\usepackage{tabularx}
\usepackage{multirow}
\usepackage{diagbox} % case d'en-tête diagonale des tableaux à double entrée
% Colonnes extensibles centrée / gauche / droite (C, L, R), souvent utilisées par les modèles
\newcolumntype{C}{>{\centering\arraybackslash}X}
\newcolumntype{L}{>{\raggedright\arraybackslash}X}
\newcolumntype{R}{>{\raggedleft\arraybackslash}X}
\usepackage{array}
\usepackage{multicol}
\usepackage{siunitx}
% parse-numbers=false : accepte aussi \SI{2,5 \times 10^{-3}}{...}
\sisetup{locale=FR,output-decimal-marker={,},group-minimum-digits=4,parse-numbers=false}
\DeclareSIUnit\ohm{\ensuremath{\symup{\Omega}}} % U+2126 absent de la police
\DeclareSIUnit\division{div} % réglages d'oscilloscope (ms/div, V/div)
\DeclareSIUnit\tr{tr} % tours (vitesse de rotation, ex. \SI{1500}{\tr\per\minute})
\DeclareSIUnit\molar{M} % concentration molaire (ex. \SI{3}{\molar}, \SI{1}{\micro\molar})
\DeclareSIUnit\an{an} % année (le modèle confond parfois avec \year, un compteur TeX) — ex. \SI{5}{\kilo\meter\per\an}
\DeclareSIUnit\FCFA{FCFA} % franc CFA (ex. \SI{500}{\FCFA\per\kilo\gram})
\DeclareSIUnit\degreeBrix{\textdegree Brix} % teneur en sucre d'un sirop/jus (réfractométrie)
\DeclareSIUnit\month{mois} % le modèle utilise parfois \month, un compteur TeX, comme unité de durée
\DeclareSIUnit\microliter{\micro\liter} % le modèle écrit parfois \microliter en un seul token
\DeclareSIUnit\million{millions} % dénombrement (ex. \SI{15}{\million\per\mL} spermatozoïdes)
\usepackage{qrcode}
% \degree : commande fréquemment utilisée par les modèles pour un angle ou une
% température en dehors de \SI{}{} (non fournie par les paquets chargés).
\providecommand{\degree}{^{\circ}}
% \qed (fin de démonstration), utilisé par les modèles sans amsthm
\providecommand{\qed}{\hfill$\square$}
% \barre : double barre des valeurs interdites dans un tableau de variations (tkz-tab)
\providecommand{\barre}{\ensuremath{\|}}
% environnement proof (amsthm non chargé) utilisé par les modèles
\ifdefined\proof\else\newenvironment{proof}[1][Démonstration]{\par\noindent\textit{#1.}\ }{\qed\par}\fi
\usepackage{lastpage}
\usepackage{hyperref}
\hypersetup{hidelinks}

% ============================================================
% Palette « Pop » OnBuch+ — mêmes hex que la classe P (plus_ui.dart)
% ============================================================
\definecolor{poporange}{HTML}{F59321}
\definecolor{popdark}{HTML}{E07A0C}
\definecolor{popcream}{HTML}{FFF6E8}
\definecolor{popink}{HTML}{1C1714}
\definecolor{poppurple}{HTML}{7A5AE0}
\definecolor{popgreen}{HTML}{2E9E6A}
\definecolor{poppink}{HTML}{E0457B}
\definecolor{popblue}{HTML}{2F7DE1}
\definecolor{popgold}{HTML}{C98A12}
\definecolor{popmuted}{HTML}{8A7F76}
\definecolor{popline}{HTML}{EADFD2}
\definecolor{popgray}{HTML}{8A7F76}
\colorlet{popgrey}{popgray}
% Alias tolérants (le modèle nomme parfois les couleurs différemment)
\colorlet{popwhite}{white}
\colorlet{popblack}{popink}
\colorlet{popred}{poppink}
\colorlet{popyellow}{popgold}
% Teintes claires (fonds de tableaux/encadrés) — le modèle nomme parfois
% les couleurs avec un suffixe L (ex. popblueL) sans les avoir définies.
\colorlet{poporangeL}{poporange!20}
\colorlet{popdarkL}{popdark!20}
\colorlet{poppurpleL}{poppurple!20}
\colorlet{popgreenL}{popgreen!20}
\colorlet{poppinkL}{poppink!20}
\colorlet{popblueL}{popblue!20}
\colorlet{popgoldL}{popgold!20}
\colorlet{popredL}{popred!20}
\colorlet{popgrayL}{popgray!20}
% Teintes claires pour fonds de tableaux / zones
\colorlet{popblueL}{popblue!10!white}
\colorlet{poporangeL}{poporange!14!white}
\colorlet{popgreenL}{popgreen!12!white}
\colorlet{poppurpleL}{poppurple!12!white}
\colorlet{poppinkL}{poppink!12!white}
\colorlet{popgoldL}{popgold!14!white}

\pagecolor{popcream}

% ============================================================
% Typographie
% ============================================================
\defaultfontfeatures{Ligatures=TeX}
\setmainfont{PlusJakartaSans-Medium.ttf}[Path=../../../fonts/,BoldFont=PlusJakartaSans-Bold.ttf]
% Maths en sans-serif (Fira Math) avec chiffres et lettres droites de Plus
% Jakarta, pour que les formules se fondent dans le texte.
\usepackage[math-style=ISO,bold-style=ISO]{unicode-math}
\setmathfont{FiraMath-Regular.otf}[Path=../../../fonts/]
\setmathfont{PlusJakartaSans-Medium.ttf}[Path=../../../fonts/,range={up/num,up/{Latin,latin},\mathup}]
\usepackage{icomma} % virgule décimale sans espace en mode math
% Caractères mathématiques Unicode absents des polices de texte (Plus Jakarta,
% Archivo Black) : rendus par leur équivalent mathématique, en texte comme en
% formule — sinon ils s'affichent en carré vide (ex. « Arithmétique dans ℤ »).
\usepackage{newunicodechar}
\newunicodechar{ℝ}{\ensuremath{\mathbb{R}}}
\newunicodechar{ℤ}{\ensuremath{\mathbb{Z}}}
\newunicodechar{ℂ}{\ensuremath{\mathbb{C}}}
\newunicodechar{ℕ}{\ensuremath{\mathbb{N}}}
\newunicodechar{ℚ}{\ensuremath{\mathbb{Q}}}
\newunicodechar{𝔻}{\ensuremath{\mathbb{D}}}
\newunicodechar{α}{\ensuremath{\alpha}}
\newunicodechar{β}{\ensuremath{\beta}}
\newunicodechar{γ}{\ensuremath{\gamma}}
\newunicodechar{δ}{\ensuremath{\delta}}
\newunicodechar{ε}{\ensuremath{\varepsilon}}
\newunicodechar{θ}{\ensuremath{\theta}}
\newunicodechar{λ}{\ensuremath{\lambda}}
\newunicodechar{μ}{\ensuremath{\mu}}
\newunicodechar{σ}{\ensuremath{\sigma}}
\newunicodechar{τ}{\ensuremath{\tau}}
\newunicodechar{φ}{\ensuremath{\varphi}}
\newunicodechar{ψ}{\ensuremath{\psi}}
\newunicodechar{ω}{\ensuremath{\omega}}
\newunicodechar{Σ}{\ensuremath{\Sigma}}
% majuscules produites par \MakeUppercase dans les titres (α-aminés -> Α)
\newunicodechar{Α}{\ensuremath{\alpha}}
\newunicodechar{Β}{\ensuremath{\beta}}
\newunicodechar{Γ}{\ensuremath{\Gamma}}
\newunicodechar{Δ}{\ensuremath{\Delta}}
\newunicodechar{Ε}{\ensuremath{\varepsilon}}
\newunicodechar{Θ}{\ensuremath{\Theta}}
\newunicodechar{Λ}{\ensuremath{\Lambda}}
\newunicodechar{Μ}{\ensuremath{\mu}}
\newunicodechar{Τ}{\ensuremath{\tau}}
\newunicodechar{Φ}{\ensuremath{\Phi}}
\newunicodechar{Ψ}{\ensuremath{\Psi}}
\newunicodechar{Ω}{\ensuremath{\Omega}}
\newunicodechar{↦}{\ensuremath{\mapsto}}
\newunicodechar{∈}{\ensuremath{\in}}
\newunicodechar{∉}{\ensuremath{\notin}}
\newunicodechar{∧}{\ensuremath{\wedge}}
\newunicodechar{⇔}{\ensuremath{\Leftrightarrow}}
\newunicodechar{⟺}{\ensuremath{\Longleftrightarrow}}
\newunicodechar{⇒}{\ensuremath{\Rightarrow}}
\newunicodechar{⟹}{\ensuremath{\Longrightarrow}}
\newunicodechar{∘}{\ensuremath{\circ}}
\newunicodechar{∪}{\ensuremath{\cup}}
\newunicodechar{∩}{\ensuremath{\cap}}
\newunicodechar{≡}{\ensuremath{\equiv}}
\newunicodechar{⊂}{\ensuremath{\subset}}
\newunicodechar{⊥}{\ensuremath{\perp}}
\newunicodechar{∃}{\ensuremath{\exists}}
\newunicodechar{∀}{\ensuremath{\forall}}
\newunicodechar{∛}{\ensuremath{\sqrt[3]{\,}}}
\newunicodechar{∜}{\ensuremath{\sqrt[4]{\,}}}
\newunicodechar{⃗}{\ensuremath{{}^{\to}}}
\newunicodechar{ⁿ}{\ifmmode^{n}\else\textsuperscript{\ensuremath{n}}\fi}
\newunicodechar{ˣ}{\ifmmode^{x}\else\textsuperscript{\ensuremath{x}}\fi}
\newunicodechar{ᵇ}{\ifmmode^{b}\else\textsuperscript{\ensuremath{b}}\fi}
\newunicodechar{ʳ}{\ifmmode^{r}\else\textsuperscript{\ensuremath{r}}\fi}
\newunicodechar{ᵘ}{\ifmmode^{u}\else\textsuperscript{\ensuremath{u}}\fi}
\newunicodechar{ʸ}{\ifmmode^{y}\else\textsuperscript{\ensuremath{y}}\fi}
\newunicodechar{ᵃ}{\ifmmode^{a}\else\textsuperscript{\ensuremath{a}}\fi}
\newunicodechar{ᵉ}{\ifmmode^{e}\else\textsuperscript{\ensuremath{e}}\fi}
\newunicodechar{ᵏ}{\ifmmode^{k}\else\textsuperscript{\ensuremath{k}}\fi}
\newunicodechar{ᵖ}{\ifmmode^{p}\else\textsuperscript{\ensuremath{p}}\fi}
\newunicodechar{ᴺ}{\ifmmode^{N}\else\textsuperscript{\ensuremath{N}}\fi}
\newunicodechar{ᶜ}{\ifmmode^{c}\else\textsuperscript{\ensuremath{c}}\fi}
\newunicodechar{ʰ}{\ifmmode^{h}\else\textsuperscript{\ensuremath{h}}\fi}
\newunicodechar{ᵠ}{\ifmmode^{q}\else\textsuperscript{\ensuremath{q}}\fi}
\newunicodechar{ₙ}{\ifmmode_{n}\else\textsubscript{\ensuremath{n}}\fi}
\newunicodechar{ₐ}{\ifmmode_{a}\else\textsubscript{\ensuremath{a}}\fi}
\newunicodechar{ᵢ}{\ifmmode_{i}\else\textsubscript{\ensuremath{i}}\fi}
\newunicodechar{ᵣ}{\ifmmode_{r}\else\textsubscript{\ensuremath{r}}\fi}
\newunicodechar{ₖ}{\ifmmode_{k}\else\textsubscript{\ensuremath{k}}\fi}
\newunicodechar{ᵦ}{\ifmmode_{\beta}\else\textsubscript{\ensuremath{\beta}}\fi}
\newunicodechar{ₓ}{\ifmmode_{x}\else\textsubscript{\ensuremath{x}}\fi}
\newfontfamily\popdisplay{ArchivoBlack-Regular.ttf}[Path=../../../fonts/]
\newfontfamily\poplogo{SpaceGrotesk-Bold.ttf}[Path=../../../fonts/]
\newfontfamily\popsemi{PlusJakartaSans-SemiBold.ttf}[Path=../../../fonts/]
\newfontfamily\popxbold{PlusJakartaSans-ExtraBold.ttf}[Path=../../../fonts/]
\newfontfamily\popxboldls{PlusJakartaSans-ExtraBold.ttf}[Path=../../../fonts/,LetterSpace=3.0]

\setlist[enumerate]{leftmargin=1.7em,itemsep=5pt,topsep=4pt}
\setlist[itemize]{leftmargin=1.7em,itemsep=4pt,topsep=4pt,label={\color{poporange}\textbullet}}
\setlength{\parindent}{0pt}
\setlength{\parskip}{5pt}
\renewcommand{\arraystretch}{1.25}

% pgfplots : style Pop par défaut
\pgfplotsset{
  every axis/.append style={axis line style={popink,line width=1pt},tick style={popink},
    grid style={popline,line width=0.5pt},label style={font=\small},tick label style={font=\footnotesize}},
  popcurve/.style={poporange,line width=1.8pt,no markers,smooth},
}
% Même style disponible dans un \draw TikZ simple (hors pgfplots)
\tikzset{popcurve/.style={poporange,line width=1.8pt}}
\tikzset{popgrid/.style={popline,very thin}}
\colorlet{popcurve}{poporange} % le nom du style est parfois utilisé comme couleur

% ============================================================
% Pill / squiggle
% ============================================================
\newcommand{\poppill}[3]{%
  \tikz[baseline=-0.55ex]\node[draw=popink,line width=1.2pt,rounded corners=2.6pt,%
    fill=#2,inner xsep=6.5pt,inner ysep=3pt]{%
    \popxboldls\fontsize{7}{7}\selectfont\color{#3}\MakeUppercase{#1}};%
}
\newcommand{\popsquiggle}[1]{%
  \tikz[baseline=-0.4ex]{
    \draw[#1,line width=2.6pt,line cap=round]
      plot [smooth] coordinates {(0,0.09) (0.34,0.01) (0.68,0.21) (1.02,0.01) (1.36,0.10)};
  }%
}
% Mot-clé surligné (marqueur orange sous le texte)
\newcommand{\cle}[1]{{\popxbold\color{popdark}#1}}

% ============================================================
% En-tête / pied
% ============================================================
\pagestyle{fancy}
\fancyhf{}
\renewcommand{\headrulewidth}{0pt}
\renewcommand{\footrulewidth}{0pt}
\lhead{\raisebox{-0.15ex}{\poplogo\fontsize{13}{13}\selectfont\color{popink}OnBuch\color{poporange}+}}
\rhead{\poppill{\DOCMATIERE\ \textbullet\ \DOCNIVEAU\ \textbullet\ Leçon \DOCLECON}{white}{popink}}
\lfoot{\popxbold\fontsize{7.6}{7.6}\selectfont\color{popmuted}\MakeUppercase{Cours signé OnBuch+}\ \textbullet\ {\popsemi\DOCTITRE}}
\rfoot{\tikz[baseline=-0.6ex]\node[rounded corners=8.5pt,fill=popink,minimum height=17pt,minimum width=17pt,inner xsep=5pt,inner sep=0]{\popxbold\fontsize{8}{8}\selectfont\color{popcream}\thepage\,/\,\pageref*{LastPage}};}

% ============================================================
% Titres
% ============================================================
\newcommand{\chaptitre}[2]{%
  \begin{tcolorbox}[enhanced,colback=poporange,colframe=popink,boxrule=2.6pt,arc=11pt,
      drop shadow=popink,left=15pt,right=15pt,top=12pt,bottom=11pt,before skip=3pt,after skip=18pt]
    \poppill{#2}{popink}{popcream}\par\vspace{9pt}
    {\raggedright\hyphenpenalty=10000\exhyphenpenalty=10000\popdisplay\fontsize{22}{24}\selectfont\color{popcream}\MakeUppercase{#1}\par}
  \end{tcolorbox}
}
% Section numérotée : pastille orange avec le numéro + titre Archivo
\newcounter{popsec}
\newcommand{\coursec}[1]{%
  \stepcounter{popsec}\setcounter{popsub}{0}%
  \par\vspace{16pt}\noindent
  \begin{minipage}{\linewidth}
  \tikz[baseline=-0.7ex]\node[draw=popink,line width=1.6pt,fill=poporange,rounded corners=4pt,minimum size=22pt,inner sep=0]{\popdisplay\fontsize{12}{12}\selectfont\color{popcream}\thepopsec};%
  \hspace{8pt}{\popdisplay\fontsize{15}{17}\selectfont\color{popink}\MakeUppercase{#1}}\par
  \vspace{4pt}\noindent{\color{poporange}\rule{36pt}{3.6pt}}
  \end{minipage}\par\nopagebreak\vspace{8pt}
}
\newcounter{popsub}
\newcommand{\courssub}[1]{%
  \stepcounter{popsub}%
  \par\vspace{9pt}\noindent{\popxbold\fontsize{11.5}{13}\selectfont\color{popink}{\color{poporange}\thepopsec.\thepopsub}\hspace{6pt}#1}\par\nopagebreak\vspace{4pt}
}

% ============================================================
% Boîtes pédagogiques — même recette : fond blanc, bordure épaisse colorée,
% ombre dure encre, badge plein. Titre optionnel [#1].
% ============================================================
\tcbset{popbase/.style={breakable,enhanced,colback=white,boxrule=2.3pt,arc=9pt,
  shadow={3pt}{-3pt}{0pt}{popink},left=14pt,right=14pt,top=11pt,bottom=11pt,before skip=11pt,after skip=15pt,
  fontupper=\popsemi\fontsize{10.6}{14.5}\selectfont\color{popink},
  fontlower=\popsemi\fontsize{10.6}{14.5}\selectfont\color{popink}}}
% Coche dessinée (le glyphe ✓ n'existe pas dans les polices du document)
\newcommand{\popcheck}[1]{\tikz[baseline=-0.5ex]\draw[#1,line width=1.6pt,line cap=round,line join=round](-0.13,0.0)--(-0.04,-0.09)--(0.13,0.1);}
\providecommand{\coche}{\popcheck{popgreen}}
\providecommand{\resultat}[1]{\par\smallskip\noindent\fcolorbox{popgreen}{popgreenL}{\parbox{\dimexpr\linewidth-2\fboxsep-2\fboxrule\relax}{\textbf{Résultat.} #1}}\par\smallskip}
\newcommand{\popboxhead}[3]{\poppill{#1}{#2}{white}\ifx\relax#3\relax\else\hspace{6pt}{\popxbold\fontsize{10.6}{12}\selectfont\color{popink}#3}\fi\par\vspace{7pt}}

\newtcolorbox{aretenir}[1][]{popbase,colframe=popgreen,before upper={\popboxhead{À retenir}{popgreen}{#1}}}
\newtcolorbox{exemplebox}[1][]{popbase,colframe=poporange,before upper={\popboxhead{Exemple}{poporange}{#1}}}
\newtcolorbox{definition}[1][]{popbase,colframe=popblue,before upper={\popboxhead{Définition}{popblue}{#1}}}
\newtcolorbox{propriete}[1][]{popbase,colframe=poppurple,before upper={\popboxhead{Propriété}{poppurple}{#1}}}
\newtcolorbox{methode}[1][]{popbase,colframe=popgold,before upper={\popboxhead{Méthode}{popgold}{#1}}}
\newtcolorbox{attention}[1][]{popbase,colframe=poppink,before upper={\popboxhead{Attention}{poppink}{#1}}}
\newtcolorbox{experience}[1][]{popbase,colframe=popblue,colback=popblueL,before upper={\popboxhead{Expérience}{popblue}{#1}}}
% « Le savais-tu ? » : carte violette pleine (contexte, culture, Cameroun)
\newtcolorbox{savaistu}[1][]{popbase,colback=poppurple,colframe=popink,
  fontupper=\popsemi\fontsize{10.4}{14.2}\selectfont\color{white},
  before upper={\poppill{Le savais-tu\,?}{white}{poppurple}\ifx\relax#1\relax\else\hspace{6pt}{\popxbold\color{white}#1}\fi\par\vspace{7pt}}}
% Exercice résolu : énoncé en haut, \tcblower puis la solution sur fond crème
\newcounter{popexo}\newcounter{popexr}
\newtcolorbox{exoresolu}[1][]{popbase,colframe=poporange,colbacklower=popcream,
  segmentation style={popink,line width=1.2pt,dashed},
  before upper={\stepcounter{popexr}\popboxhead{Exercice résolu \thepopexr}{poporange}{#1}},
  before lower={\poppill{Solution}{popink}{popcream}\par\vspace{6pt}}}
% Exercice (non résolu en place) — la correction est dans la partie Corrigés
\newtcolorbox{exercice}[1][]{popbase,colframe=popink,
  before upper={\stepcounter{popexo}\popboxhead{Exercice \thepopexo}{popink}{#1}}}
% Corrigé
\newtcolorbox{corrige}[1][]{popbase,colframe=popgreen,colback=popgreenL,
  before upper={\popboxhead{Corrigé}{popgreen}{#1}}}
% Objectifs / prérequis / bilan
\newtcolorbox{objectifs}{popbase,colframe=popink,colback=poporangeL,
  before upper={\popboxhead{Objectifs de la leçon}{popink}{}}}
\newtcolorbox{prerequis}{popbase,colframe=popink,colback=popblueL,
  before upper={\popboxhead{Prérequis}{popblue}{}}}
\newtcolorbox{bilan}{popbase,colframe=popink,colback=white,boxrule=2.8pt,
  before upper={\popboxhead{Fiche bilan}{poporange}{L'essentiel à connaître pour l'examen}}}
% Figure encadrée avec légende
\newtcolorbox{popfigure}[1][]{enhanced,breakable=false,colback=white,colframe=popline,boxrule=1.4pt,arc=8pt,
  left=10pt,right=10pt,top=10pt,bottom=8pt,before skip=10pt,after skip=12pt,halign=center,
  lower separated=false,halign lower=center,
  fontlower=\popsemi\fontsize{9}{11.5}\selectfont\color{popmuted},
  #1}
\newcommand{\legende}[1]{\par\vspace{6pt}{\popsemi\fontsize{9}{11.5}\selectfont\color{popmuted}#1}}

% ============================================================
% Signature OnBuch+
% ============================================================
\newcommand{\poplogoinline}[1]{{\poplogo\fontsize{#1}{#1}\selectfont\color{popink}OnBuch\color{poporange}+}}
\newcommand{\signatureonbuch}{%
  \begin{tcolorbox}[enhanced,colback=popink,colframe=popink,boxrule=2.2pt,arc=10pt,
      drop shadow=poporange,left=14pt,right=14pt,top=10pt,bottom=10pt,before skip=0pt,after skip=0pt]
    \tikz[baseline=-0.7ex]\node[circle,fill=popgreen,draw=popcream,line width=1.3pt,minimum size=22pt,inner sep=0]{\popcheck{white}};%
    \hspace{9pt}\begin{minipage}[c]{0.62\linewidth}
      {\popxbold\fontsize{12}{13}\selectfont\color{popcream}Signé\ \textbullet\ {\poplogo OnBuch\color{poporange}+}}\par\vspace{2pt}
      {\raggedright\popsemi\fontsize{8}{10.5}\selectfont\color{popline}\MakeUppercase{Équipe pédagogique OnBuch+}\\\MakeUppercase{Conforme au programme officiel MINESEC}\par}
    \end{minipage}\hfill
    {\poplogo\fontsize{22}{22}\selectfont\color{popcream}OnBuch\color{poporange}+}
  \end{tcolorbox}%
}

% ============================================================
% Page de garde
% #1 titre · #2 matière · #3 niveau · #4 module · #5 n° leçon · #6 durée ·
% #7 description / objectifs (texte ou itemize)
% ============================================================
\newcommand{\pagegarde}[7]{%
  \thispagestyle{empty}
  \newgeometry{a4paper,margin=1.7cm,top=1.6cm,bottom=1.4cm}%
  \begin{tikzpicture}[remember picture,overlay]
    \fill[poporange!18] ([xshift=-3.2cm,yshift=-3.6cm]current page.north east) circle (4.6cm);
    \fill[poppurple!14] ([xshift=2.4cm,yshift=7.6cm]current page.south west) circle (3.8cm);
  \end{tikzpicture}%
  {\poplogo\fontsize{30}{30}\selectfont\color{popink}OnBuch\color{poporange}+}\hfill
  \raisebox{0.6ex}{\poppill{Cours signé \textbullet\ Premium}{popink}{popcream}}\par
  \vspace{1pt}\popsquiggle{poppurple}\par
  \vspace{8pt}
  \begin{tcolorbox}[enhanced,colback=poporange,colframe=popink,boxrule=2.8pt,arc=13pt,
      drop shadow=popink,left=18pt,right=18pt,top=12pt,bottom=13pt,after skip=10pt]
    \poppill{#2 \textbullet\ #3}{popink}{popcream}\hspace{5pt}\poppill{#4}{white}{popink}\par\vspace{8pt}
    {\popdisplay\fontsize{14}{14}\selectfont\color{popink}LEÇON #5}\par\vspace{4pt}
    {\raggedright\hyphenpenalty=10000\exhyphenpenalty=10000\popdisplay\fontsize{25}{27}\selectfont\color{popcream}\MakeUppercase{#1}\par}
    \vspace{8pt}
    {\popxbold\fontsize{9.5}{9.5}\selectfont\color{popink}\MakeUppercase{Durée conseillée : #6}}
  \end{tcolorbox}
  \begin{tcolorbox}[enhanced,colback=white,colframe=popink,boxrule=2.2pt,arc=10pt,
      drop shadow=popink,left=16pt,right=16pt,top=10pt,bottom=10pt,after skip=8pt]
    \poppill{Dans cette leçon}{poppurple}{white}\par\vspace{5pt}
    \popsemi\fontsize{9.3}{12.6}\selectfont\color{popink}#7
  \end{tcolorbox}
  \noindent\begin{minipage}[t]{0.487\linewidth}
    \begin{tcolorbox}[enhanced,colback=popgreen,colframe=popink,boxrule=2.2pt,arc=10pt,
        drop shadow=popink,left=11pt,right=11pt,top=10pt,bottom=10pt,height=3.1cm,valign=center]
      \tcbox[colback=white,colframe=popink,boxrule=1.3pt,arc=5pt,boxsep=0pt,nobeforeafter,
        left=3pt,right=3pt,top=3pt,bottom=3pt,valign=center]{\qrcode[height=1.9cm]{https://play.google.com/store/apps/details?id=com.luvvix.onbuch&hl=fr}}%
      \hspace{8pt}\begin{minipage}[c]{0.5\linewidth}
        {\popdisplay\fontsize{11}{12.5}\selectfont\color{white}\MakeUppercase{Télécharge OnBuch}}\par\vspace{4pt}
        {\raggedright\popsemi\fontsize{8.4}{11}\selectfont\color{white}Gratuit \textbullet\ Google Play\\Tuteur IA, annales, concours, résultats\par}
      \end{minipage}
    \end{tcolorbox}
  \end{minipage}\hfill
  \begin{minipage}[t]{0.487\linewidth}
    \begin{tcolorbox}[enhanced,colback=poppurple,colframe=popink,boxrule=2.2pt,arc=10pt,
        drop shadow=popink,left=11pt,right=11pt,top=10pt,bottom=10pt,height=3.1cm,valign=center]
      \tikz[baseline=-0.7ex]\node[circle,fill=white,draw=popink,line width=1.3pt,minimum size=1.6cm,inner sep=0]{\popdisplay\fontsize{24}{24}\selectfont\color{poppurple}+};%
      \hspace{8pt}\begin{minipage}[c]{0.6\linewidth}
        {\popdisplay\fontsize{11}{12.5}\selectfont\color{white}\MakeUppercase{Passe à OnBuch+}}\par\vspace{4pt}
        {\raggedright\popsemi\fontsize{8.4}{11}\selectfont\color{white}Tous les cours signés, Tuteur IA illimité, fiches PDF — onglet OnBuch+ de l'app\par}
      \end{minipage}
    \end{tcolorbox}
  \end{minipage}
  \vspace{8pt}
  \popcontenu
  \vfill
  \signatureonbuch
  \restoregeometry
  \newpage
}

% Rangée « Ce cours contient » (page de garde)
\newcommand{\poptile}[3]{% #1 couleur #2 gros texte #3 légende
  \begin{tcolorbox}[enhanced,colback=#1,colframe=popink,boxrule=2pt,arc=9pt,shadow={2.5pt}{-2.5pt}{0pt}{popink},
      left=6pt,right=6pt,top=9pt,bottom=9pt,halign=center,valign=center,height=1.75cm,width=\linewidth,nobeforeafter]
    {\popdisplay\fontsize{12.5}{13}\selectfont\color{white}\MakeUppercase{#2}}\par\vspace{3pt}
    {\centering\popsemi\fontsize{7.8}{9.5}\selectfont\color{white}#3\par}
  \end{tcolorbox}}
\newcommand{\popcontenu}{%
  \par\noindent{\popxboldls\fontsize{7.5}{7.5}\selectfont\color{popmuted}CE COURS SIGNÉ CONTIENT}\par\vspace{7pt}
  \noindent\begin{minipage}[t]{0.235\linewidth}\poptile{popblue}{Cours}{complet, illustré, pas à pas}\end{minipage}\hfill
  \begin{minipage}[t]{0.235\linewidth}\poptile{poporange}{Méthodes}{et exercices résolus}\end{minipage}\hfill
  \begin{minipage}[t]{0.235\linewidth}\poptile{poppink}{Exercices}{type examen + corrigés}\end{minipage}\hfill
  \begin{minipage}[t]{0.235\linewidth}\poptile{popgreen}{Fiche bilan}{l'essentiel à retenir}\end{minipage}\par}

% Sommaire
\newtcolorbox{sommaire}{enhanced,colback=white,colframe=popink,boxrule=2.2pt,arc=10pt,
  drop shadow=popink,left=18pt,right=18pt,top=13pt,bottom=13pt,before skip=6pt,after skip=20pt,
  before upper={\poppill{Sommaire}{poppurple}{white}\par\vspace{9pt}\popsemi\fontsize{10.6}{14.5}\selectfont\color{popink}}}

% Signature de fin de cours — poussée tout en bas de la dernière page
\newcommand{\findecours}{%
  \par\vfill\nopagebreak
  \begin{center}\popsquiggle{poporange}\end{center}\vspace{4pt}
  \signatureonbuch
}
\AtBeginDocument{\def\llbracket{\mathopen{[\mkern-3.5mu[}}\def\rrbracket{\mathclose{]\mkern-3.5mu]}}} % intervalles d'entiers (glyphe ⟦ absent de la police)
% Glyphes absents de Fira Math : redéfinis à partir de glyphes disponibles
\AtBeginDocument{%
  \def\setminus{\mathbin{\backslash}}\let\smallsetminus\setminus
  \def\vdots{\mathord{\vbox{\baselineskip3.2pt\lineskiplimit0pt\kern4pt\hbox{.}\hbox{.}\hbox{.}}}}%
  \def\ddots{\mathinner{\mkern1mu\raise6.5pt\hbox{.}\mkern2mu\raise3.6pt\hbox{.}\mkern2mu\raise0.7pt\hbox{.}\mkern1mu}}%
  \def\triangle{\mathord{\scalebox{0.9}{$\symup{\Delta}$}}}%
  \def\nabla{\mathord{\raisebox{\depth}{\scalebox{1}[-1]{$\symup{\Delta}$}}}}%
  \def\checkmark{\popcheck{popdark}}%
  \def\star{\mathbin{\ast}}\def\bigstar{\mathord{\ast}}%
  \def\diamond{\mathbin{\scalebox{0.6}{$\Diamond$}}}%
  \def\bigcup{\mathop{\mathchoice{\raisebox{-0.3em}{\scalebox{1.7}{$\cup$}}}{\scalebox{1.25}{$\cup$}}{\cup}{\cup}}\displaylimits}%
  \def\bigcap{\mathop{\mathchoice{\raisebox{-0.3em}{\scalebox{1.7}{$\cap$}}}{\scalebox{1.25}{$\cap$}}{\cap}{\cap}}\displaylimits}%
  \def\lesseqqgtr{\mathrel{\lessgtr}}\def\bigoplus{\mathop{\oplus}\displaylimits}%
}
\providecommand{\ssi}{si et seulement si} % abréviation fréquente
\AtBeginDocument{\providecommand{\wideparen}{\overparen}} % arc de cercle

\begin{document}

\pagegarde{Contraction de texte : reformuler les idées essentielles en vue d’un résumé}{Français}{Tle Tech}{Français – classe de Terminale technique}{N03}{15 séances (fiche de progression harmonisée)}{Cette leçon outille l'élève de Terminale technique pour dégager la structure argumentative d'un texte, en réduire la longueur par des techniques précises et reformuler les idées essentielles en vue d'un résumé fidèle et personnel. Elle s'appuie sur la lecture méthodique du Vieux nègre et la médaille de Ferdinand Oyono et mobilise des acquis de grammaire, de lexicologie et de stylistique.
\vspace{4pt}
\begin{multicols}{2}\small
\begin{itemize}
\item À la fin de cette leçon, je sais identifier la thèse, les arguments et les exemples dans un texte argumentatif
\item À la fin de cette leçon, je sais dégager la structure argumentative d'un texte (introduction, développement, conclusion) sous forme de plan schématique
\item À la fin de cette leçon, je sais appliquer les techniques de réduction (suppression, condensation, généralisation) pour respecter un nombre de mots imposé
\item À la fin de cette leçon, je sais reformuler les idées essentielles avec mes propres mots sans déformer la pensée de l'auteur
\item À la fin de cette leçon, je sais distinguer et employer les énoncés constatifs et les énoncés performatifs dans un texte
\item À la fin de cette leçon, je sais reconnaître l'ironie et le paradoxe et en rendre compte dans un résumé
\end{itemize}\end{multicols}}

\begin{sommaire}
\begin{multicols}{2}
\begin{enumerate}
\item La structure de la phrase : outil de la contraction
\item Énoncés constatifs et performatifs
\item Ironie et paradoxe : figures de pensée
\item Les variétés lexicales du français
\item Dégager la structure argumentative du texte
\item Techniques de réduction et reformulation des idées essentielles
\item Application au Vieux nègre et la médaille : lectures méthodiques 5 et 6, exposés
\item Activité d'intégration
\item Méthodes et exercices résolus
\item Exercices
\item Corrigés des exercices
\item Fiche bilan
\end{enumerate}
\end{multicols}
\end{sommaire}

\begin{objectifs}
\begin{itemize}
\item À la fin de cette leçon, je sais identifier la thèse, les arguments et les exemples dans un texte argumentatif
\item À la fin de cette leçon, je sais dégager la structure argumentative d'un texte (introduction, développement, conclusion) sous forme de plan schématique
\item À la fin de cette leçon, je sais appliquer les techniques de réduction (suppression, condensation, généralisation) pour respecter un nombre de mots imposé
\item À la fin de cette leçon, je sais reformuler les idées essentielles avec mes propres mots sans déformer la pensée de l'auteur
\item À la fin de cette leçon, je sais distinguer et employer les énoncés constatifs et les énoncés performatifs dans un texte
\item À la fin de cette leçon, je sais reconnaître l'ironie et le paradoxe et en rendre compte dans un résumé
\item À la fin de cette leçon, je sais analyser la structure de la phrase (simple, complexe) pour condenser efficacement
\item À la fin de cette leçon, je sais identifier les variétés lexicales du français (français standard, régional, camerounais) et choisir le registre adapté au résumé
\item À la fin de cette leçon, je sais rédiger, corriger et justifier un résumé conforme à une consigne de longueur
\end{itemize}
\end{objectifs}

\begin{prerequis}
\begin{itemize}
\item La phrase simple et la phrase complexe (classes de 4e et 3e)
\item Les connecteurs logiques et l'organisation du texte argumentatif (Première)
\item Les figures de style de base : comparaison, métaphore, hyperbole (Seconde)
\item La lecture méthodique d'un extrait de roman (séances précédentes sur Le Vieux nègre et la médaille)
\item Les notions de thème, de thèse et d'argument (Première)
\end{itemize}
\end{prerequis}

\newpage

\coursec{Leçon 3 : Contraction de texte : reformuler les idées essentielles en vue d'un résumé}

\courssub{Situation d'accroche et problématique}

Imaginez la scène : nous sommes à la veille du concours d'entrée à l'\textsc{enam}. Marlyse, élève de Terminale technique au Lycée technique de Yaoundé, est penchée sur sa copie. Elle doit rendre un résumé de 100 mots d'un article de 800 mots paru dans \emph{Cameroon Tribune} sur la décentralisation. Son premier essai dépasse 300 mots : elle a recopié des phrases entières, gardé des exemples illustratifs, des adjectifs qualificatifs et des subordonnées relatives qui alourdissent le texte sans en porter le sens. Le jury ne lira pas au-delà de la limite imposée. Marlyse a confondu \emph{recopier} et \emph{résumer}.

Comment repérer les idées essentielles d'un texte ? Comment les reformuler sans trahir l'auteur, en respectant la contrainte de longueur ? Cette leçon donne les techniques professionnelles de la contraction de texte, compétence exigée au Baccalauréat technique et dans la vie administrative camerounaise (rapports, notes de synthèse, comptes rendus). Le point de départ est la maîtrise de la syntaxe : comprendre la structure de la phrase pour oser la transformer.

\coursec{La structure de la phrase : outil de la contraction}

\courssub{Rappel : phrase simple et phrase complexe}

\begin{definition}[La proposition]
La \cle{proposition} est un ensemble de mots organisé autour d'un \textbf{verbe conjugué} (ou d'une forme verbale équivalente). C'est l'unité de base de la syntaxe. Une phrase contient au moins une proposition.
\end{definition}

\begin{exemplebox}[Phrase simple vs phrase complexe]
\textbf{Phrase simple (une seule proposition) :} « Le délégué du gouvernement a signé l'arrêté. » \\
\textbf{Phrase complexe (plusieurs propositions) :} « Le délégué du gouvernement a signé l'arrêté \textbf{que} le ministre \textbf{avait} préparé \textbf{afin que} la rentrée \textbf{se déroule} normalement. »
\end{exemplebox}

Dans la phrase complexe, on distingue la \cle{proposition principale} (qui peut exister seule) et les \cle{propositions subordonnées} (qui dépendent d'un mot de la principale, appelé \cle{antécédent} ou \cle{régissant}). Il existe aussi des \cle{propositions juxtaposées} (séparées par un point-virgule, deux-points ou virgule) et des \cle{propositions coordonnées} (liées par \emph{mais, ou, et, donc, or, ni, car}).

\begin{popfigure}
\begin{tikzpicture}[
    level distance=1.8cm,
    sibling distance=3.5cm,
    every node/.style={draw, rounded corners, fill=popblueL, thick, inner sep=4pt, font=\small},
    edge from parent/.style={draw, thick, -{Stealth[length=3mm]}, popdark},
    level 1/.style={sibling distance=4cm},
    level 2/.style={sibling distance=2.5cm}
]
\node {Phrase complexe}
    child { node {Proposition principale} 
        child { node {Sujet : Le délégué} }
        child { node {Verbe : a signé} }
        child { node {COD : l'arrêté} }
    }
    child { node {Subordonnée relative} 
        child { node {Antécédent : l'arrêté} }
        child { node {Pronom relatif : que} }
        child { node {Sujet : le ministre} }
        child { node {Verbe : avait préparé} }
    }
    child { node {Subordonnée de but} 
        child { node {Subordonnant : afin que} }
        child { node {Sujet : la rentrée} }
        child { node {Verbe : se déroule (subj.)} }
        child { node {CCM : normalement} }
    };
\end{tikzpicture}
\legende{Arbre syntaxique d'une phrase complexe : identification de la principale et des subordonnées (relative, de but). La principale porte l'information centrale ; les subordonnées apportent des précisions (définition, cause, but, conséquence...).}
\end{popfigure}

\courssub{Typologie des propositions : repérage dans un texte}

Pour contracter, il faut d'abord \textbf{segmenter} le texte en propositions et \textbf{hiérarchiser} l'information.

\begin{center}
\begin{tabularx}{\linewidth}{|l|X|X|l|}
\hline
\rowcolor{popblueL}
\textbf{Type} & \textbf{Lien / Marqueur} & \textbf{Fonction sémantique} & \textbf{Exemple (Cameroon Tribune)} \\
\hline
Indépendante & Ponctuation forte (.) & Affirmation autonomes & Le Ministre a présidé la cérémonie. \\
\hline
Juxtaposée & ; : , & Parataxe (égalité) & Le Ministre est arrivé ; la cérémonie a commencé. \\
\hline
Coordonnée & \emph{mais, ou, et, donc, or, ni, car} & Relation logique explicite & Le Ministre a parlé \textbf{et} les élèves ont applaudi. \\
\hline
Subordonnée relative & \emph{qui, que, dont, où, lequel...} & Définition / Précision & Le texte \textbf{que} vous lisez est officiel. \\
\hline
Subordonnée conjonctive & \emph{que, si, quand, comme, puisque, afin que...} & Circonstance (cause, temps, but...) & \textbf{Parce que} la loi l'exige, nous obéissons. \\
\hline
Subordonnée participiale & Participe présent / passé & Circonstance condensée & \textbf{Signé} le décret, le Ministre est parti. \\
\hline
\end{tabularx}
\end{center}

\begin{methode}[Repérer la structure argumentative d'un paragraphe]
\begin{enumerate}
\item \textbf{Segmenter :} souligner tous les verbes conjugués $\rightarrow$ compter les propositions.
\item \textbf{Identifier la principale :} celle qui reste si on supprime les autres (test de l'autonomie).
\item \textbf{Classifier les subordonnées :} relative (précision), conjonctive (circonstance : cause, but, conséquence, condition), participiale (condensation native).
\item \textbf{Évaluer le poids informationnel :} une subordonnée de but (\emph{afin que}) ou de cause (\emph{parce que}) porte souvent l'idée force ; une relative explicative (\emph{qui était...}) est souvent accessoire.
\end{enumerate}
\end{methode}

\courssub{La subordination : hiérarchie de l'information}

La subordination n'est pas qu'une contrainte grammaticale : c'est un \textbf{choix rhétorique}. L'auteur place en principale ce qu'il veut imposer comme fait majeur ; il relègue en subordonnée ce qui est présupposé, accessoire ou justificatif.

\begin{propriete}[Hiérarchie implicite]
\begin{itemize}
\item \textbf{Principale} = Information nouvelle, thèse, fait saillant.
\item \textbf{Subordonnée conjonctive (cause, concession, condition)} = Argumentation, justification.
\item \textbf{Subordonnée relative (restrictive)} = Identification nécessaire de l'antécédent.
\item \textbf{Subordonnée relative (explicative / épithète)} = Information accessoire, souvent entre virgules $\rightarrow$ \textbf{cible prioritaire de suppression}.
\item \textbf{Participiale} = Déjà condensée ; à conserver ou nominaliser.
\end{itemize}
\end{propriete}

\begin{exemplebox}[Effet de mise en relief par la subordination]
\textit{Version A (Cause en principale) :} « \textbf{La décentralisation a échoué} parce que les fonds n'ont pas été transférés. » (Thèse : l'échec).
\textit{Version B (Cause en subordonnée) :} « \textbf{Parce que} les fonds n'ont pas été transférés, la décentralisation a échoué. » (Thèse : le non-transfert des fonds).
\textbf{Conséquence pour le résumé :} Selon la question posée (« Pourquoi l'échec ? » vs « Quel constat ? »), on gardera la principale correspondante.
\end{exemplebox}

\courssub{Expansions du nom et compléments circonstanciels : ce qui est supprimable}

La \cle{contraction} (ou \cle{réduction}) vise à passer d'une phrase complexe à une phrase simple, voire à un groupe nominal, sans perte de sens \emph{essentiel}. Les cibles naturelles sont les \textbf{expansions du nom} et les \textbf{compléments circonstanciels} (CC) non obligatoires.

\begin{center}
\begin{tabularx}{\linewidth}{|l|X|X|}
\hline
\rowcolor{popblueL}
\textbf{Élément syntaxique} & \textbf{Rôle} & \textbf{Stratégie de contraction} \\
\hline
Adjectif qualificatif / Apposition & Qualification & Supprimer si redondant ou inférentiel (ex: \emph{le Ministre \textbf{actuel}} $\rightarrow$ le Ministre) \\
\hline
Complément du nom (prep. \emph{de}) & Précision / Possession & Conserver si identifiant (ex: \emph{la signature \textbf{du Ministre}}) ; supprimer si générique \\
\hline
Relative explicative (\emph{qui, lequel...} + virgules) & Info accessoire & \textbf{Supprimer purement et simplement} \\
\hline
CC Lieu / Temps / Manière & Circonstance & Supprimer si le contexte le rend évident ; garder si datation cruciale \\
\hline
CC Cause / But / Condition & Argumentation & \textbf{Ne jamais supprimer sans les nominaliser} (voir piège) \\
\hline
\end{tabularx}
\end{center}

\begin{attention}[Piège majeur : supprimer l'idée force]
Ne supprimez \textbf{jamais} une subordonnée de cause (\emph{parce que, comme, puisque}), de but (\emph{afin que, pour que}) ou de condition (\emph{si}) sans en avoir récupéré le contenu sous forme nominale.
\begin{itemize}
\item \textbf{Texte :} « Le projet a été rejeté \textbf{parce qu'il était trop coûteux}. »
\item \textbf{Mauvais résumé :} « Le projet a été rejeté. » (Perte de l'explication $\rightarrow$ hors sujet).
\item \textbf{Bon résumé :} « Le projet a été rejeté \textbf{en raison de son coût excessif} » / « \textbf{Le coût excessif} a causé le rejet du projet. »
\end{itemize}
\end{attention}

\courssub{Transformation : de la phrase complexe à la phrase simple}

C'est le cœur technique de la contraction. Trois opérations majeures s'articulent :

\begin{enumerate}
\item \textbf{Suppression} des expansions accessoires (adjectifs, relatives explicatives, CC lieu/temps redondants).
\item \textbf{Nominalisation} : transformer un verbe (processus) en nom (chose/résultat). C'est l'outil n°1 de la condensation administrative.
\item \textbf{Intégration (incorporation)} : transformer une subordonnée en complément du nom (préposition + nom) ou en participe.
\end{enumerate}

\begin{methode}[Algorithme de condensation d'une phrase]
\begin{enumerate}
\item \textbf{Isoler} le noyau : Sujet + Verbe + Compléments \textbf{obligatoires} (COD, COI, COS).
\item \textbf{Lister} les propositions subordonnées et expansions.
\item \textbf{Trier} : 
    \begin{itemize}
    \item \textsc{Accessoire} (relatives explicatives, CC lieu/temps connus, adjectifs génériques) $\rightarrow$ \textbf{Supprimer}.
    \item \textsc{Argumentatif} (cause, but, concession, condition, relative restrictive) $\rightarrow$ \textbf{Nominaliser / Intégrer}.
    \end{itemize}
\item \textbf{Réécrire} la phrase simple en vérifiant la correction syntaxique (accords, prépositions).
\item \textbf{Contrôler} : le sens est-il préservé ? La phrase est-elle plus courte ? Est-elle fluide ?
\end{enumerate}
\end{methode}

\begin{exemplebox}[Exemple chiffré : condensation progressive]
\textbf{Original (26 mots) :} « Le vieux Meka, \textbf{qui était épuisé par une longue attente sous le soleil brûlant}, s'endormit \textbf{profondément} sur sa chaise \textbf{en bois}. »

\begin{enumerate}
\item \textit{Suppression relative explicative + CC lieu générique + adjectif « profondément » + « en bois » :} « Le vieux Meka, épuisé par une longue attente sous le soleil, s'endormit sur sa chaise. » (16 mots)
\item \textit{Nominalisation « était épuisé par » $\rightarrow$ « épuisé par » (participe) / « à cause de » :} « Épuisé par une longue attente sous le soleil, le vieux Meka s'endormit sur sa chaise. » (14 mots)
\item \textit{Condensation maximale (contexte connu = soleil/attente) :} « \textbf{Épuisé, Meka s'endormit.} » (\textbf{4 mots})
\end{enumerate}
\textbf{Taux de compression :} 84 \%. L'idée essentielle (État $\rightarrow$ Action) est intacte.
\end{exemplebox}

\begin{popfigure}
\begin{tikzpicture}[
    node distance=1.2cm and 2cm,
    box/.style={draw, rounded corners, thick, inner sep=6pt, minimum width=2.5cm, align=center, font=\small},
    arrow/.style={-Stealth, thick, popdark},
    main/.style={box, fill=popblueL},
    trans/.style={box, fill=poporangeL},
    final/.style={box, fill=popgreenL},
]
\node[main, align=center] (orig){Phrase complexe\\« Le vieux Meka, qui était\\épuisé par une longue attente\\sous le soleil, s'endormit\\profondément sur sa chaise en bois. »};
\node[trans, below=of orig, align=center] (step1){Étape 1 : Élagage\\« Le vieux Meka, épuisé par\\une longue attente sous le soleil,\\s'endormit sur sa chaise. »};
\node[trans, below=of step1, align=center] (step2){Étape 2 : Nominalisation\\« Épuisé par une longue attente,\\Meka s'endormit. »};
\node[final, below=of step2, align=center] (final){Phrase simple condensée\\« \textbf{Épuisé, Meka s'endormit.} »};
\draw[arrow] (orig) -- (step1) node[midway, right, font=\tiny, text=popdark, align=center]{Supprimer : relative expl., CC lieu,\\adjectifs, déterminants};
\draw[arrow] (step1) -- (step2) node[midway, right, font=\tiny, text=popdark] {Verbe $\rightarrow$ Participe / Nom};
\draw[arrow] (step2) -- (final) node[midway, right, font=\tiny, text=popdark] {Contexte $\rightarrow$ Ellipse};
\end{tikzpicture}
\legende{Schéma de condensation : de la phrase complexe (26 mots) à la phrase simple noyau (4 mots) par élagage syntaxique, nominalisation et ellipse contextuelle.}
\end{popfigure}

\courssub{La nominalisation : technique reine du résumé administratif}

\begin{savaistu}[La nominalisation dans l'administration camerounaise]
Dans les rapports de la \textsc{fonction publique}, les comptes rendus de Conseil des Ministres ou les notes de synthèse de l'\textsc{enam}, la nominalisation est systématique. Elle permet de transformer des \textbf{processus} (verbes d'action) en \textbf{objets} (noms de concepts) manipulables, classables, chiffrables.
\begin{itemize}
\item Verbe \emph{décider} $\rightarrow$ Nom \emph{la décision} (objet juridique).
\item Verbe \emph{réunir} $\rightarrow$ Nom \emph{la réunion} (objet agenda).
\item Verbe \emph{résister} $\rightarrow$ Nom \emph{la résistance} (concept politique/historique).
\item Verbe \emph{décentraliser} $\rightarrow$ Nom \emph{la décentralisation} (politique publique).
\end{itemize}
Cela permet d'écrire : « \textbf{La décentralisation} bute sur \textbf{l'insuffisance des transferts} » au lieu de « \textbf{On décentralise} mais \textbf{l'État ne transfère pas assez} ». Le style devient \textbf{impersonnel, dense, objectif} — le style de l'expertise technique.
\end{savaistu}

\begin{exoresolu}[Condensation d'un extrait de Cameroon Tribune]
\textbf{Consigne :} Condenser les trois phrases suivantes en \textbf{une seule phrase correcte} (max 25 mots).

\textit{Extrait :} « Le Ministre de la Décentralisation et du Développement Local a présidé hier la cérémonie d'installation des nouveaux maires. \textbf{Il a rappelé} que la loi leur confie la gestion des affaires locales. \textbf{Il les a exhortés} à faire preuve de transparence dans la gestion des finances communales. »

\textbf{Analyse structurelle :}
\begin{itemize}
\item P1 (Principale 1) : Le Ministre a présidé la cérémonie.
\item P2 (Sub. conjonctive objet de « rappelé ») : que la loi leur confie la gestion...
\item P3 (Principale 2, coordonnée à P1 par juxtaposition/paragraphe) : Il les a exhortés à faire preuve...
\item P4 (Sub. infinitif objet de « exhortés ») : à faire preuve de transparence...
\end{itemize}

\textbf{Application de la méthode :}
\begin{enumerate}
\item \textbf{Noyaux} : Ministre / présider / cérémonie ; Ministre / rappeler / loi / gestion ; Ministre / exhorter / maires / transparence / finances.
\item \textbf{Suppressions} : « hier » (contexte), « nouveaux » (inféré par installation), « dans la gestion des finances communales » (précisé par « transparence » + « communes »).
\item \textbf{Nominalisations} : 
    \begin{itemize}
    \item « a présidé la cérémonie d'installation » $\rightarrow$ « a installé » (verbe) ou « l'installation » (nom).
    \item « a rappelé que la loi leur confie » $\rightarrow$ « a rappelé leur mandat légal » / « le mandat légal ».
    \item « a exhortés à faire preuve de transparence » $\rightarrow$ « a exigé la transparence » / « l'exigence de transparence ».
    \end{itemize}
\item \textbf{Fusion} : Le sujet unique (Le Ministre) permet la coordination des verbes ou la nominalisation globale.
\end{enumerate}

\textbf{Propositions de réponses (choisir selon la nuance voulue) :}
\begin{enumerate}
\item \textit{Verbal (action) :} « Le Ministre a installé les maires, rappelé leur mandat légal et exigé une gestion transparente. » (16 mots)
\item \textit{Nominalisé (style rapport) :} « Lors de l'installation des maires, le Ministre a rappelé leur mandat légal et exigé la transparence financière. » (17 mots)
\item \textit{Très condensé (titre/lead) :} « Installation des maires : le Ministre impose mandat légal et transparence financière. » (11 mots)
\end{enumerate}

\textbf{Vérification :} Sens conservé ? Oui. Structure argumentative (Installation $\rightarrow$ Rappel cadre $\rightarrow$ Exigence éthique) respectée ? Oui. Syntaxe correcte ? Oui.
\tcblower
\textbf{Solution détaillée :} La réponse 2 est la plus équilibrée pour un résumé administratif : elle conserve la chronologie (circonstance temporelle « Lors de... »), le sujet agent (le Ministre), et les deux idées forces (mandat légal, transparence) sous forme nominale dense. La réponse 1 est plus narrative. La réponse 3 est un titre, pas une phrase de résumé.
\end{exoresolu}

\coursec{Énoncés constatifs et performatifs}

\courssub{Distinction fondamentale (Théorie des actes de langage : Austin, Searle)}

\begin{definition}[Énoncé constatif]
Un \cle{énoncé constatif} (ou assertif) décrit un état de choses, affirme une proposition qui peut être \textbf{vraie ou fausse}. Il engage la responsabilité du locuteur sur la véracité du contenu.
\begin{itemize}
\item \textit{Exemple :} « Le taux de réussite au Probatoire est de 65 \%. » (Vérifiable).
\item \textit{Exemple :} « Meka a reçu une médaille. » (Fait narratif).
\end{itemize}
\end{definition}

\begin{definition}[Énoncé performatif]
Un \cle{énoncé performatif} \textbf{réalise l'action qu'il énonce} par le seul fait d'être prononcé (dans des conditions de réception appropriées : autorité, contexte, forme). Il n'est ni vrai ni faux, mais \textbf{heureux} (réussi) ou \textbf{inheureux} (raté).
\begin{itemize}
\item \textit{Exemple :} « Je déclare la séance ouverte. » (Le dire, c'est l'ouvrir).
\item \textit{Exemple :} « Je vous nomme délégué de classe. » (La nomination est effective).
\item \textit{Exemple :} « Je promets de rendre le devoir demain. » (Engagement créé).
\end{itemize}
\end{definition}

\begin{propriete}[Critères de reconnaissance du performatif explicite]
\begin{enumerate}
\item Verbe performatif à la \textbf{1ère personne du singulier} (je).
\item \textbf{Présent de l'indicatif} (ou présent de performation).
\item Verbe d'action langagière : \emph{déclarer, nommer, promettre, jurer, ordonner, interdire, baptiser, léguer, parier}...
\item Possibilité d'insérer \emph{« par la présente »} : « Je \textbf{par la présente} vous nomme... »
\end{enumerate}
\end{propriete}

\begin{exemplebox}[Constatif vs Performatif dans \emph{Le Vieux nègre et la médaille}]
\begin{itemize}
\item \textbf{Constatif (narration) :} « Meka \textbf{attendait} depuis l'aube. » $\rightarrow$ Description d'un fait.
\item \textbf{Performatif (dialogue) :} « Je \textbf{vous remets} cette médaille au nom du Gouverneur. » $\rightarrow$ L'acte de remise est accompli par la parole du chef de poste.
\item \textbf{Performatif implicite (ordre) :} « Sortez ! » $\rightarrow$ Équivaut à « Je vous ordonne de sortir. »
\end{itemize}
\end{exemplebox}

\begin{attention}[Piège en contraction de texte]
Ne transformez pas un performatif en constatif (perte de la force illocutoire).
\begin{itemize}
\item \textbf{Texte :} « Le Directeur \textbf{a déclaré} : "Je \textbf{nomme} M. X directeur adjoint." »
\item \textbf{Mauvais résumé :} « Le Directeur a dit que M. X était nommé. » (Constatif affaibli).
\item \textbf{Bon résumé :} « Le Directeur \textbf{a nommé} M. X directeur adjoint. » (Verbe performatif rapporté au passé, force conservée).
\end{itemize}
\end{attention}

\coursec{Ironie et paradoxe : figures de pensée}

\courssub{L'ironie : dire le contraire de ce qu'on pense}

\begin{definition}[Ironie]
L'\cle{ironie} est une figure de pensée qui consiste à \textbf{substituer à une idée l'expression de son contraire}, dans une intention moqueuse, critique ou amusée. Elle suppose une \textbf{complicité} entre l'émetteur et le récepteur (partage de connaissances contextuelles).
\end{definition}

\begin{propriete}[Marqueurs de l'ironie]
\begin{itemize}
\item \textbf{Contexte} : décalage entre le ton et la situation (éloge pour un échec, plainte pour un succès).
\item \textbf{Hyperbole / Litote} : exagération ou atténuation volontaire.
\item \textbf{Guillemets / Italique / Ponctuation} : « Il a fait un \emph{merveilleux} travail » (sarcasme écrit).
\item \textbf{Antiphrase} : forme lexicale de l'ironie (mot utilisé en sens inverse).
\end{itemize}
\end{propriete}

\begin{exemplebox}[Ironie chez Oyono (\emph{Le Vieux nègre et la médaille})]
« Meka était \textbf{fier} de sa médaille. Il la montrait à tout le monde : "Regardez, \textbf{les Blancs} m'aiment !" »
$\rightarrow$ L'ironie dramatique : le lecteur sait que la médaille est une manipulation coloniale (mépris déguisé en honneur). Meka est dupe ; le lecteur ne l'est pas.
\textbf{Enjeu pour le résumé :} Ne pas écrire « Meka était heureux de sa médaille » (prise au premier degré), mais « Meka \textbf{s'enorgueillissait naïvement} d'une médaille \textbf{ironique} / \textbf{instrumentalisée}. »
\end{exemplebox}

\courssub{Le paradoxe : l'évidence contrariée}

\begin{definition}[Paradoxe]
Le \cle{paradoxe} est une affirmation \textbf{apparemment contradictoire ou absurde}, mais qui révèle une \textbf{vérité profonde} en forçant la réflexion. Contrairement à l'ironie, il ne vise pas la moquerie mais l'étonnement intellectuel.
\end{definition}

\begin{exemplebox}[Exemples de paradoxes]
\begin{itemize}
\item « \textbf{Il faut se hâter lentement. »} (Auguste / Boileau) — Contradiction vitesse/lenteur $\rightarrow$ Sagesse de la précipitation réfléchie.
\item « \textbf{Je ne sais qu'une chose, c'est que je ne sais rien.} » (Socrate) — Auto-référence négative $\rightarrow$ Revendication de l'ignorance comme savoir premier.
\item Dans \emph{Le Vieux nègre et la médaille} : « \textbf{Pour être respecté des Blancs, Meka a dû trahir les siens.} » — Paradoxe de la collaboration : l'honneur apparent (médaille) scelle la honte réelle.
\end{itemize}
\end{exemplebox}

\begin{methode}[Traiter ironie et paradoxe en contraction]
\begin{enumerate}
\item \textbf{Repérer} le décalage (ton vs contexte, affirmation vs évidence).
\item \textbf{Identifier} la cible (qui/quoi est visé ?) et la thèse réelle de l'auteur.
\item \textbf{Reformuler} la thèse réelle (sens profond) sans conserver l'énoncé littéral ironique/paradoxal, sauf citation brève entre guillemets pour l'analyse.
\end{enumerate}
\end{methode}

\coursec{Les variétés lexicales du français}

\courssub{Diaphasie, diatopie, diastratie : trois axes de variation}

\begin{definition}[Variété de langue]
Une \cle{variété du français} est une forme de la langue caractérisée par des traits lexicaux, syntaxiques ou phonétiques propres, liée à l'usage (diaphasie), à l'espace (diatopie) ou au groupe social (diastratie). Le \textbf{lexique} est le niveau le plus visible de cette variation.
\end{definition}

\begin{center}
\begin{tabularx}{\linewidth}{|l|X|X|}
\hline
\rowcolor{popblueL}
\textbf{Axe} & \textbf{Variété (Exemples camerounais)} & \textbf{Traits lexicaux saillants} \\
\hline
\textbf{Diaphasique} (Registre) & Soutenu / Courant / Familier / Argots & \emph{Soutenu :} « décéder, s'abstenir, par conséquent ». \emph{Familier :} « claquer, se barrer, du coup ». \emph{Argot (Camfranglais) :} « do (dormir), nga (fille), tchop (manger) ». \\
\hline
\textbf{Diatopique} (Géographique) & Français du Cameroun, de Côte d'Ivoire, du Québec, de Belgique & \emph{Cameroun :} « un taxi-clando, un appelant, un dossier (affaire), se pointer (arriver), le go (la fille), le pain (l'argent) ». \emph{Québec :} « un char (voiture), un dépanneur (épicerie), magasiner (faire les courses) ». \\
\hline
\textbf{Diastratique} (Social) & Jargon professionnel, sociolecte jeunes, langue administrative & \emph{Admin :} « par la présente, ci-joint, ouï le rapport, vu la loi... ». \emph{Médical :} « hypertension, hémorragie, protocole ». \emph{Jeunes :} « wesh, ça va le faire, c'est tendu ». \\
\hline
\end{tabularx}
\end{center}

\courssub{Emprunts, calques et glissements sémantiques}

\begin{propriete}[Mécanismes lexicaux au Cameroun]
\begin{itemize}
\item \textbf{Emprunt intégral :} \emph{clando} (anglais \emph{clandestine}), \emph{bush taxi}, \emph{call-box}.
\item \textbf{Calque sémantique :} « \emph{Appeler} quelqu'un » (au sens de « téléphoner » $\leftarrow$ anglais \emph{to call}) $\rightarrow$ « Je t'\textbf{appelle} » au lieu de « Je te \textbf{téléphone} ».
\item \textbf{Glissement de sens (métaphore/métonymie) :} « \emph{Un dossier} » = un problème, une affaire judiciaire ou administrative. « \emph{Manger} un dossier » = le traiter (ou le détourner). « \emph{Le pain} » = l'argent, le salaire.
\item \textbf{Néologismes de nécessité :} \emph{Décentralisation, déconcentration, commune d'arrondissement, délégué du gouvernement}.
\end{itemize}
\end{propriete}

\begin{attention}[Variété et contraction : vigilance]
\begin{itemize}
\item \textbf{Identifiez le registre du texte source} : un rapport administratif (soutenu/technique) ne se résume pas en langage familier.
\item \textbf{Conservez les termes techniques/juridiques} (ex: \emph{délégation de compétences, transfert de compétences, redevabilité}) : ce sont des \textbf{concepts} non synonymiques.
\item \textbf{Neutralisez le familier/argot} sauf si le sujet est l'étude de la langue elle-même : « Il a \emph{claqué} la porte » $\rightarrow$ « Il est \textbf{parti brutalement} / a \textbf{démissionné} ».
\end{itemize}
\end{attention}

\coursec{Dégager la structure argumentative du texte}

\courssub{Au-delà de la syntaxe : la macrostructure}

La contraction de texte ne se fait pas phrase à phrase, mais \textbf{paragraphe par paragraphe} (ou séquence par séquence). Il faut repérer l'\cle{architecture argumentative} : comment l'auteur enchaîne ses idées pour convaincre, informer ou narrer.

\begin{propriete}[Schémas argumentatifs classiques]
\begin{itemize}
\item \textbf{Déductif (Thèse $\rightarrow$ Arguments) :} Idée générale annoncée d'emblée, puis développée. (Rapport, éditorial).
\item \textbf{Inductif (Arguments $\rightarrow$ Thèse) :} Faits, exemples, statistiques $\rightarrow$ Conclusion finale. (Enquête, article d'analyse).
\item \textbf{Dialectique (Thèse / Antithèse / Synthèse) :} Présentation d'un point de vue, objections, dépassement. (Dissertation, discours politique).
\item \textbf{Chronologique / Narratif :} Situation initiale $\rightarrow$ Événement déclencheur $\rightarrow$ Péripéties $\rightarrow$ Dénouement. (Roman, compte-rendu d'activité).
\item \textbf{Analytique (Description / Analyse / Proposition) :} Constat $\rightarrow$ Causes $\rightarrow$ Solutions. (Note de synthèse, rapport technique).
\end{itemize}
\end{propriete}

\begin{methode}[Dégager la structure argumentative (Méthode des 4 lectures)]
\begin{enumerate}
\item \textbf{Lecture 1 (Globale) :} Identifier le \textbf{sujet}, la \textbf{thèse principale} (souvent en titre, intro ou conclusion), le \textbf{type de texte} et le \textbf{ton}.
\item \textbf{Lecture 2 (Découpage) :} Repérer les \textbf{paragraphes} (ou séquences) et leur \textbf{fonction} : \emph{Introduction, Thèse 1, Argument 1, Exemple, Transition, Thèse 2, Conclusion...}. Numéroter les parties.
\item \textbf{Lecture 3 (Hiérarchisation) :} Dans chaque partie, distinguer \textbf{Idée force} (noyau) et \textbf{Illustrations/Explications} (périphérie). Souligner les \textbf{connecteurs logiques} (\emph{donc, cependant, en effet, par conséquent, en somme}).
\item \textbf{Lecture 4 (Synthèse) :} Rédiger le \textbf{plan détaillé} du texte (titres de parties + idées forces en une phrase chacune). C'est la "charpente" du futur résumé.
\end{enumerate}
\end{methode}

\begin{exemplebox}[Application à un éditorial type \emph{Cameroon Tribune}]
\textbf{Titre :} « Décentralisation : l'urgence des moyens. »
\begin{enumerate}
\item \textbf{Intro (§1) :} Constat : Loi 2004 votée, mais application lente. \textbf{Thèse :} « Sans moyens, la décentralisation reste un vœu pieux. »
\item \textbf{§2 (Argument 1 - Juridique) :} Texte de loi clair (compétences transférées). \textbf{Idée force :} Cadre légal existant.
\item \textbf{§3 (Argument 2 - Financier - Noyau) :} Fonds non suivis / retenus au central. Chiffres à l'appui. \textbf{Idée force :} Blocage budgétaire = cause principale.
\item \textbf{§4 (Argument 3 - Humain) :} Manque de formation des élus locaux. \textbf{Idée force :} Incompétence technique aggravante.
\item \textbf{Conclusion (§5) :} Appel à l'État : transferts effectifs + formation. \textbf{Synthèse :} la décentralisation exige des moyens réels.
\end{enumerate}
\end{exemplebox}


\coursec{Énoncés constatifs et performatifs}

\courssub{Transition : de la structure de la phrase à la force de l'énoncé}

Dans la section précédente, nous avons vu que la \cle{structure de la phrase} (sujet, verbe, compléments, subordination) est le squelette sur lequel se greffe la contraction de texte : supprimer des expansions, nominaliser des propositions, fusionner des phrases simples en phrases complexes. Mais une phrase n'est pas seulement une architecture syntaxique ; c'est un \cle{acte de langage}. Avant de réduire un texte, il faut comprendre \emph{ce que fait} l'énoncé : décrit-il le monde ou le transforme-t-il ? C'est la distinction fondamentale entre \cle{énoncé constatif} et \cle{énoncé performatif}, théorisée par le philosophe John Austin. Maîtriser cette distinction est indispensable pour \emph{reformuler} sans trahir l'intention de l'auteur, notamment dans les passages dialogués ou officiels du \emph{Vieux nègre et la médaille} de Ferdinand Oyono.

\courssub{L'énoncé constatif : dire le vrai ou le faux}

\begin{definition}[Énoncé constatif]
Un \textbf{énoncé constatif} (ou assertif) est un énoncé qui \textbf{décrit un état de choses}, qui \textbf{représente une réalité} extérieure au langage. Il porte sur un contenu propositionnel qui peut être jugé \textbf{vrai} ou \textbf{faux}. C'est le mode par défaut de l'information factuelle.
\end{definition}

\begin{exemplebox}[Exemples de constatifs]
\begin{itemize}
    \item « Le ciel est bleu. » (Vrai ou faux selon la météo)
    \item « Meka a reçu la médaille. » (Fait vérifiable dans le roman)
    \item « Le commandant arrive demain. » (Prédiction vérifiable a posteriori)
\end{itemize}
\end{exemplebox}

\begin{propriete}[Test de vérité]
Un énoncé est constatif s'il accepte l'interrogation : « \emph{Est-ce vrai que... ?} » ou « \emph{Est-ce faux que... ?} ». Si la réponse a un sens, c'est un constatif.
\end{propriete}

Dans un résumé, les constatifs se traitent par \textbf{condensation factuelle} : on regroupe les faits, on élimine les détails, on conserve la chronologie ou la logique causale. Exemple : « Meka a travaillé, a attendu, a reçu la médaille » $\rightarrow$ « Meka reçoit la médaille après une longue attente. »

\courssub{L'énoncé performatif : faire par le dire}

\begin{definition}[Énoncé performatif]
Un \textbf{énoncé performatif} est un énoncé qui \textbf{accomplit l'action qu'il énonce} par le seul fait d'être proféré dans des conditions appropriées. Il ne \emph{décrit} pas une action, il \emph{la réalise}. On ne peut pas le juger « vrai » ou « faux », mais « \textbf{heureux} » (réussi) ou « \textbf{malheureux} » (raté, inopérant).
\end{definition}

\begin{exemplebox}[Exemples canoniques de performatifs]
\begin{itemize}
    \item « \textbf{Je déclare} la séance ouverte. » (L'énoncé \emph{est} l'ouverture)
    \item « \textbf{Je vous promets} de venir. » (L'énoncé \emph{crée} la promesse)
    \item « \textbf{Je te baptise} Jean. » (L'énoncé \emph{effectue} le baptême)
    \item « \textbf{Je vous déclare} mari et femme. » (L'énoncé \emph{scelle} le mariage)
\end{itemize}
\end{exemplebox}

\begin{aretenir}[Critères formels du performatif explicite (Austin)]
Souvent (mais pas toujours), le performatif explicite présente cette structure :
\begin{enumerate}
    \item \textbf{Première personne du singulier} (Je) ou du pluriel (Nous).
    \item \textbf{Verbe performatif} au \textbf{présent de l'indicatif} (ou présentatif : « Je promets », « Nous décrétons »).
    \item Possibilité d'insérer \textbf{\emph{par la présente}} : « Je \emph{par la présente} vous nomme directeur. »
\end{enumerate}
\end{aretenir}

\begin{popfigure}
\begin{center}
\begin{tikzpicture}[
    node distance=1.2cm and 2cm,
    box/.style={draw=popblue, thick, rounded corners, fill=popblue!5, minimum width=3.5cm, minimum height=1.2cm, align=center, font=\small},
    arrow/.style={-Stealth, thick, popdark},
    labelbox/.style={draw=poporange, thick, rounded corners, fill=poporange!10, minimum width=4cm, minimum height=1cm, align=center, font=\small\bfseries}
]
% Nœuds
\node[labelbox] (const) {ÉNONCÉ CONSTATIF};
\node[labelbox, right=of const] (perf) {ÉNONCÉ PERFORMATIF};

\node[box, below=of const] (c1) {Décrit une réalité};
\node[box, below=of c1] (c2) {Vérifiable : VRAI / FAUX};
\node[box, below=of c2] (c3) {Ex : « Il pleut. »};

\node[box, below=of perf] (p1) {Accomplit une action};
\node[box, below=of p1] (p2) {Évaluable : HEUREUX / MALHEUREUX};
\node[box, below=of p2] (p3) {Ex : « Je promets. »};

% Flèches
\draw[arrow] (const.south) -- (c1.north);
\draw[arrow] (c1.south) -- (c2.north);
\draw[arrow] (c2.south) -- (c3.north);

\draw[arrow] (perf.south) -- (p1.north);
\draw[arrow] (p1.south) -- (p2.north);
\draw[arrow] (p2.south) -- (p3.north);
\end{tikzpicture}
\end{center}
\legende{Opposition fondamentale entre énoncé constatif (description, vérité) et énoncé performatif (action, félicité).}
\end{popfigure}

\courssub{Les verbes performatifs : un répertoire pour l'action langagière}

Tous les verbes ne sont pas performatifs. Seuls les \cle{verbes performatifs} (ou verbes d'action langagière) permettent de réaliser l'acte en le nommant. On les classe par \textbf{force illocutoire} (type d'action accomplie) selon la typologie de Searle, héritière d'Austin.

\begin{center}
\begin{tabularx}{\linewidth}{|l|X|l|}
\hline
\rowcolor{popblueL}
\textbf{Classe (Searle)} & \textbf{Action accomplie} & \textbf{Verbes performatifs types} \\
\hline
\textbf{Assertifs} & Engager le locuteur sur la vérité & \emph{affirmer, jurer, certifier, déposer, avouer} \\
\hline
\textbf{Directifs} & Faire faire quelque chose à l'autre & \emph{ordonner, commander, prier, supplier, conseiller, inviter} \\
\hline
\textbf{Commissifs} & Engager le locuteur pour l'avenir & \emph{promettre, jurer, s'engager, garantir, menacer} \\
\hline
\textbf{Expressifs} & Exprimer un état psychologique & \emph{remercier, féliciter, présenter ses excuses, déplorer, souhaiter} \\
\hline
\textbf{Déclaratifs} & Changer la réalité institutionnelle & \emph{déclarer, nommer, baptiser, condamner (justice), absoudre, décréter, ouvrir/clore (séance)} \\
\hline
\end{tabularx}
\end{center}

\begin{attention}[Piège : performatif ou constatif au futur ?]
Ne confondez pas :
\begin{itemize}
    \item \textbf{Performatif} : « \textbf{Je promets} de venir. » $\rightarrow$ L'acte de promettre est fait \emph{maintenant}.
    \item \textbf{Constatif (prédiction)} : « \textbf{Je viendrai} demain. » $\rightarrow$ Je décris une action future ; je ne m'engage pas formellement par la forme performative (pas de verbe « promettre »).
\end{itemize}
\textbf{Test} : Ajoutez « \emph{par la présente} ». « Je \emph{par la présente} viendrai demain » $\rightarrow$ \textbf{Impossible}. « Je \emph{par la présente} promets de venir » $\rightarrow$ \textbf{Correct}.
\end{attention}

\courssub{Conditions de réussite (félicité) d'un acte performatif}

Dire « Je vous nomme ministre » ne suffit pas pour en faire un ministre. Austin distingue trois conditions de \cle{félicité} (conditions de réussite) :

\begin{methode}[Vérifier la félicité d'un performatif (conditions d'Austin)]
Pour qu'un énoncé performatif soit \textbf{heureux} (opérant), trois types de conditions doivent être réunis :
\begin{enumerate}
    \item \textbf{Conditions de procédure} : il doit exister une \textbf{procédure conventionnelle} acceptée, et les \textbf{circonstances} (personnes, lieu, moment) doivent être appropriées.
        \begin{itemize}
            \item \emph{Ex.} : seul un maire (ou officier d'état civil) peut marier ; pas n'importe qui, pas n'importe où.
        \end{itemize}
    \item \textbf{Conditions d'exécution} : la procédure doit être \textbf{exécutée correctement} (formule consacrée, gestes rituels) et \textbf{complètement} (jusqu'au bout).
    \item \textbf{Conditions de sincérité} : le locuteur doit avoir les \textbf{intentions, sentiments, croyances} requis, et \textbf{agir ensuite} conformément à l'engagement.
        \begin{itemize}
            \item \emph{Ex.} : promettre sans intention de tenir $\rightarrow$ performatif \textbf{malheureux} (abus de procédure), mais l'acte existe (la promesse a été faite).
        \end{itemize}
\end{enumerate}
\textbf{Conséquence pour le résumé} : si un performatif est \emph{malheureux} dans le texte (ironie, usurpation), le résumé doit le signaler : « Le commandant \emph{prétend} décorer Meka » et non « Le commandant décore Meka ».
\end{methode}

\courssub{Repérage dans les discours officiels camerounais et dans \emph{Le Vieux nègre et la médaille}}

Les discours officiels (discours présidentiels, arrêtés préfectoraux, jugements, cérémonies protocolaires) sont le terrain par excellence des \textbf{performatifs déclaratifs} et \textbf{expressifs}. L'autorité du locuteur (Président, Préfet, Commandant) est la condition \emph{sine qua non} de leur félicité.

\begin{exemplebox}[Discours officiel camerounais : types d'énoncés]
\begin{itemize}
    \item \textbf{Constatif} : « Le taux de croissance est de 4\,\%. » (Rapport d'activité, vérifiable).
    \item \textbf{Performatif déclaratif} : « \textbf{Je déclare} ouverte la session parlementaire. » (Président de l'Assemblée, formule consacrée).
    \item \textbf{Performatif déclaratif (arrêté)} : « \textbf{Arrêté} : Article 1 : \textbf{Est nommé} Directeur général... M. X. » (Premier Ministre / Président, pouvoir réglementaire).
    \item \textbf{Performatif expressif} : « \textbf{Je félicite} les lauréats. » (Autorité hiérarchique en cérémonie).
\end{itemize}
\end{exemplebox}

Dans \emph{Le Vieux nègre et la médaille} (Ferdinand Oyono, 1956), la scène de la remise de la médaille est une \textbf{machine performative géante}. Le Commandant d'arrondissement utilise son autorité institutionnelle pour transformer Meka en « chevalier » par la parole.

\begin{exemplebox}[Analyse performative de la scène de remise (extrait reconstitué)]
\begin{itemize}
    \item \textbf{Constatif (préambule)} : « Meka a servi la France pendant de longues années. » $\rightarrow$ Fait présenté comme vérifiable (mais contesté par l'ironie du roman).
    \item \textbf{Performatif expressif} : « \textbf{Je vous félicite}, Meka. » $\rightarrow$ Crée la reconnaissance officielle.
    \item \textbf{Performatif déclaratif} : « Au nom du Président de la République, \textbf{je vous remets} la médaille. » $\rightarrow$ \textbf{Change le statut} de Meka (civil $\rightarrow$ décoré). Condition : autorité déléguée du Commandant.
    \item \textbf{Performatif expressif} : « \textbf{Je vous remercie} pour vos loyaux services. » $\rightarrow$ Clôture symbolique de l'échange.
\end{itemize}
\end{exemplebox}

\begin{exemplebox}[Densité performative dans la scène]
Dans le passage de la remise de médaille, le discours du Commandant est \textbf{saturé d'énoncés performatifs explicites} : féliciter, remettre, décorer, nommer, remercier, souhaiter. Cette saturation performative \textbf{construit l'ironie} : la puissance du dire officiel masque le vide humain (Meka ne comprend pas le français, sa famille est humiliée, la médaille est une coquille vide). Le résumé doit rendre cette \emph{force illocutoire} sans la noyer dans le récit.
\end{exemplebox}

\begin{popfigure}
\begin{center}
\begin{tikzpicture}[
    node distance=0.8cm and 1cm,
    header/.style={draw=popdark, fill=popblueL, thick, minimum height=0.9cm, font=\bfseries\small, align=center},
    cell/.style={draw=popline, minimum height=0.85cm, font=\small, align=left, text width=7cm},
    cellc/.style={draw=popline, minimum height=0.85cm, font=\small, align=center, text width=2.5cm},
    rowodd/.style={fill=popcream},
    roweven/.style={fill=white}
]
% En-têtes
\node[header, text width=3cm] (h1) {Énoncé extrait / reconstitué};
\node[header, text width=2.5cm, right=0pt of h1] (h2) {Type};
\node[header, text width=2.5cm, right=0pt of h2] (h3) {Force illocutoire};
\node[header, text width=2.5cm, right=0pt of h3] (h4) {Félicité dans le roman};

% Ligne 1
\node[cell, rowodd, below=0pt of h1] (r1c1) {« Meka, vous avez bien servi. »};
\node[cellc, rowodd, right=0pt of r1c1] (r1c2) {Constatif};
\node[cellc, rowodd, right=0pt of r1c2] (r1c3) {Assertif};
\node[cellc, rowodd, right=0pt of r1c3] (r1c4) {Vrai (selon l'administration)};

% Ligne 2
\node[cell, roweven, below=0pt of r1c1] (r2c1) {« Je vous félicite de tout cœur. »};
\node[cellc, roweven, right=0pt of r2c1] (r2c2) {Performatif};
\node[cellc, roweven, right=0pt of r2c2] (r2c3) {Expressif};
\node[cellc, roweven, right=0pt of r2c3] (r2c4) {Heureuse (autorité OK)};

% Ligne 3
\node[cell, rowodd, below=0pt of r2c1] (r3c1) {« Au nom du Président, je vous remets cette médaille. »};
\node[cellc, rowodd, right=0pt of r3c1] (r3c2) {Performatif};
\node[cellc, rowodd, right=0pt of r3c2] (r3c3) {Déclaratif};
\node[cellc, rowodd, right=0pt of r3c3] (r3c4) {Heureuse (délégation OK)};

% Ligne 4
\node[cell, roweven, below=0pt of r3c1] (r4c1) {« Je vous nomme Chevalier de la Légion d'honneur. »};
\node[cellc, roweven, right=0pt of r4c1] (r4c2) {Performatif};
\node[cellc, roweven, right=0pt of r4c2] (r4c3) {Déclaratif};
\node[cellc, roweven, right=0pt of r4c3] (r4c4) {Heureuse (formule sacramentelle)};

% Ligne 5
\node[cell, rowodd, below=0pt of r4c1] (r5c1) {« Je souhaite que cet exemple soit suivi. »};
\node[cellc, rowodd, right=0pt of r5c1] (r5c2) {Performatif};
\node[cellc, rowodd, right=0pt of r5c2] (r5c3) {Expressif / Directif};
\node[cellc, rowodd, right=0pt of r5c3] (r5c4) {Heureuse (souhait émis)};

% Ligne 6 (Ironie interne)
\node[cell, roweven, below=0pt of r5c1] (r6c1) {« Meka ne comprend pas un mot du discours. » (Narrateur)};
\node[cellc, roweven, right=0pt of r6c1] (r6c2) {Constatif};
\node[cellc, roweven, right=0pt of r6c2] (r6c3) {Assertif};
\node[cellc, roweven, right=0pt of r6c3] (r6c4) {Vrai (réalité du personnage)};
\end{tikzpicture}
\end{center}
\legende{Classification des énoncés de la scène de remise de médaille (Commandant vs Narrateur). Les performatifs du Commandant sont « heureux » institutionnellement mais ironiques humainement.}
\end{popfigure}

\courssub{Intérêt majeur pour la contraction de texte (résumé)}

C'est le point nodal pour votre épreuve. Un résumé n'est pas une simple liste de phrases ; c'est une \textbf{reformulation de la structure argumentative et énonciative}.

\begin{methode}[Reformuler un performatif dans un résumé]
Ne recopiez \textbf{jamais} le performatif tel quel (discours direct) sauf citation brève stratégique. Transformez-le en \textbf{constatif méta-langagier} (discours indirect rapporté) qui explicite l'acte accompli.
\begin{enumerate}
    \item \textbf{Identifiez} le verbe performatif et sa force (déclarer, promettre, féliciter, ordonner...).
    \item \textbf{Identifiez} le locuteur (qui a l'autorité ?) et le destinataire.
    \item \textbf{Reformulez} selon le schéma : \textbf{[Locuteur] + [Verbe de parole performatif au passé/imparfait] + [Contenu de l'acte]}.
\end{enumerate}
\end{methode}

\begin{exoresolu}[Transformation performatif $\rightarrow$ résumé]
\textbf{Texte source (extrait) :}
\begin{quote}
« Mesdames, Messieurs, \textbf{je déclare} la session 2025 du Baccalauréat technique officiellement ouverte. \textbf{Je souhaite} plein succès à tous les candidats. \textbf{Je vous prie} de respecter le silence dans les salles d'examen. »
\end{quote}

\textbf{Analyse :}
\begin{itemize}
    \item « Je déclare... ouverte » $\rightarrow$ Performatif \textbf{Déclaratif} (change l'état du monde : session fermée $\rightarrow$ ouverte). Locuteur : Ministre/Président de jury.
    \item « Je souhaite... » $\rightarrow$ Performatif \textbf{Expressif}.
    \item « Je vous prie de... » $\rightarrow$ Performatif \textbf{Directif} (ordre déguisé en politesse).
\end{itemize}
\tcblower
\textbf{Résumé correct (niveau Probatoire/Bac Tech) :}
\begin{quote}
« Le Ministre \textbf{a déclaré ouverte} la session 2025 du Baccalauréat technique, \textbf{a souhaité} succès aux candidats et \textbf{a enjoint} au silence dans les salles. »
\end{quote}

\textbf{Erreurs à ne pas commettre :}
\begin{itemize}
    \item[\textbf{×}] « Le Ministre a dit : "Je déclare la session ouverte." » (Citation inutile, alourdit).
    \item[\textbf{×}] « La session a commencé. » (Perte de l'acte d'autorité, du locuteur, de la solennité).
    \item[\textbf{✓}] « Le Ministre \textbf{a ouvert} la session par une déclaration officielle. » (Synthèse acceptable si la place manque).
\end{itemize}
\end{exoresolu}

\begin{attention}[Cas délicat : performatif \emph{malheureux} ou ironique]
Dans \emph{Le Vieux nègre et la médaille}, le performatif « Je vous nomme Chevalier » est institutionnellement heureux mais \textbf{humainement vide} (Meka ne comprend pas, la médaille est volée ensuite).
\begin{itemize}
    \item \textbf{Mauvais résumé} : « Le Commandant nomme Meka Chevalier. » (Valide le mensonge social).
    \item \textbf{Bon résumé} : « Le Commandant \textbf{décerne officiellement} la médaille à Meka, dans une cérémonie que le héros ne comprend pas. » (Marque l'acte \emph{et} son échec communicationnel).
    \item \textbf{Résumé analytique (niveau excellence)} : « La remise de médaille repose sur une \textbf{saturation performative} (féliciter, décorer, nommer) qui masque l'incommunicabilité et l'ironie du sort de Meka. »
\end{itemize}
\end{attention}

\courssub{Exercice d'application guidé}

\begin{exercice}[Repérage et reformulation]
Voici un extrait mêlant constatifs et performatifs (inspiré d'un arrêté préfectoral fictif).
\begin{quote}
« \textbf{Considérant} la loi n°2023/001 du 15 janvier 2023 ; \\
\textbf{Considérant} la nécessité de préserver l'ordre public ; \\
\textbf{Arrête} : \\
Article 1 : \textbf{Est interdite} la circulation des poids lourds sur le pont du Wouri de 6h à 20h. \\
Article 2 : \textbf{Est chargé} de l'exécution du présent arrêté le Commandant de la Légion de Gendarmerie. \\
\textbf{Fait} à Douala, le 10 octobre 2025. \\
\textbf{Le Préfet de la Wouri,} \\
\emph{Signé :} N. M. »
\end{quote}

\textbf{Questions :}
\begin{enumerate}
    \item Classez chaque segment souligné en \textbf{Constatif} ou \textbf{Performatif} (précisez la classe : Déclaratif, Directif...).
    \item Reformulez l'essentiel de cet arrêté en \textbf{3 phrases maximum} pour un résumé administratif.
\end{enumerate}
\end{exercice}

\begin{corrige}[Exercice n°1]
\textbf{1. Classification :}
\begin{itemize}
    \item « Considérant la loi... » / « Considérant la nécessité... » $\rightarrow$ \textbf{Constatifs} (Assertifs : exposent les motifs juridiques et factuels, base du raisonnement).
    \item « Arrête » (verbe performatif implicite : « J'arrête ») $\rightarrow$ \textbf{Performatif Déclaratif} (création de la norme).
    \item « Est interdite » $\rightarrow$ \textbf{Performatif Déclaratif} (création d'une obligation/interdiction juridique). Structure passive impersonnelle typique du style administratif = performatif de l'autorité.
    \item « Est chargé de l'exécution » $\rightarrow$ \textbf{Performatif Déclaratif} (délégation de compétence / nomination de fait).
    \item « Fait à Douala... Le Préfet » $\rightarrow$ \textbf{Constatifs} (circonstances de lieu, temps, auteur ; la signature vaut performatif d'authentification).
\end{itemize}

\textbf{2. Résumé administratif (3 phrases) :}
\begin{quote}
« Par arrêté du 10 octobre 2025, le Préfet de la Wouri \textbf{interdit} la circulation des poids lourds sur le pont du Wouri (6h-20h) pour préserver l'ordre public. Il \textbf{charge} le Commandant de la Légion de Gendarmerie de faire appliquer cette mesure. »
\end{quote}
\begin{itemize}
    \item \emph{Gain :} suppression des « Considérant » (motifs $\rightarrow$ cause finale « pour préserver... »), fusion « Arrête + Article 1 » $\rightarrow$ « interdit », « Article 2 » $\rightarrow$ « charge ». Les verbes performatifs sont conservés comme verbes d'action du résumé.
\end{itemize}
\end{corrige}

\begin{savaistu}[John Austin et la naissance de la pragmatique]
C'est le philosophe britannique \textbf{John Langshaw Austin} (1911--1960), professeur à Oxford, qui a révolutionné la philosophie du langage en montrant que « dire, c'est faire ». Dans ses cours de 1955, publiés posthumes sous le titre \emph{Comment faire des choses avec des mots} (\emph{How to Do Things with Words}, 1962), il démonte le mythe selon lequel le langage ne sert qu'à \emph{constater} (vérité/fausseté). Il invente les termes \textbf{performatif}, \textbf{constatif}, \textbf{locutionnaire}, \textbf{illocutoire}, \textbf{perlocutoire}. Son élève \textbf{John Searle} systématisera la classification des actes de langage (1969). Aujourd'hui, cette théorie fonde l'analyse du discours juridique, politique, littéraire et l'intelligence artificielle (agents conversationnels).
\end{savaistu}

\courssub{Synthèse de la section : ce qu'il faut retenir pour l'épreuve}

\begin{aretenir}[Check-list « Constatif / Performatif » pour le résumé]
\begin{enumerate}
    \item \textbf{Distinction} : Constatif = Vrai/Faux (description) ; Performatif = Heureux/Malheureux (action).
    \item \textbf{Repérage} : verbe à la 1\iere{} personne au présent + « par la présente » $\rightarrow$ Performatif explicite. Style administratif (passif « est interdit », « est nommé ») $\rightarrow$ Performatif institutionnel.
    \item \textbf{Conditions} : Autorité (Qui parle ?), Procédure (Formule), Sincérité (Intention). Vérifier si le texte les respecte ou les détourne (ironie).
    \item \textbf{Reformulation (règle d'or)} : \textbf{Locuteur + Verbe d'acte (passé) + Contenu}.
        \begin{itemize}
            \item « Je promets » $\rightarrow$ « Il a promis de... »
            \item « Je déclare ouvert » $\rightarrow$ « Il a ouvert / a déclaré ouverte... »
            \item « Je vous ordonne de partir » $\rightarrow$ « Il a ordonné le départ / a enjoint de partir. »
        \end{itemize}
    \item \textbf{Attention au contexte} : dans \emph{Le Vieux nègre et la médaille}, la prolifération des performatifs officiels (Commandant) contraste avec le silence constatif de Meka (qui ne peut que constater : « Je ne comprends pas »). Le résumé doit rendre ce \textbf{rapport de force énonciatif}.
\end{enumerate}
\end{aretenir}

\coursec{Ironie et paradoxe : figures de pensée}

Dans la section précédente, nous avons distingué les \cle{énoncés constatifs} (qui décrivent un état de fait et peuvent être vrais ou faux) des \cle{énoncés performatifs} (qui accomplissent une action par le seul fait d'être prononcés). Mais il arrive que le sens d'un énoncé ne coïncide pas du tout avec ce que les mots semblent dire : un constatif peut signifier le \emph{contraire} de ce qu'il affirme, et une affirmation peut heurter l'opinion commune tout en étant vraie. C'est le domaine des \cle{figures de pensée}, et en particulier de deux figures essentielles pour comprendre \emph{Le Vieux nègre et la médaille} de Ferdinand Oyono : l'\cle{ironie} et le \cle{paradoxe}. Or, pour résumer un texte argumentatif ou littéraire, il est capital de ne pas trahir ces figures : un résumé qui prend l'ironie au premier degré inverse complètement la pensée de l'auteur.

\courssub{L'ironie : dire le contraire de ce qu'on pense}

\textbf{Intuition d'abord.} Imaginez qu'un camarade arrive en cours avec une heure de retard, alors que le surveillant l'attendait à la porte. Vous lui dites : « Bravo, quelle ponctualité exemplaire ! » Personne ne croit que vous le félicitez réellement : tout le monde comprend que vous le raillez. Vous venez de produire une \cle{ironie}. Le mécanisme repose sur un \emph{écart} : écart entre ce que les mots disent (l'éloge) et ce que la situation réelle impose de penser (le reproche). C'est cet écart que le lecteur ou l'auditeur doit savoir détecter.

\begin{definition}[L'ironie]
L'\cle{ironie} est une figure de pensée qui consiste à dire le \textbf{contraire} de ce que l'on pense (ou à dire \textbf{plus} ou \textbf{moins} que ce que l'on pense), dans une intention de \textbf{raillerie}, de moquerie ou de critique. Celui qui emploie l'ironie compte sur l'intelligence du destinataire pour renverser le sens apparent et saisir le sens réel.
\end{definition}

L'ironie n'est donc pas un mensonge : le menteur veut être cru, tandis que l'ironiste veut être \emph{démasqué}. Toute sa stratégie consiste à laisser des indices (le ton, l'exagération, le contexte) qui signalent au lecteur : « ne me croyez pas au premier degré ».

\begin{definition}[L'antiphrase]
L'\cle{antiphrase} est une figure qui fait entendre le \textbf{contraire} de ce que disent les mots. C'est le procédé central de l'ironie : « Quel courage ! » dit d'un fuyard ; « Voilà un beau travail ! » dit d'une copie bâclée. L'antiphrase est une ironie concentrée en une phrase ou en quelques mots.
\end{definition}

\textbf{Les procédés de l'ironie.} L'ironiste dispose de plusieurs outils pour créer l'écart entre le dit et le pensé :

\begin{itemize}
\item L'\cle{antiphrase} : dire exactement l'inverse de ce qu'on pense (« Quelle belle médaille ! » pour une récompense dérisoire).
\item L'\cle{exagération} (ou hyperbole ironique) : pousser l'éloge ou la critique jusqu'à l'invraisemblance, de sorte que l'excès lui-même trahit la moquerie (« Ce fonctionnaire est un génie : il a retrouvé le dossier en seulement huit mois »).
\item Le \cle{faux éloge} : louer avec des termes apparemment flatteurs une personne ou une institution que l'on condamne en réalité. C'est le procédé favori d'Oyono envers l'administration coloniale.
\item Le \cle{ton décalé} : adopter un ton solennel, admiratif ou naïf pour raconter des faits injustes ou absurdes ; le contraste entre le ton et le contenu produit l'effet ironique.
\end{itemize}

\begin{methode}[Repérer l'ironie dans un texte]
Pour repérer l'ironie, cherchez l'\textbf{écart} entre ce qui est dit et la situation réelle décrite.
\begin{enumerate}
\item Relevez l'énoncé à valeur élogieuse ou admirative (« quelle belle récompense », « quel honneur »).
\item Confrontez-le aux faits que le texte rapporte : le personnage a-t-il réellement reçu quelque chose de précieux ? A-t-il été réellement honoré ?
\item Si l'éloge est disproportionné par rapport à la réalité (une médaille contre toute une vie de labeur et d'humiliations), vous êtes en présence d'une ironie.
\item Formulez alors le sens réel : « l'auteur veut dire que cette récompense est dérisoire ».
\end{enumerate}
\end{methode}

\begin{exemplebox}[« Quelle belle médaille ! »]
Dans \emph{Le Vieux nègre et la médaille}, Meka reçoit une médaille après des décennies de labeur, de terres confisquées et d'humiliations. Dire « Quelle belle médaille ! » au premier degré serait absurde : une simple médaille ne pèse rien face à une vie entière de souffrances. L'éloge apparent cache donc une \textbf{dénonciation} : l'ironie fait mesurer, par le contraste, l'immensité de l'injustice. Plus l'éloge est chaleureux, plus la critique est cinglante.
\end{exemplebox}

\courssub{Le paradoxe : contre l'opinion commune, mais vrai}

\textbf{Intuition d'abord.} On entend souvent : « Pour apprendre, il faut d'abord accepter de ne pas savoir. » Cette phrase semble se contredire : comment l'ignorance pourrait-elle mener au savoir ? Pourtant, à la réflexion, elle est profondément vraie : celui qui croit tout savoir n'apprend plus rien. Voilà un \cle{paradoxe} : une affirmation qui choque l'opinion commune, mais qui résiste à l'examen.

\begin{definition}[Le paradoxe]
Le \cle{paradoxe} est une affirmation \textbf{contraire à l'opinion commune} (à la \emph{doxa}), qui paraît absurde ou contradictoire à première vue, mais qui peut se révéler \textbf{vraie} ou féconde lorsqu'on l'examine de près. Contrairement au mensonge ou à l'erreur, le paradoxe invite à réfléchir : il force l'auditeur à dépasser les idées reçues.
\end{definition}

Il faut bien distinguer les deux figures :

\begin{center}
\begin{tabularx}{\linewidth}{|l|X|X|}
\hline
\rowcolor{popblueL} & \textbf{Ironie} & \textbf{Paradoxe} \\
\hline
Rapport au sens & Dit le contraire de ce qu'on pense & Affirme ce qui semble contraire à la vérité commune \\
\hline
Intention & Railler, critiquer, dénoncer & Faire réfléchir, ébranler les préjugés \\
\hline
Vérité & Le sens réel est l'inverse du sens littéral & L'affirmation elle-même peut être vraie \\
\hline
Exemple & « Quelle belle médaille ! » (elle est dérisoire) & « Cette médaille, censée honorer Meka, l'humilie » \\
\hline
\end{tabularx}
\end{center}

Remarquez le dernier exemple : « Cette médaille, censée honorer Meka, l'humilie » est un paradoxe. L'opinion commune veut qu'une récompense honore celui qui la reçoit ; le roman montre au contraire que cette récompense, par son insignifiance et par la manière dont elle est accordée, \emph{aggrave} l'humiliation. L'affirmation choque, mais elle est vraie dans l'univers du roman.

\begin{popfigure}
\begin{center}
\begin{tikzpicture}[
box/.style={draw=popdark, rounded corners, fill=popcream, text width=3.6cm, align=center, minimum height=1.4cm, font=\small},
sens/.style={draw=popink, rounded corners, fill=poppinkL, text width=3.6cm, align=center, minimum height=1.4cm, font=\small},
fleche/.style={-{Stealth[length=3mm]}, very thick, poporange},
etq/.style={font=\small\bfseries, text=popdark}
]
% --- Ironie ---
\node[box, align=center] (lit) at (0,2.2){\textbf{Énoncé littéral}\\ « Quelle belle médaille ! »};
\node[sens, align=center] (vis) at (6.2,2.2){\textbf{Sens visé}\\ « Cette récompense est dérisoire »};
\draw[fleche] (lit) -- (vis);
\node[etq] at (3.1,2.9) {IRONIE};
\node[font=\small\itshape, text=popmuted] at (3.1,1.6) {renversement du sens};

% --- Paradoxe ---
\node[box, align=center] (op) at (0,-1.2){\textbf{Opinion commune}\\ « Une médaille honore celui qui la reçoit »};
\node[sens, align=center] (par) at (6.2,-1.2){\textbf{Affirmation paradoxale}\\ « Cette médaille humilie Meka »};
\draw[fleche] (op) -- (par);
\node[etq] at (3.1,-0.5) {PARADOXE};
\node[font=\small\itshape, text=popmuted] at (3.1,-1.8) {contre l'opinion, mais vrai};
\end{tikzpicture}
\end{center}
\legende{Schéma comparatif des deux figures de pensée. En haut : l'ironie fait passer de l'énoncé littéral à son contraire (sens visé). En bas : le paradoxe part de l'opinion commune pour aboutir à une affirmation choquante mais vraie.}
\end{popfigure}

\courssub{Ironie et paradoxe dans \emph{Le Vieux nègre et la médaille}}

Le roman de Ferdinand Oyono (1956) est traversé par ces deux figures, qui en constituent le ressort critique principal.

\textbf{La médaille, récompense dérisoire.} Meka, vieux paysan camerounais, a tout donné à la France coloniale : ses terres, ses fils morts à la guerre, sa foi, son travail. En échange, l'administration coloniale lui décerne... une médaille, au cours d'une cérémonie pompeuse du 14 juillet. Tout le dispositif romanesque est ironique : les discours officiels célèbrent la « fidélité » et le « dévouement » du vieil homme, alors même que cette fidélité n'a été payée que d'humiliations. Le faux éloge des discours coloniaux est démonté par la réalité des faits que le lecteur connaît.

\textbf{Le paradoxe de la reconnaissance.} Le roman repose sur un paradoxe central : la plus haute distinction que la colonisation puisse offrir à un « indigène » est en réalité la preuve de son aliénation. Être « récompensé » par l'oppresseur, c'est être confirmé dans sa soumission. La médaille, symbole d'honneur aux yeux de l'opinion commune, devient symbole de la dépossession. Après la cérémonie, Meka, ivre, est arrêté et humilié par les mêmes forces de l'ordre qui l'avaient décoré quelques heures plus tôt : le paradoxe éclate alors au grand jour — l'honneur colonial n'a duré que l'espace d'un discours.

\begin{propriete}[Fonction critique de l'ironie et du paradoxe]
Dans un texte à visée argumentative ou satirique, l'ironie et le paradoxe ont une \textbf{fonction de dénonciation}. Chez Oyono, ils servent à :
\begin{itemize}
\item dénoncer l'\textbf{administration coloniale}, dont la générosité affichée masque l'exploitation ;
\item dénoncer l'\textbf{illusion de la reconnaissance} : la médaille fait croire à une égalité et à une gratitude qui n'existent pas ;
\item éveiller la \textbf{lucidité du lecteur}, qui doit renverser les éloges apparents pour saisir la critique réelle.
\end{itemize}
\end{propriete}

\begin{savaistu}[Ferdinand Oyono, de l'ironie romanesque à la lucidité politique]
Ferdinand Oyono, auteur du \emph{Vieux nègre et la médaille} (1956) et d'\emph{Une vie de boy} (1956), est devenu après l'indépendance du Cameroun un diplomate de haut rang (ambassadeur dans plusieurs pays, représentant à l'ONU) puis ministre. L'ironie mordante de ses romans de jeunesse, qui démontaient déjà les discours officiels et les cérémonies creuses, annonçait la lucidité politique de l'homme d'État : savoir lire derrière les mots est une compétence de l'écrivain comme du diplomate.
\end{savaistu}

\courssub{Traitement de l'ironie et du paradoxe dans le résumé}

C'est ici que cette leçon rejoint directement notre objet d'étude : la \cle{contraction de texte}. Lorsque vous résumez un texte ironique ou paradoxal, vous ne pouvez pas supprimer la valeur critique de ces figures, sous peine de trahir l'auteur.

\begin{attention}[Erreur fréquente : résumer au premier degré]
L'erreur la plus grave consiste à résumer une phrase ironique \textbf{au premier degré} et à attribuer à l'auteur une opinion qu'il critique. Écrire : « Selon Oyono, la médaille est une belle récompense qui honore dignement Meka » est un \textbf{contresens complet} : vous faites dire à l'auteur l'inverse de sa pensée. À l'examen, cette erreur est lourdement sanctionnée, car elle prouve que le candidat n'a pas compris le texte.
\end{attention}

\begin{methode}[Reformuler l'ironie et le paradoxe dans un résumé]
\begin{enumerate}
\item \textbf{Repérez} la figure (écart entre le dit et la situation, affirmation contraire à l'opinion commune).
\item \textbf{Traduisez} le sens réel en langage direct : remplacez l'éloge ironique par la critique qu'il cache.
\item \textbf{Signalez} la démarche de l'auteur par des formules explicites : « ironiquement », « avec ironie », « l'auteur dénonce », « l'auteur raille », « par ce paradoxe, Oyono montre que... ».
\item \textbf{Conservez} l'enjeu de la figure : le résumé doit garder l'idée que la récompense est dérisoire et que la reconnaissance coloniale est une illusion.
\end{enumerate}
\end{methode}

\begin{exemplebox}[Reformulation correcte]
Phrase du texte (premier degré apparent) : « Meka fut comblé : la France reconnaissante lui décerna une magnifique médaille. »\par
Résumé fautif : « L'auteur montre que la France a généreusement récompensé Meka. »\par
Résumé correct : « Ironiquement, l'auteur montre que la France se contente d'offrir à Meka une médaille dérisoire après une vie de sacrifices ; il dénonce ainsi l'illusion de la reconnaissance coloniale. »
\end{exemplebox}

\begin{exoresolu}[Repérer et reformuler une ironie]
Énoncé. Voici un passage réécrit dans l'esprit du roman : « Quelle magnifique journée pour Meka ! Le commandant en personne lui a accroché à la poitrine une médaille en échange de ses terres, de ses deux fils et de cinquante ans de labeur. Le vieil homme pouvait mourir tranquille : la patrie ne l'oublierait pas. » 1) Relevez deux procédés ironiques. 2) Reformulez ce passage en deux phrases pour un résumé, sans trahir l'auteur.
\tcblower
Solution détaillée.
1) Procédés ironiques :
\begin{itemize}
\item \textbf{Antiphrase} : « Quelle magnifique journée ! » fait entendre le contraire — c'est une journée qui scelle une vie de spoliation.
\item \textbf{Faux éloge et exagération} : « la patrie ne l'oublierait pas » est démenti par la réalité (une médaille contre des terres, deux fils et cinquante ans de labeur) ; l'énumération chiffrée crée une disproportion qui trahit la moquerie. Le ton décalé, solennel et admiratif, renforce l'écart.
\end{itemize}
2) Reformulation pour le résumé : « Avec ironie, l'auteur souligne la disproportion entre les sacrifices de Meka — ses terres, ses fils, cinquante ans de travail — et la simple médaille reçue en échange. Il dénonce ainsi l'ingratitude de l'administration coloniale et le caractère illusoire de sa reconnaissance. »\par
On remarque que le sens réel (critique) est exprimé directement, et que la démarche de l'auteur est signalée par « avec ironie » et « il dénonce ».
\end{exoresolu}

\begin{exercice}[À vous de jouer]
Dans chacune des phrases suivantes, dites s'il s'agit d'ironie, de paradoxe ou d'un énoncé au premier degré, puis reformulez-la pour un résumé :
\begin{enumerate}
\item « Le colonisateur est généreux : il prend vos terres et vous offre en retour un morceau de métal à épingler sur la poitrine. »
\item « C'est en perdant tout ce qu'il possédait que Meka a découvert ce que valait réellement la reconnaissance des hommes. »
\item « La cérémonie a eu lieu le 14 juillet, en présence des autorités. »
\end{enumerate}
\end{exercice}

\begin{corrige}[Exercice : ironie, paradoxe, premier degré]
1) \textbf{Ironie} (faux éloge, antiphrase : « généreux » signifie « spoliateur »). Reformulation : « L'auteur dénonce ironiquement la prétendue générosité coloniale, qui confisque les terres et n'offre en échange qu'une médaille sans valeur. »\par
2) \textbf{Paradoxe} : l'opinion commune veut que la perte appauvrisse ; ici, perdre tout rend lucide. Reformulation : « Par ce paradoxe, l'auteur montre que la dépossession totale a ouvert les yeux de Meka sur le néant de la reconnaissance coloniale. »\par
3) \textbf{Énoncé au premier degré} (constatif, simple information factuelle). Reformulation : « La remise de la médaille s'est déroulée le 14 juillet devant les autorités coloniales. » Aucune valeur critique à signaler ici.
\end{corrige}

\begin{aretenir}[L'essentiel de la section]
\begin{itemize}
\item L'\cle{ironie} dit le contraire (ou plus/moins) de ce qu'on pense pour railler ; ses procédés : antiphrase, exagération, faux éloge, ton décalé.
\item Le \cle{paradoxe} affirme contre l'opinion commune une idée qui peut être vraie.
\item Chez Oyono, ces figures dénoncent l'administration coloniale et l'illusion de la reconnaissance : la médaille est une récompense dérisoire.
\item Dans un résumé, ne jamais prendre l'ironie au premier degré : traduire le sens réel et signaler la démarche (« ironiquement », « l'auteur dénonce »).
\end{itemize}
\end{aretenir}

\coursec{Les variétés lexicales du français}

Après avoir analysé l'ironie et le paradoxe comme figures de pensée qui détournent le sens littéral pour en suggérer un autre, plus profond ou plus critique, nous devons maintenant porter notre attention sur les mots eux-mêmes. Une même réalité peut être désignée par des termes radicalement différents selon qui parle, d'où il parle et à qui il s'adresse. Maîtriser les \cle{variétés lexicales}, c'est savoir naviguer entre ces mondes lexicaux pour, in fine, les unifier dans le français standard exigé par l'exercice de contraction.

\courssub{Les trois registres de la langue française : repères lexicaux}

La langue n'est pas un bloc monolithique ; elle se déploie en \cle{registres de langue} (ou niveaux de langue), distingués principalement par le choix du vocabulaire, mais aussi par la syntaxe et la phonétique. En situation de rédaction administrative, scolaire ou journalistique — et donc dans le résumé —, seul le \cle{français standard} (ou courant) est la norme, le français soutenu étant réservé à l'écrit littéraire ou solennel, et le français familier à l'oral intime.

\begin{definition}[Variété lexicale et registre de langue]
Une \textbf{variété lexicale} est un ensemble de mots propres à une région, un groupe social, un corps de métier ou un niveau de langue (registre). Le \textbf{registre de langue} correspond à la variation de la langue selon la situation de communication (contexte, interlocuteurs, intention). On distingue classiquement trois registres majeurs :
\begin{itemize}
    \item \textbf{Soutenu} : vocabulaire rare, précis, syntaxe complexe (ex. : \emph{déambuler}, \emph{s'alimenter}, \emph{quitter les lieux}).
    \item \textbf{Standard (courant)} : vocabulaire usuel, syntaxe correcte, communication efficace (ex. : \emph{marcher}, \emph{manger}, \emph{partir}). \cle{C'est la norme du résumé.}
    \item \textbf{Familier} : vocabulaire expressif, imagé, syntaxe relâchée, toléré à l'oral (ex. : \emph{se balader}, \emph{bouffer}, \emph{filer / se casser}).
\end{itemize}
\end{definition}

\begin{exemplebox}[Glissement de registre : le verbe « manger »]
\begin{center}
\begin{tabularx}{\linewidth}{|l|X|X|X|}\hline
\textbf{Registre} & \textbf{Verbe} & \textbf{Exemple} & \textbf{Contexte} \\ \hline
Soutenu & s'alimenter, prendre un repas & « Les convives \emph{prirent leur repas} dans le salon. » & Roman, discours officiel \\ \hline
Standard & manger, déjeuner, dîner & « Nous \emph{avons mangé} à midi. » & Journal, résumé, lettre administrative \\ \hline
Familier & bouffer, grailler, se goinfrer & « J'ai \emph{bouffé} un sandwich sur le pouce. » & Discussion entre amis, SMS \\ \hline
\end{tabularx}
\end{center}
\legende{Le sens référentiel (l'action de se nourrir) est identique ; la connotation (le ton, l'image) change radicalement.}
\end{exemplebox}

\courssub{Les variétés régionales du français : le cas camerounais}

Le français est une langue \cle{pluricentrique} : elle possède plusieurs centres normatifs (Paris, Québec, Abidjan, Yaoundé, Kinshasa, etc.). Au Cameroun, pays bilingue officiel (français/anglais) doté de plus de 250 langues nationales, le français a développé une \cle{variété régionale} normative, le \textbf{français camerounais}, enrichie d'emprunts lexicaux indispensables pour nommer des réalités culturelles, culinaires, sociales ou institutionnelles locales.

\begin{savaistu}[Le français camerounais : une norme à part entière]
Le français est langue officielle du Cameroun aux côtés de l'anglais. Mais loin d'être une simple copie du français hexagonal, il s'est \emph{indigénisé}. Plus de 250 langues nationales (bassa, bamileke, ewondo, douala, fulfulde, etc.) ont versé des centaines d'emprunts dans le lexique courant : \emph{ndolé, mbanga, njangi, ngondo, mvet, fongbé, tamtam, chefferie, notable, élite, deuxième bureau, etc.} Ces mots ne sont pas des « fautes » ; ce sont des \cle{lexèmes de remplissage} (ils comblent un vide lexical du français standard) ou des \cle{lexèmes de nuance} (ils portent une culture spécifique). Ils sont consignés dans des dictionnaires de référence comme le \emph{Dictionnaire du français camerounais} (Nguemba, Kouega, etc.).
\end{savaistu}

\begin{popfigure}
\begin{tikzpicture}[scale=0.95, every node/.style={font=\small}]
    % Couleurs
    \colorlet{std}{popblue!70}
    \colorlet{reg}{popgreen!70}
    \colorlet{regis}{poporange!70}
    \colorlet{camf}{poppurple!70}
    
    % 1. Le grand cercle : Langue Française
    \draw[thick, std, fill=std!8] (0,0) circle (5cm);
    \node[std, align=center, font=\bfseries\large] at (0, 3.8) {LANGUE FRANÇAISE (Système commun)};
    \node[std, align=center, font=\tiny] at (0, 3.2) {Grammaire, syntaxe, lexique de base partagés};
    
    % 2. Variétés Régionales (Cercle intermédiaire)
    \draw[thick, reg, fill=reg!8] (0,0) circle (3.5cm);
    \node[reg, align=center, font=\bfseries] at (0, 2.4) {VARIÉTÉS RÉGIONALES (Normes nationales)};
    \node[reg, align=center] at (0, 1.7) {Français camerounais, ivoirien, québécois, belge, suisse, etc.};
    \node[reg, align=center] at (0, 1.0) {Emprunts aux langues locales : \emph{ndolé, mbanga, njangi, chefferie}};
    
    % 3. Registres de langue (Cercle central)
    \draw[thick, regis, fill=regis!8] (0,0) circle (2cm);
    \node[regis, align=center, font=\bfseries] at (0, 0.7) {REGISTRES DE LANGUE};
    \node[regis, align=center] at (0, 0.0) {Soutenu \textbullet\ Standard \textbullet\ Familier};
    \node[regis, align=center] at (0, -0.7) {Variation situationnelle (contexte, interlocuteurs)};
    
    % 4. Camfranglais (Ellipse transversale chevauchant Régional et Registres)
    \draw[thick, dashed, camf, fill=camf!5, rotate=15] (0,-1.2) ellipse (3.2cm and 1.6cm);
    \node[camf, align=center, font=\bfseries, rotate=15] at (0, -1.2) {CAMFRANGLIS};
    \node[camf, align=center, font=\small, rotate=15] at (0, -1.8) {Français + Anglais + Langues nationales + Argot};
    \node[camf, align=center, font=\small, rotate=15] at (0, -2.4) {Variété hybride, orale, identitaire (jeunes urbains)};
    
    % Flèche de convergence vers le centre (Résumé)
    \draw[->, popdark, thick, shorten >=2pt] (0, -3.5) -- (0, -2.8) node[below, font=\bfseries\small, popdark] {CONTRACTION / RÉSUMÉ : Neutralisation $\to$ FRANÇAIS STANDARD};
    
    % Légendes externes
    \node[std, anchor=west, font=\tiny] at (5.3, 4.0) {Noyau dur invariant};
    \node[reg, anchor=west, font=\tiny] at (5.3, 1.5) {Légitimité nationale};
    \node[regis, anchor=west, font=\tiny] at (5.3, -0.5) {Choix énonciatif};
    \node[camf, anchor=west, font=\tiny] at (5.3, -2.5) {Mélange codique};
\end{tikzpicture}
\legende{Architecture des variétés lexicales : du système commun (français) aux variétés régionales (camerounais), traversées par les registres (soutenu/standard/familier). Le camfranglais est une variété hybride transversale (chevauchement). En contraction, on converge vers le centre : le français standard.}
\end{popfigure}

\courssub{Zoom sur le lexique camerounais : réalités culturelles et emprunts}

Certains mots camerounais n'ont pas d'équivalent exact en français standard sans périphrase. Ils sont donc \emph{intraduisibles} en un seul mot. Le résumeur doit les connaître pour les \cle{reformuler} (périphrase standard) et non les \emph{supprimer} (perte de sens).

\begin{center}
\begin{tabularx}{\linewidth}{|l|X|X|X|}\hline
\textbf{Camerounisme (Lexème)} & \textbf{Origine / Nature} & \textbf{Sens précis (Définition)} & \textbf{Reformulation standard (Résumé)} \\ \hline
\emph{Ndolé} & Emprunt (bassa/duala) & Plat national : feuilles de \emph{Vernonia amygdalina} (amères), lavées, cuites avec arachide, viande/poisson, crevettes. & « Plat traditionnel à base de feuilles amères et d'arachide » ou conserver \emph{ndolé} entre guillemets si terme culturel clé. \\ \hline
\emph{Mbanga} (ou \emph{mbongo}) & Emprunt (bassa/duala) & Sauce/condiment à base d'écorces de \emph{mbongo} (arbre), donnant une couleur noire et un goût fumé ; aussi le plat (poisson mbanga). & « Sauce traditionnelle noire à base d'écorces » ou « plat en sauce noire ». \\ \hline
\emph{Njangi} (ou \emph{tontine}) & Emprunt (bamileke/bassa) / Français régional & Association d'épargne et de crédit rotatif (tontine) très structurée, pilier de l'économie informelle. & « Tontine », « association d'épargne rotative », « groupement d'épargne ». \\ \hline
« \emph{Il est allé au village} » & Syntagme figé (Calque structurel) & Ne signifie pas un simple déplacement géographique. Sens local : \textbf{il est mort} (retour aux ancêtres, terre natale). Euphémisme culturel majeur. & « Il est décédé », « il est mort », « il a rejoint ses ancêtres » (selon ton). \\ \hline
\emph{Deuxième bureau} & Métaphore administrative détournée & Lieu officieux (souvent un bar, un domicile) où se traitent les vrais dossiers, par opposition au « premier bureau » (officiel, lent, bureaucratique). & « Lieu de décision informel », « coulisses administratives », « réseau d'influence ». \\ \hline
\emph{Appeler} (au sens de « payer ») & Extension de sens (Calque anglais \emph{to call} / langues locales) & « Il a \emph{appelé} la bouteille » = Il a payé la tournée / offert le verre. & « Il a payé la tournée », « il a offert à boire ». \\ \hline
\end{tabularx}
\end{center}
\legende{Principaux camerounismes lexicaux et leurs équivalents standard pour la contraction. La colonne « Reformulation » montre le travail de \emph{paraphrase} nécessaire.}

\courssub{Le lexique colonial et administratif dans \emph{Le Vieux nègre et la médaille}}

Le roman de Ferdinand Oyono est une satire féroce de l'administration coloniale. Le lexique utilisé n'est pas neutre ; il charrie l'idéologie coloniale. Le contracteur doit identifier ces termes pour les restituer avec la distance critique appropriée (guillemets, verbes d'opinion, reformulation neutre).

\begin{exemplebox}[Lexique colonial / administratif chez Oyono]
\begin{itemize}
    \item \textbf{Commandant} : Représentant de l'autorité coloniale (chef de circonscription). \emph{Résumé :} « l'administrateur colonial », « le chef de l'administration ».
    \item \textbf{Évolué} : Terme colonial péjoratif/paternaliste désignant l'Africain « assimilé », scolarisé, coupé de ses racines. \emph{Résumé :} « l'Africain occidentalisé » (terme neutre), « l'assimilé » (avec guillemets), « l'élite lettrée ».
    \item \textbf{Médaille} : Symbole de la reconnaissance coloniale, outil de domination symbolique (acheter la loyauté). \emph{Résumé :} « la décoration coloniale », « la médaille honorifique », « le symbole de la collaboration ».
    \item \textbf{Indigène / Naturel} : Catégories juridiques coloniales (Code de l'indigénat). \emph{Résumé :} « le sujet colonial », « l'Africain colonisé ».
    \item \textbf{Chef de canton / Chef de village} : Auxiliaires de l'administration, souvent imposés. \emph{Résumé :} « le chef traditionnel nommé par l'administration », « l'auxiliaire de l'administration ».
\end{itemize}
\end{exemplebox}

\begin{attention}[Piège : Confondre variété régionale et faute de français]
\emph{Erreur fréquente :} Croire que « ndolé », « njangi », « il est allé au village » ou le camfranglais sont des « fautes de français » à corriger en les biffant.
\emph{Réalité linguistique :} Ce sont des \textbf{variétés légitimes} du français camerounais (norme régionale) ou des variétés sociales (camfranglais). Elles ont leur grammaire, leur lexique, leur fonction identitaire.
\emph{Règle du résumé :} On ne « corrige » pas, on \textbf{reformule en standard}. Le résumé est un texte \emph{métalinguistique} : il parle \emph{du} texte source en français standard. On neutralise la marque locale/familière pour garantir l'intelligibilité universelle (le correcteur du Probatoire/Bac n'est pas nécessairement camerounais).
\end{attention}

\courssub{Le camfranglais : variété hybride et identité jeune}

Le \cle{camfranglais} (ou \emph{Camfranglais}, \emph{Frananglais}) est une variété linguistique urbaine, parlée principalement par les jeunes à Yaoundé, Douala, Bafoussam. C'est un \cle{méta-langage} identitaire : il mélange structurellement le français (base syntaxique), l'anglais (apports lexicaux massifs), les langues nationales (bassa, ewondo, bamileke, fulfulde) et l'argot.

\begin{exoresolu}[Analyse et reformulation d'énoncés camfranglais]
\textbf{Énoncé source (oral, camfranglais) :}
« \emph{Mon gars, le prof a trop \textbf{coincé} le \textbf{contrôle}. Il a \textbf{dash} tout le monde zéro. On a \textbf{walk} pour aller \textbf{chiller} au \textbf{quartier latin}. »}

\textbf{Analyse lexicale :}
\begin{itemize}
    \item \emph{Mon gars} (familier/argot) $\to$ Mon ami / Mon camarade.
    \item \emph{Coincé} (argot : rendre difficile, faire rater) $\to$ Rendu difficile / A fait échouer.
    \item \emph{Contrôle} (standard, gardé).
    \item \emph{Dash} (pidgin/argot : donner, attribuer ; ici « coller » une note) $\to$ Attribué / Mis / Infligé.
    \item \emph{Walk} (anglais \emph{walk} : marcher, partir) $\to$ Sont partis / Ont marché.
    \item \emph{Chiller} (anglais \emph{chill} : se détendre) $\to$ Se détendre / Se reposer / Prendre un verre.
    \item \emph{Quartier Latin} (Toponyme local : quartier estudiantin/festif de Yaoundé) $\to$ Le quartier étudiant / Un lieu de détente estudiantin.
\end{itemize}

\textbf{Reformulation en français standard (style résumé) :}
\begin{quote}
« Son camarade lui raconte que l'enseignant a rendu l'épreuve excessivement difficile, sanctionnant tous les élèves par la note de zéro. Le groupe a alors quitté les lieux pour se détendre dans un quartier animé fréquenté par les étudiants. »
\end{quote}

\textbf{Observation :} La reformulation \textbf{ne conserve pas le ton familier/argotique} (pas de « mon gars », « walk », « chill »). Elle \textbf{conserve l'information factuelle} (prof, contrôle difficile, zéro pour tous, départ, détente, lieu). C'est la \cle{neutralisation du registre}.
\end{exoresolu}

\courssub{Choix du registre dans le résumé : la neutralisation obligatoire}

Le résumé (ou contraction) est un texte \textbf{hétérogène} : il dépend d'un texte source mais doit être \textbf{autonome} et \textbf{universellement lisible}. La règle d'or est l'\cle{uniformisation au français standard}.

\begin{methode}[Neutraliser les marques de registre et régionales en 3 étapes]
\begin{enumerate}
    \item \textbf{Repérage (Lecture active) :} Surligner \emph{tous} les mots relevant du familier, de l'argot, du camfranglais, du lexique régional (camerounismes) ou du lexique idéologique (colonial) dans le texte source.
    \item \textbf{Qualification :} Pour chaque mot surligné, identifier sa variété : \emph{Familier ? Régional (Cameroun) ? Camfranglais ? Colonial / Idéologique ? Soutenu (trop littéraire) ?}
    \item \textbf{Substitution standard :} Remplacer par l'équivalent le plus neutre, le plus précis et le plus concis du français standard (dictionnaire de référence : \emph{Larousse}, \emph{Robert}, \emph{Dictionnaire du français camerounais} pour vérifier le sens exact du régionalisme).
    \begin{itemize}
        \item Familier $\to$ Standard : \emph{bouffer} $\to$ \emph{manger} ; \emph{se casser} $\to$ \emph{partir} ; \emph{fric} $\to$ \emph{argent}.
        \item Camerounisme (intraduisible en 1 mot) $\to$ Périphrase standard : \emph{njangi} $\to$ \emph{tontine / association d'épargne} ; \emph{allé au village} $\to$ \emph{décédé}.
        \item Camfranglais $\to$ Standard : \emph{dash} $\to$ \emph{donner} ; \emph{walk} $\to$ \emph{partir}.
        \item Colonial/Idéologique $\to$ Neutre critique : \emph{évolué} $\to$ \emph{occidentalisé / assimilé (guillemets)} ; \emph{indigène} $\to$ \emph{colonisé / sujet colonial}.
    \end{itemize}
\end{enumerate}
\end{methode}

\courssub{Application guidée : réécrire en français standard}

Appliquons la méthode à cinq phrases caractéristiques de l'oralité ou de l'écrit informel camerounais.

\begin{exercice}[Mise en situation : Neutralisation lexicale]
Réécrire les phrases suivantes en \textbf{français standard}, apte à figurer dans un résumé formel (examen Probatoire/Baccalauréat). Conserver l'information exacte, supprimer la marque de registre ou de variété.
\begin{enumerate}
    \item « Le vieux a \textbf{quadrillé} son \textbf{mbanga} hier soir, il a trop \textbf{chopé} la \textbf{go}. »
    \item « Depuis qu'il a \textbf{dash} la \textbf{médaille}, il se la \textbf{pète} : il dit qu'il est \textbf{évolué} maintenant. »
    \item « Mon \textbf{pote} a \textbf{walk} au \textbf{village} le mois dernier ; sa \textbf{njangi} a \textbf{coincé} pour les \textbf{funérailles}. »
    \item « Le \textbf{commandant} a \textbf{appelé} la \textbf{bière} au \textbf{deuxième bureau} pour \textbf{ngambler} les \textbf{notables}. »
    \item « Elle a \textbf{go} au \textbf{marché} \textbf{acheter} le \textbf{ndolé} et le \textbf{poisson braisé} pour le \textbf{baptême}. »
\end{enumerate}
\end{exercice}

\begin{corrige}[Exercice n°1 : Neutralisation lexicale]
\begin{enumerate}
    \item \textbf{Source :} « Le vieux a \emph{quadrillé} son \emph{mbanga} hier soir, il a trop \emph{chopé} la \emph{go}. »
    \begin{itemize}
        \item \emph{Quadrillé} (argot : bien garni, copieux) $\to$ \textbf{bien garni / copieux}.
        \item \emph{Mbanga} (camerounisme : sauce noire / plat) $\to$ \textbf{sauce traditionnelle (mbanga) / plat en sauce noire}.
        \item \emph{Chopé} (argot/camfranglais : attrapé, eu, mangé avec appétit / séduit — ambiguïté contextuelle) $\to$ \textbf{mangé avec appétit / consommé} (contexte culinaire privilégié).
        \item \emph{Go} (camfranglais/argot : fille, compagne / ou bouche, repas) $\to$ \textbf{compagne / petite amie} ou \textbf{repas}. \textbf{Contexte :} « sa compagne » ou « le repas ».
    \end{itemize}
    \textbf{Standard :} « L'ancien a \textbf{bien garni sa sauce traditionnelle hier soir ; il a mangé avec appétit / il a régalé sa compagne.} » (Le contexte tranche l'ambiguïté de « chopé la go »).

    \item \textbf{Source :} « Depuis qu'il a \emph{dash} la \emph{médaille}, il se la \emph{pète} : il dit qu'il est \emph{évolué} maintenant. »
    \begin{itemize}
        \item \emph{Dash} (camfranglais : donner, attribuer ; ici « recevoir/décrocher ») $\to$ \textbf{reçu / obtenu}.
        \item \emph{Se la pète} (familier : se vanter, faire le fier) $\to$ \textbf{se vanter / affiche une fierté excessive}.
        \item \emph{Évolué} (colonial/idéologique) $\to$ \textbf{« évolué » (guillemets) / occidentalisé / assimilé}.
    \end{itemize}
    \textbf{Standard :} « Depuis qu'il a \textbf{reçu la médaille}, il \textbf{se vante / affiche une fierté outrancière}, se présentant comme un \textbf{« évolué » / un occidentalisé}. »

    \item \textbf{Source :} « Mon \emph{pote} a \emph{walk} au \emph{village} le mois dernier ; sa \emph{njangi} a \emph{coincé} pour les \emph{funérailles}. »
    \begin{itemize}
        \item \emph{Pote} (familier : ami) $\to$ \textbf{ami / camarade}.
        \item \emph{Walk} (camfranglais : partir, aller) $\to$ \textbf{est parti / s'est rendu}.
        \item \emph{Au village} (euphémisme culturel : mourir) $\to$ \textbf{est décédé / est mort}.
        \item \emph{Njangi} (camerounisme : tontine) $\to$ \textbf{tontine / association d'épargne}.
        \item \emph{Coincé} (argot jeune : ici sens positif « a contribué, a payé, a honoré sa part » par extension de « mettre la main à la poche » / « débloquer ») $\to$ \textbf{a cotisé / a financé / a honoré sa contribution}.
        \item \emph{Funérailles} (standard, gardé).
    \end{itemize}
    \textbf{Standard :} « Mon \textbf{ami est décédé} le mois dernier ; sa \textbf{tontine a financé / assuré les frais des funérailles.} »

    \item \textbf{Source :} « Le \emph{commandant} a \emph{appelé} la \emph{bière} au \emph{deuxième bureau} pour \emph{ngambler} les \emph{notables}. »
    \begin{itemize}
        \item \emph{Commandant} (colonial/admin) $\to$ \textbf{administrateur / chef de l'administration}.
        \item \emph{Appelé} (camerounisme : payer, offrir) $\to$ \textbf{offert / payé / réglé}.
        \item \emph{Deuxième bureau} (camerounisme : lieu informel de pouvoir) $\to$ \textbf{coulisses / lieu de concertation informel / en dehors des circuits officiels}.
        \item \emph{Ngambler} (camfranglais/argot : \emph{ngamb} = parler, discuter, négocier, flatter, soudoyer) $\to$ \textbf{négocier avec / influencer / soudoyer}.
        \item \emph{Notables} (standard/admin, gardé ou « chefs traditionnels / élites locales »).
    \end{itemize}
    \textbf{Standard :} « L'\textbf{administrateur a offert des boissons lors d'une concertation informelle pour \textbf{négocier / influencer les notables.}} »

    \item \textbf{Source :} « Elle a \emph{go} au \emph{marché} \emph{acheter} le \emph{ndolé} et le \emph{poisson braisé} pour le \emph{baptême}. »
    \begin{itemize}
        \item \emph{A go} (camfranglais : \emph{go} = aller, partir ; structure \emph{sujet + go + lieu}) $\to$ \textbf{est allée / s'est rendue}.
        \item \emph{Ndolé} (camerounisme : plat/feuille) $\to$ \textbf{ndolé (terme culturel conservé entre guillemets) / feuilles de vernonia / plat traditionnel ndolé}.
        \item \emph{Poisson braisé} (standard régional, accepté dans standard large, ou \textbf{poisson grillé}).
        \item \emph{Baptême} (standard, gardé).
    \end{itemize}
    \textbf{Standard :} « Elle \textbf{s'est rendue au marché pour acheter du ndolé et du poisson grillé pour le baptême.} »
\end{enumerate}
\end{corrige}

\begin{aretenir}[Règle d'or pour la contraction]
Dans un résumé (contraction de texte), \textbf{tout mot relevant du familier, de l'argot, du camfranglais, du lexique régional non lexicalisé dans le standard national, ou du lexique idéologique (colonial, politique) doit être remplacé par son équivalent en français standard}.
\begin{itemize}
    \item On \textbf{ne cite} le terme source (entre guillemets) que s'il est \emph{l'objet même de l'étude} (ex. : « Oyono utilise le terme "évolué" pour dénoncer... »).
    \item On \textbf{conserve} le terme régional (ndolé, njangi) \textbf{entre guillemets} s'il n'a pas d'équivalent monosémique court, en l'expliquant brièvement si nécessaire.
    \item Le \textbf{verbe de reformulation} (dire, affirmer, montrer, décrire, ironiser) doit être au \textbf{standard} (pas de « il balance que », « il dit que »).
\end{itemize}
\end{aretenir}

\courssub{Synthèse : de la diversité lexicale à l'unité du résumé}

\begin{center}
\begin{tabularx}{\linewidth}{|l|X|X|}\hline
\textbf{Niveau d'analyse} & \textbf{Question opératoire} & \textbf{Action du contracteur} \\ \hline
\textbf{Mot (Lexique)} & Ce mot est-il standard ? Familier ? Régional (Cameroun) ? Camfranglais ? Colonial ? & \textbf{Substituer} par l'équivalent standard (périphrase si intraduisible). \\ \hline
\textbf{Syntagme / Expression figée} & Y a-t-il un sens culturel caché ? (ex. « allé au village ») & \textbf{Décoder} le sens réel $\to$ \textbf{Reformuler} littéralement en standard. \\ \hline
\textbf{Phrase / Énoncé} & Quel est le registre global ? (Ironique ? Performatif ? Constatif ?) & \textbf{Neutraliser} la tonalité (ironie $\to$ constat ; performatif $\to$ rapport d'acte). \\ \hline
\textbf{Texte (Global)} & La reformulation assure-t-elle l'intelligibilité pour tout francophone ? & \textbf{Relire} en « mode correcteur extérieur » : supprimer toute trace de « moi, ici, maintenant, entre nous ». \\ \hline
\end{tabularx}
\end{center}
\legende{Grille de neutralisation lexicale et énonciative : du mot au texte, le fil conducteur est la convergence vers le français standard.}

Cette maîtrise des variétés lexicales est le socle technique qui nous permettra, dans la section suivante, de \textbf{dégager la structure argumentative du texte} : car on ne peut analyser la logique d'un raisonnement (thèse, arguments, connecteurs) si l'on reste prisonnier du vocabulaire subjectif, local ou familier de l'énonciateur. La neutralisation lexicale est la condition \emph{sine qua non} de l'objectivation du résumé.

\coursec{La structure de la phrase : outil de la contraction}

\courssub{Rappel morphosyntaxique : de la phrase simple à la phrase complexe}

Maîtriser la contraction, c'est d'abord maîtriser la syntaxe. Réduire un texte sans le trahir suppose de savoir manipuler les briques de la phrase : les propositions. En Terminale technique, vous devez opérer des transformations syntaxiques conscientes et justifiées.

\begin{definition}[Propositions et types de phrases]
\begin{itemize}
    \item \textbf{Phrase simple :} Une seule proposition (un seul verbe conjugué). Ex. : \emph{« L'ingénieur calcule la résistance. »}
    \item \textbf{Phrase complexe :} Au moins deux propositions. Elles peuvent être :
    \begin{itemize}
        \item \cle{Indépendantes} (juxtaposées ou coordonnées) : \emph{« L'ingénieur calcule, le technicien exécute. »}
        \item \cle{Subordonnées} (dépendantes d'une principale) : \emph{« L'ingénieur calcule \textbf{pendant que} le technicien exécute. »} (subordonnée circonstancielle de temps).
    \end{itemize}
\end{itemize}
\end{definition}

\courssub{Transformations syntaxiques au service de la réduction}

La contraction repose sur trois opérations morphosyntaxiques majeures. Elles permettent de passer du \emph{parlant} (verbal, explicite, long) à l'\emph{écrit dense} (nominal, implicite, court).

\begin{methode}[Trois opérations de réduction syntaxique]
\begin{enumerate}
    \item \textbf{Nominalisation (verbe $\rightarrow$ nom) :} Transformer une proposition en groupe nominal (GN).
    \begin{exemplebox}[Nominalisation]
    \emph{Proposition :} « \textbf{Parce que} le coût des matières premières \textbf{augmente}, l'entreprise \textbf{réduit} sa production. » (2 verbes)
    \emph{GN :} « \textbf{L'augmentation} du coût des matières premières \textbf{entraîne} la réduction de la production. » (1 verbe + 2 GN).
    \end{exemplebox}
    \item \textbf{Suppression de l'accessoire (élision) :} Supprimer les propositions subordonnées relatives explicatives, les incises, les compléments de phrase redondants.
    \begin{exemplebox}[Suppression]
    \emph{Source :} « Le rapport, \textbf{qui a été publié hier par le MINESEC}, confirme la réforme. »
    \emph{Contracté :} « Le rapport \textbf{publié hier par le MINESEC} confirme la réforme. » (Participe passé adjectivé) ou « Le \textbf{dernier rapport du MINESEC} confirme la réforme. »
    \end{exemplebox}
    \item \textbf{Fusion / Intégration (coordination $\rightarrow$ subordination ou participe) :} Lier deux idées fortes en une seule structure.
    \begin{exemplebox}[Fusion]
    \emph{Source :} « Les élèves sont motivés. \textbf{Et} ils réussissent mieux. »
    \emph{Contracté :} « \textbf{Motivés,} les élèves réussissent mieux. » (Gérondif/Participe) ou « La motivation des élèves \textbf{favorise} leur réussite. » (Nominalisation).
    \end{exemplebox}
\end{enumerate}
\end{methode}

\begin{attention}[Piège de la sur-nominalisation]
Un excès de noms abstraits (\emph{« la mise en œuvre de la réalisation de l'évaluation... »}) alourdit le style et floute le sens. \textbf{Alternance :} Gardez le verbe pour l'action principale, nominalisez les circonstances et les causes secondaires.
\end{attention}

\coursec{Énoncés constatifs et performatifs}

\courssub{Distinction fondamentale (Théorie des actes de langage : Austin, Searle)}

Tout énoncé n'a pas la même fonction langagière. En contraction, il est vital de distinguer ce qui \emph{constate} (informations à garder) de ce qui \emph{agit} (souvent à reformuler ou condenser).

\begin{definition}[Énoncé constatif vs Performatif]
\begin{itemize}
    \item \textbf{Énoncé \cle{constatif} :} Décrit un état de fait, vrai ou faux. Il \textbf{rend compte} de la réalité. \emph{Marqueurs :} verbes déclaratifs (affirmer, constater, décrire, expliquer), indicatif présent/passé, absence de « je » performatif.
    \item \textbf{Énoncé \cle{performatif} :} \textbf{Réalise l'action} qu'il énonce par le seul fait d'être proféré (dans des conditions de réception appropriées). Il \textbf{transforme} la réalité sociale.
\end{itemize}
\end{definition}

\begin{center}
\begin{tabularx}{\linewidth}{|l|X|X|}\hline
\rowcolor{popblueL}
\textbf{Critère} & \textbf{Constatif} & \textbf{Performatif (Explicite)} \\ \hline
\textbf{Verbe type} & Constater, dire, penser, écrire, affirmer & Promettre, jurer, nommer, baptiser, déclarer, ordonner, voter, pardonner \\ \hline
\textbf{Structure canonique} & Sujet + Verbe (indicatif) + Complément & \textbf{Je (Nous)} + \textbf{Verbe performatif (1\textsuperscript{re} pers. sing./plur. indicatif présent)} + Complément \\ \hline
\textbf{Test de vérité} & Vrai / Faux & \textbf{Heureux / Inheureux} (conditions de félicité : autorité, contexte, sincérité) \\ \hline
\textbf{Exemple} & « Le directeur \emph{a signé} le décret. » & « Je \emph{déclare} la session ouverte. » / « Nous \emph{promettons} d'agir. » \\ \hline
\end{tabularx}
\end{center}

\courssub{Identification et traitement en contraction}

\begin{methode}[Repérer et contracter les performatifs]
\begin{enumerate}
    \item \textbf{Repérage :} Cherchez les verbes de la \emph{liste performative} (déclarer, décider, nommer, interdire, autoriser, promettre, s'engager) à la 1\textsuperscript{re} personne présent.
    \item \textbf{Analyse :} L'énoncé engage-t-il la responsabilité de l'énonciateur ? Change-t-il l'état du monde (nomination, engagement contractuel) ?
    \item \textbf{Contraction :}
    \begin{itemize}
        \item \textbf{Performatif explicite $\rightarrow$ Constatif rapporté :} « Je promets de venir » $\rightarrow$ « Il s'engage à venir » / « Sa promesse de venir ».
        \item \textbf{Performatif implicite (ordre, conseil) $\rightarrow$ Infinitif / Nominalisation :} « Je vous ordonne de partir » $\rightarrow$ « Ordre de départ » / « Il ordonne le départ ».
    \end{itemize}
\end{enumerate}
\end{methode}

\begin{exemplebox}[Application à un texte administratif/technique]
\textbf{Texte :} « \emph{Le Ministre arrête :} Article 1 : \emph{Je nomme} M. X directeur. Article 2 : \emph{Le présent arrêté sera publié} au Journal Officiel. »
\textbf{Analyse :} « J'arrête » et « Je nomme » = performatifs explicites (pouvoir réglementaire). « Sera publié » = constatif futur (passif) ou performatif implicite (décision de publication).
\textbf{Contraction :} « Par arrêté ministériel, M. X est nommé directeur ; publication au Journal Officiel. » (Les performatifs deviennent des noms d'actes : \emph{arrêté, nomination, publication}).
\end{exemplebox}

\coursec{Ironie et paradoxe : figures de pensée}

\courssub{Définition et repérage}

Ces figures disent le contraire de ce qu'elles montrent (ironie) ou unissent l'inconciliable (paradoxe). Elles sont fréquentes dans les textes d'opinion, éditoriaux, essais (corpus Bac). Ne pas les repérer conduit à un contresens total en contraction.

\begin{definition}[Ironie]
Procédé qui consiste à \textbf{dire le contraire de ce que l'on pense}, en comptant sur l'intelligence du lecteur (complicité) pour décoder le sens réel. Elle repose sur un \textbf{décalage} entre l'énoncé littéral et le contexte/les connaissances partagées.
\begin{itemize}
    \item \textbf{Marqueurs :} Ton surévalué (éloges excessifs), guillemets ironiques, antiphrase (mot utilisé en sens inverse), hyperboles, contexte contradictoire.
    \item \textbf{Effet :} Dénonciation, critique virulente, rire jaune, complicité lecteur/auteur.
\end{itemize}
\end{definition}

\begin{definition}[Paradoxe]
Affirmation \textbf{apparemment contradictoire, absurde ou contre-intuitive}, mais qui \textbf{révèle une vérité profonde} sur le réel après réflexion. Ce n'est pas une erreur logique, mais une tension heuristique.
\begin{itemize}
    \item \textbf{Structure :} « A est non-A » / « Le pire est le meilleur » / « Il faut perdre pour gagner ».
    \item \textbf{Effet :} Provocation intellectuelle, ouverture de la pensée, remise en cause des évidences.
\end{itemize}
\end{definition}

\begin{exemplebox}[Ironie dans un éditorial technique]
« \emph{Quelle magnifique idée} que de former des ingénieurs sans laboratoires ! \emph{On admire} la clairvoyance de nos décideurs qui \emph{offrent} à la jeunesse le luxe du chômage intellectuel. »
\begin{itemize}
    \item \textbf{Littéral :} Éloge de l'absence de labos, admiration pour décideurs, luxe du chômage.
    \item \textbf{Sens réel (Ironie) :} Critique violente du manque d'équipements, incompétence des décideurs, drame du chômage.
    \item \textbf{Contraction :} \textbf{Ne gardez JAMAIS le littéral.} Reformulez le sens réel : « L'absence de laboratoires compromet la formation d'ingénieurs employables, révélant l'irresponsabilité politique. »
\end{itemize}
\end{exemplebox}

\begin{exemplebox}[Paradoxe dans un essai sur la technique]
« \emph{C'est en acceptant l'échec que l'ingénieur construit la réussite.} » (Paradoxe : Échec $\neq$ Réussite).
\textbf{Explication :} L'essai-erreur, le prototypage, l'analyse des modes de défaillance (AMDEC) sont le cœur de la démarche qualité.
\textbf{Contraction :} « L'analyse de l'échec fonde la fiabilité technique. » (On garde la tension conceptuelle mais on explicite le mécanisme).
\end{exemplebox}

\begin{attention}[Piège Bac : Ironie vs Mensonge / Paradoxe vs Contradiction]
\begin{itemize}
    \item \textbf{Ironie} = Intention critique partagée (complicité). \textbf{Mensonge} = Tromperie unilatérale.
    \item \textbf{Paradoxe} = Vérité profonde masquée. \textbf{Contradiction} = Erreur logique (A et non-A en même temps, même sens).
\end{itemize}
En contraction : \textbf{Traduisez l'ironie en énoncé direct (constatif critique). Explicitez le paradoxe en relation causale ou dialectique.}
\end{attention}

\coursec{Les variétés lexicales du français}

\courssub{Diaphasie, Diastratie, Diatopie : le français en mouvement}

Le choix des mots n'est jamais neutre. Il situe l'auteur (qui parle ?), le destinataire (à qui ?), le contexte (où ? quand ?) et le registre (comment ?). En contraction, le vocabulaire technique, les néologismes, les emprunts sont des \textbf{unités de sens non négociables} (on les garde), tandis que le vocabulaire familier ou régional peut être standardisé.

\begin{propriete}[Trois axes de variation lexicale (G. Louis, M. Francard)]
\begin{enumerate}
    \item \textbf{Variation \cle{diaphasique} (selon la situation / registre) :} Même locuteur, contextes différents.
    \begin{itemize}
        \item \emph{Registre soutenu :} « \textbf{Acquérir} des compétences », « \textbf{Dysfonctionnement} ».
        \item \emph{Registre courant :} « \textbf{Apprendre} un métier », « \textbf{Panne} ».
        \item \emph{Registre familier / argotique :} « \textbf{Choper} les bases », « \textbf{Planté} ».
    \end{itemize}
    \item \textbf{Variation \cle{diastratique} (selon le groupe social / sociolecte) :} Jargon professionnel, langue de jeunesse, langue administrative.
    \begin{itemize}
        \item \emph{Jargon technique (BTP) :} « \textbf{Ferraillage}, \textbf{banchage}, \textbf{fluage}, \textbf{fléau} ».
        \item \emph{Admin. camerounaise :} « \textbf{Arrêté}, \textbf{note de service}, \textbf{visa}, \textbf{ordonnateur} ».
    \end{itemize}
    \item \textbf{Variation \cle{diatopique} (selon l'espace géographique) :} Français standard vs variétés régionales.
    \begin{itemize}
        \item \emph{Cameroun (Français camerounais) :} « \textbf{Appeler} (pour téléphoner), \textbf{descendre} (pour aller au centre-ville), \textbf{le courant} (l'électricité), \textbf{un taxi} (moto-taxi \emph{benskin} ou voiture), \textbf{le quartier} (arrondissement) ».
        \item \emph{Camfranglais (jeunes urbains) :} « \emph{Ça chauffe} (c'est difficile), \emph{je wanda} (je marche/je vais), \emph{le djo} (l'ami), \emph{la go} (la fille) ».
    \end{itemize}
\end{enumerate}
\end{propriete}

\begin{methode}[Gestion du lexique en contraction]
\begin{enumerate}
    \item \textbf{Identifiez la variété :} Technique ? Administrative ? Familière ? Régionale (Cameroun) ?
    \item \textbf{Conservez (Sacralisé) :} Termes techniques précis (\emph{alternance, CQP, BTS, fluage, AMDEC}), noms propres (MINESEC, ENEO, OHADA), chiffres, sigles officiels.
    \item \textbf{Standardisez (Reformulez) :} Familier $\rightarrow$ Courant/Soutenu (\emph{« ça chauffe » $\rightarrow$ « c'est ardu »}). Régional marqué $\rightarrow$ Standard (\emph{« le courant » $\rightarrow$ « l'électricité » ; « descendre en ville » $\rightarrow$ « se rendre au centre-ville »}).
    \item \textbf{Expliquez si nécessaire :} Camfranglais ou néologisme récent $\rightarrow$ gardez le terme \emph{entre guillemets} + définition brève si crucial pour l'argument.
\end{enumerate}
\end{methode}

\begin{exemplebox}[Variétés dans un texte sur l'apprentissage dual]
« Les \emph{apprentis} de \emph{Mboa Soudure} (Camfranglais : \emph{Mboa} = pays/patrie) ont \emph{géré} (familier : réussi) leur CQP. Le \emph{patron} (familier/argotique pour employeur) a \emph{validé} le \emph{processus} (terme technique/admin). »
$\rightarrow$ \textbf{Contraction :} « Les apprentis de l'atelier Mboa Soudure ont obtenu leur CQP ; l'employeur a validé la procédure. » (\emph{Mboa} gardé comme nom propre, \emph{géré/patron} standardisés, \emph{processus/CQP} gardés).
\end{exemplebox}

\coursec{Dégager la structure argumentative du texte}

Après avoir exploré les \emph{variétés lexicales du français} et compris comment le choix des mots ancre un texte dans un registre, une époque ou une aire culturelle, nous devons maintenant nous pencher sur l'architecture même du discours. Une contraction de texte réussie ne se contente pas de raccourcir des phrases : elle préserve le \textbf{squelette logique} de l'argumentation. C'est pourquoi, avant toute reformulation, il est impératif de savoir \emph{radiographier} un texte pour en extraire la thèse, les arguments et leur articulation.

\courssub{Rappel : les composantes de l'argumentation}

Tout texte argumentatif, qu'il s'agisse d'un éditorial, d'un essai philosophique ou d'une lettre de motivation, repose sur quatre piliers indissociables. Les identifier avec précision est la première étape de toute contraction rigoureuse.

\begin{definition}[Structure argumentative]
La \cle{structure argumentative} est l'organisation logique qui relie la thèse aux arguments qui la soutiennent. Elle assure la cohérence du discours et guide le lecteur vers l'adhésion au point de vue défendu.
\end{definition}

\begin{itemize}
    \item Le \cle{thème} : le sujet général abordé (ex. : « l'éducation technique au Cameroun »). Il est neutre, il ne prend pas parti.
    \item La \cle{thèse} : le \textbf{point de vue défendu} par l'auteur sur ce thème (ex. : « L'éducation technique est le levier principal de l'industrialisation du Cameroun »). C'est la réponse de l'auteur à la problématique implicite ou explicite.
    \item Les \cle{arguments} : les \textbf{raisons logiques} avancées pour justifier la thèse (ex. : « Elle forme une main-d'œuvre qualifiée immédiatement employable », « Elle réduit la dépendance aux expatriés »).
    \item Les \cle{exemples} (ou illustrations) : des \textbf{faits concrets, statistiques, anecdotes, citations} qui donnent du poids aux arguments (ex. : « Le taux d'insertion des BTS en Génie civil est de 85\% selon le MINEFOP 2023 »).
\end{itemize}

\begin{attention}[Hiérarchie des idées]
Ne confondez jamais l'argument (l'idée force, abstraite et générale) et l'exemple (la preuve, singulière et concrète). Dans une contraction, \textbf{l'argument est sacro-saint, l'exemple est souvent condensé ou supprimé} s'il n'apporte pas une information statistique cruciale.
\end{attention}

\courssub{Les étapes types d'un texte argumentatif}

La très grande majorité des textes soumis au Probatoire / Baccalauréat technique suivent une architecture classique en trois temps. La repérer permet de segmenter le texte pour le traiter par blocs.

\begin{propriete}[Architecture standard d'un texte argumentatif]
\begin{enumerate}
    \item \textbf{Introduction} : Accroche $\rightarrow$ Présentation du thème $\rightarrow$ \cle{Problématique} (souvent sous forme de question) $\rightarrow$ Annonce de la thèse et du plan (parfois implicite).
    \item \textbf{Développement} : Suite de \cle{paragraphes arguments}. Chaque paragraphe = 1 argument principal + ses illustrations + connecteurs de liaison. Structure interne : \emph{Thèse partielle $\rightarrow$ Explication $\rightarrow$ Exemple $\rightarrow$ Mini-conclusion / Transition}.
    \item \textbf{Conclusion} : \cle{Bilan} (réponse à la problématique / reformulation de la thèse) $\rightarrow$ \cle{Ouverture} (perspective, conséquence future, question nouvelle).
\end{enumerate}
\end{propriete}

\begin{methode}[Repérer la structure en 3 lectures]
\begin{enumerate}
    \item \textbf{Lecture 1 (globale) :} Quel est le thème ? Quelle est la thèse ? (Souvent dans l'intro ou la conclusion).
    \item \textbf{Lecture 2 (structurale) :} Numérotez les paragraphes. Identifiez l'idée-force de chacun (souvent la 1ère phrase). S'agit-il d'un argument pour la thèse, d'un contre-argument (concession), d'une illustration ?
    \item \textbf{Lecture 3 (fine) :} Soulignez les connecteurs logiques. Repérez la modalisation (verbes modaux, adverbes, tournures impersonnelles) qui signale l'engagement de l'auteur.
\end{enumerate}
\end{methode}

\courssub{Les marqueurs de l'argumentation : la piste des connecteurs et de l'énonciateur}

L'argumentation ne se cache pas ; elle se \emph{signale} par des indices linguistiques précis. En Terminale technique, vous devez savoir les nommer et les exploiter pour tracer le cheminement de la pensée.

\begin{center}
\begin{tabularx}{\linewidth}{|l|X|X|}\hline
\rowcolor{popblueL}
\textbf{Type de marqueur} & \textbf{Exemples} & \textbf{Fonction argumentative} \\ \hline
\cle{Connecteurs logiques} (organisateurs) & \emph{D'abord, en premier lieu, ensuite, de plus, par ailleurs, en effet, car, parce que, donc, par conséquent, c'est pourquoi, cependant, toutefois, au contraire, en revanche, en définitive, pour conclure} & Structurent le raisonnement : énumération, cause, conséquence, opposition, concession, conclusion. \\ \hline
\cle{Présence de l'énonciateur} & \emph{Je pense que, à mon avis, nous devons, il est clair que, force est de constater que, il est indéniable que} & Marque l'engagement subjectif (je, nous) ou l'objectivation (il est... que). \\ \hline
\cle{Modalisation} (degré de certitude) & \emph{Certainement, sans doute, probablement, peut-être, il faut que, il est nécessaire que, devoir, pouvoir (obligation/possibilité)} & Nuance la force de l'argument : affirmation forte, probabilité, hypothèse, obligation morale. \\ \hline
\cle{Marqueurs de la concession} & \emph{Il est vrai que, admettons que, certes, on pourrait objecter que... mais} & Intègre un point de vue adverse pour mieux le réfuter (argument d'autorité retourné). \\ \hline
\end{tabularx}
\end{center}

\begin{exemplebox}[Analyse de marqueurs]
\textbf{Texte :} « \emph{Certes}, l'enseignement général forme des esprits critiques. \emph{Cependant}, \emph{il est indéniable que} l'enseignement technique \emph{répond directement} aux besoins du marché de l'emploi. \emph{En effet}, les entreprises recrutent massivement des techniciens supérieurs. \emph{Par conséquent}, orienter les élèves vers les filières techniques \emph{apparaît comme une nécessité}. »

\begin{itemize}
    \item \emph{Certes / Cependant} : Concession $\rightarrow$ Opposition (structure dialectique).
    \item \emph{Il est indéniable que} : Modalisation forte (certitude objective).
    \item \emph{En effet} : Introduction d'une justification (argument d'autorité / fait).
    \item \emph{Par conséquent} : Conséquence logique $\rightarrow$ Thèse finale.
\end{itemize}
\end{exemplebox}

\courssub{Distinguer l'essentiel de l'accessoire : l'art du tri sélectif}

C'est le cœur du métier de la contraction. Le texte source contient toujours une part significative de « matière grasse » : exemples développés, répétitions rhétoriques, citations d'autorité, digressions, transitions lourdes. Votre mission : \textbf{garder la substantifique moelle}.

\begin{methode}[Technique du « Surlignage Différencié » (avant rédaction)]
\begin{enumerate}
    \item \textbf{Surlignez en JAUNE} la \cle{Thèse} (souvent 1 phrase).
    \item \textbf{Surlignez en VERT} les \cle{Arguments principaux} (1 par paragraphe de développement, souvent la phrase-topic).
    \item \textbf{Barrez au STYLO ROUGE} (ou entourez pour suppression) :
    \begin{itemize}
        \item Les exemples longs (récits, anecdotes, études de cas détaillées).
        \item Les citations (sauf si l'auteur \emph{est} l'argument, ex. : « Selon X... » $\rightarrow$ garder « X affirme que... »).
        \item Les répétitions / reformulations internes (« En d'autres termes », « C'est-à-dire »).
        \item Les transitions lourdes (« Nous venons de voir que..., nous allons maintenant examiner... »).
        \item Les concessions secondaires qui ne sont pas le cœur du débat.
    \end{itemize}
    \item \textbf{Conservez} les connecteurs logiques essentiels (donc, mais, or, car) pour maintenir la cohérence du résumé.
\end{enumerate}
\end{methode}

\begin{exemplebox}[Piège majeur : l'exemple « frappant »]
\textbf{Texte source (80 mots) :} « L'éducation technique est vitale. Prenons le cas de Jean, jeune BTS Électrotechnique de Douala. Après son stage chez ENEO, il a été embauché CDI à 22 ans, achète sa moto, soutient sa famille. C'est la preuve que ce secteur insère. »
\textbf{Mauvais résumé (élève) :} « L'éducation technique est vitale. Jean, BTS Électrotechnique à Douala, embauché chez ENEO après son stage, achète une moto et soutient sa famille. » (L'exemple a mangé l'argument !)
\textbf{Bon résumé :} « L'éducation technique assure une insertion professionnelle rapide, comme en témoigne l'embauche massive des techniciens supérieurs par les grandes entreprises. » (L'argument est gardé, l'exemple condensé en fait général).
\end{exemplebox}

\begin{exemplebox}[Anatomie d'un éditorial type (600 mots)]
\begin{center}
\begin{tabular}{|l|c|c|}\hline
\rowcolor{popblueL}
\textbf{Composante} & \textbf{Volume approx.} & \textbf{Devenir en contraction (1/4 = 150 mots)} \\ \hline
Thèse (intro + conclusion) & 50 mots (8\%) & \textbf{Intégrale} (reformulée) \\ \hline
Arguments principaux (3 à 4) & 150 mots (25\%) & \textbf{Intégrale} (reformulée, fusionnée) \\ \hline
Exemples / Illustrations / Citations & 250 mots (42\%) & \textbf{Supprimés ou condensés} en 1 phrase globale \\ \hline
Transitions / Connecteurs / Redites / Accroche / Ouverture & 150 mots (25\%) & \textbf{Supprimés} (remplacés par connecteurs courts) \\ \hline
\end{tabular}
\end{center}
\legende{Règle empirique : jusqu'à 60\% du texte source est « accessoire » au sens strict de la contraction argumentative.}
\end{exemplebox}

\courssub{Construction d'un schéma argumentatif : l'arbre de la pensée}

Avant d'écrire une seule ligne du résumé, le meilleur élève trace un \textbf{schéma argumentatif} (carte mentale / arbre). Cela force à hiérarchiser et à vérifier qu'aucun argument n'est oublié.

\begin{popfigure}
\begin{tikzpicture}[
    mindmap,
    grow cyclic,
    every node/.style={concept, execute at begin node=\hskip0pt},
    concept color=popblue,
    level 1/.style={level distance=4.5cm, sibling angle=120, concept color=popblue!70},
    level 2/.style={level distance=3.5cm, sibling angle=60, concept color=popgreen!70},
    level 3/.style={level distance=3cm, sibling angle=40, concept color=poporange!70},
    text=white,
    font=\small\bfseries
]
\node[root concept, concept color=popdark, text width=3.5cm, align=center] {THÈSE\\L'éducation technique\\moteur de l'industrialisation}
    child[concept color=popblue] { node[concept, text width=3cm, align=center] {ARGUMENT 1\\Insertion pro.\\rapide}
        child { node[concept, text width=2.5cm, align=center] {Ex: Taux insertion BTS 85\%} }
        child { node[concept, text width=2.5cm, align=center] {Ex: Partenariats entreprises} }
    }
    child[concept color=poppurple] { node[concept, text width=3cm, align=center] {ARGUMENT 2\\Réponse aux\\besoins locaux}
        child { node[concept, text width=2.5cm, align=center] {Ex: Formation maintenance\\industrielle} }
        child { node[concept, text width=2.5cm, align=center] {Ex: BTP, Agro-alimentaire} }
    }
    child[concept color=popgreen] { node[concept, text width=3cm, align=center] {ARGUMENT 3\\Réduction chômage\\jeunes}
        child { node[concept, text width=2.5cm, align=center] {Ex: Statistiques MINJEC} }
        child { node[concept, text width=2.5cm, align=center] {Ex: Auto-emploi facilité} }
    };
\end{tikzpicture}
\legende{Schéma argumentatif type (Arbre) : Thèse centrale, 3 arguments branches, exemples feuilles. Ce schéma guide la rédaction du résumé : on ne rédige que les nœuds « Thèse » et « Arguments ».}
\end{popfigure}

\begin{aretenir}[Règle d'or du schéma]
Si vous ne pouvez pas dessiner l'arbre en 2 minutes, vous n'avez pas compris la structure du texte. Ne commencez \textbf{jamais} la rédaction du résumé sans ce schéma validé.
\end{aretenir}

\courssub{Application guidée : article de presse camerounais sur l'éducation technique}

Appliquons la méthode sur un texte simulé, représentatif des corpus du Bac Technique (éditorial de \emph{Cameroon Tribune} ou \emph{Mutations}, ~450 mots).

\begin{exoresolu}[Dégager la structure : « Technique : l'urgence d'un pari gagnant »]
\textbf{Énoncé :} Voici un extrait d'éditorial. 1) Identifiez thème, thèse, arguments, exemples. 2) Tracez le schéma argumentatif. 3) Proposez une contraction à 110 mots ($\pm$ 10\%).

\textbf{Texte source (432 mots) :}
« L'avenir économique du Cameroun se joue aujourd'hui dans les ateliers de nos lycées techniques. \emph{C'est le constat lucide} que dresse le dernier rapport de la Banque mondiale sur l'emploi des jeunes. \emph{En effet}, alors que le taux de chômage des diplômés de l'enseignement général dépasse 30\%, les titulaires de BTS et de DUT affichent un taux d'insertion à 6 mois supérieur à 75\%. \emph{Certes}, certains parents persistent à voir dans les filières techniques une voie de garage pour élèves en difficulté. \emph{Cependant}, \emph{il faut souligner que} les programmes ont été réformés en 2018 pour intégrer l'alternance et les TIC. \emph{Prenons l'exemple} du Lycée Technique de New-Bell : sur la session 2023, 92\% des élèves en Maintenance Industrielle ont signé un contrat avant même la proclamation des résultats. \emph{De plus}, l'État a lancé le programme « Trois ans pour l'emploi » qui finance l'équipement de 50 ateliers modernes. \emph{Néanmoins}, le défi reste la formation des formateurs. \emph{En définitive}, faire de l'enseignement technique la priorité nationale n'est plus une option, \emph{c'est une nécessité vitale} pour l'émergence à l'horizon 2035. »

\tcblower
\textbf{Solution détaillée :}

\textbf{1. Analyse structurelle}
\begin{itemize}
    \item \textbf{Thème :} L'éducation technique au Cameroun / Emploi des jeunes.
    \item \textbf{Thèse :} L'enseignement technique doit devenir la priorité nationale (nécessité vitale pour l'émergence 2035).
    \item \textbf{Argument 1 (Économique/Statistique) :} Insertion professionnelle massive et rapide (75\% vs 30\% chômage général). \emph{Exemple :} Rapport Banque Mondiale ; Lycée New-Bell 92\%.
    \item \textbf{Argument 2 (Pédagogique/Structurel) :} Modernisation des cursus (Réforme 2018, alternance, TIC). \emph{Exemple :} Programme « 3 ans pour l'emploi », 50 ateliers.
    \item \textbf{Argument 3 (Réponse à l'objection) :} Levée du préjugé « voie de garage » par la réforme et les résultats concrets.
    \item \textbf{Argument 4 (Limite/Realisme) :} Défi de la formation des formateurs (concession lucide).
\end{itemize}

\textbf{2. Schéma argumentatif (simplifié pour la rédaction)}
\begin{center}
\textsc{Thèse} : Priorité nationale à l'enseignement technique $\rightarrow$ Émergence 2035.
$\downarrow$
\begin{tabular}{ll}
\textsc{Arg 1} : Insertion pro. supérieure (Stats BM, New-Bell). & (Garder chiffre clé 75\%/92\%) \\
\textsc{Arg 2} : Adaptation offre/besoins (Réforme 2018, Alternance, TIC). & (Garder idée force) \\
\textsc{Arg 3} : Soutien étatique concret (Programme 3 ans, 50 ateliers). & (Garder chiffre 50) \\
\textsc{Arg 4} : Vigilance nécessaire (Formation formateurs). & (Garder comme condition de succès)
\end{tabular}
\end{center}

\textbf{3. Contraction proposée (108 mots)}
« Le rapport de la Banque mondiale confirme que l'enseignement technique, avec 75\% d'insertion à six mois contre 30\% de chômage pour le général, est le levier de l'émergence 2035. La réforme de 2018 (alternance, TIC) et le programme « Trois ans pour l'emploi » (50 ateliers modernes) alignent l'offre de formation sur les besoins réels, comme le prouve les 92\% d'embauche précoce à New-Bell. Lever le préjugé de la « voie de garage » et former les formateurs restent les conditions \emph{sine qua non} pour faire de cette filière la priorité nationale absolue. »

\textbf{Commentaire de la correction :}
\begin{itemize}
    \item Thèse reformulée en incipit.
    \item Arg 1 + Exemples fusionnés : stats globales + exemple local condensé en preuve unique.
    \item Arg 2 + Arg 3 fusionnés (réforme + moyens étatiques = adaptation).
    \item Arg 4 (concession) gardé comme ouverture/condition.
    \item Connecteurs : « confirme que », « et », « comme le prouve », « restent ».
    \item Vocabulaire technique conservé : « insertion », « alternance », « émergence », « sine qua non ».
\end{itemize}
\end{exoresolu}

\courssub{Du texte complet au résumé : l'entonnoir de la contraction}

Visualisons le processus global. La contraction n'est pas une coupe aléatoire, c'est un \textbf{entonnage progressif} de l'information.

\begin{popfigure}
\begin{tikzpicture}[
    node distance=1.2cm and 2cm,
    box/.style={draw=popblue, thick, rounded corners, fill=popblue!10, text width=4.5cm, align=center, minimum height=1.8cm, font=\small},
    arrow/.style={-Stealth, thick, popdark},
    labelbox/.style={font=\footnotesize\itshape, text=popmuted}
]
\node[box, align=center] (source){\textbf{TEXTE SOURCE}\\(~450 mots)\\\emph{Thème, Thèse, Args, Exemples, Citations, Redites, Transitions, Ouverture}};
\node[labelbox, below=0.1cm of source] {Lecture analytique + Surlignage différencié};
\node[box, below=1.8cm of source, align=center] (essentiel){\textbf{IDÉES ESSENTIELLES (Brut)}\\(~150 mots)\\\emph{Thèse + Args principaux (phrases-topic) + Chiffres clés + Concession majeure}};
\node[labelbox, below=0.1cm of essentiel] {Rédaction / Reformulation + Liaisons logiques};
\node[box, below=1.8cm of essentiel, align=center] (resume){\textbf{RÉSUMÉ / CONTRACTION FINAL}\\(110 mots $\pm$ 10\%)\\\emph{Texte autonome, cohérent, personnel, respectant la contrainte de longueur}};
\node[labelbox, below=0.1cm of resume] {Vérification : Fidélité + Nombre de mots + Corrections};
\draw[arrow] (source) -- (essentiel);
\draw[arrow] (essentiel) -- (resume);
% Entonnoir visuel derrière
\begin{scope}[on background layer]
\draw[poporange!50, line width=12cm] (source.north) -- (resume.south);
\draw[poporange, thick] (source.north east) -- (resume.south east);
\draw[poporange, thick] (source.north west) -- (resume.south west);
\end{scope}
\end{tikzpicture}
\legende{L'entonnoir de la contraction : du texte brut (large) aux idées essentielles (filtrées) vers le résumé final (dense et reformulé). La perte de volume (mots) ne doit jamais signifier perte de structure (logique).}
\end{popfigure}

\begin{savaistu}[Au Probatoire / Baccalauréat Technique]
L'épreuve de contraction (souvent couplée au résumé) est notée sur trois critères majeurs, pondérés environ ainsi :
\begin{enumerate}
    \item \textbf{Fidélité aux idées essentielles (40\%) :} Thèse présente ? Arguments principaux tous présents ? Pas de contresens ? Pas d'ajout personnel (opinion, connaissance extérieure) ?
    \item \textbf{Reformulation personnelle et qualité de la langue (40\%) :} Pas de « copier-coller » de phrases entières (tolérance : termes techniques, chiffres, noms propres). Phrases correctes, connecteurs fluides, vocabulaire précis.
    \item \textbf{Respect de la contrainte quantitative (20\%) :} Nombre de mots imposé $\pm$ 10\%. Dépasser = pénalité ; trop court = idées manquantes.
\end{enumerate}
\emph{Astuce :} Comptez vos mots \textbf{avant} la copie propre. Une ligne manuscrite $\approx$ 8-10 mots. Ciblez le nombre de lignes.
\end{savaistu}

\coursec{Techniques de réduction et reformulation des idées essentielles}

\courssub{De l'analyse à la rédaction : les opérations de reformulation}

Une fois le schéma argumentatif validé, il faut \textbf{réécrire}. La reformulation n'est pas une traduction mot à mot, c'est une \emph{re-création} syntaxique et lexicale qui préserve le sens exact.

\begin{propriete}[Les 4 piliers de la reformulation réussie]
\begin{enumerate}
    \item \textbf{Substitution lexicale (Synonymie / Hyperonymie) :} Remplacer un mot par son synonyme précis ou son hyperonyme (terme plus général).
    \begin{exemplebox}[Substitution]
    \emph{Source :} « L'entreprise \textbf{recrute} massivement des \textbf{soudeurs}. »
    \emph{Reformulation :} « La société \textbf{embauche} en nombre des \textbf{techniciens de soudage}. » (Synonymie) ou « Le secteur \textbf{absorbe} de la \textbf{main-d'œuvre qualifiée}. » (Hyperonymie).
    \end{exemplebox}
    \item \textbf{Transposition grammaticale (Changement de nature) :} Verbe $\leftrightarrow$ Nom $\leftrightarrow$ Adjectif $\leftrightarrow$ Adverbe.
    \begin{exemplebox}[Transposition]
    \emph{Source :} « \textbf{Parce qu'il a plu}, le sol \textbf{est devenu} boueux. »
    \emph{Reformulation :} « \textbf{La pluie} a \textbf{boueux} le sol. » (Nominalisation cause) ou « \textbf{Pluvieux}, le sol est boueux. » (Adjectivation).
    \end{exemplebox}
    \item \textbf{Changement de voix / diathèse :} Actif $\leftrightarrow$ Passif $\leftrightarrow$ Réfléchi $\leftrightarrow$ Impersonnel.
    \begin{exemplebox}[Diathèse]
    \emph{Source :} « \textbf{Le ministre a signé} l'arrêté. » (Actif)
    \emph{Reformulation :} « L'arrêté \textbf{a été signé par le ministre} » (Passif) ou « \textbf{La signature} de l'arrêté \textbf{par le ministre} » (Nominalisation) ou « \textbf{Il a été signé} l'arrêté ministériel. » (Impersonnel).
    \end{exemplebox}
    \item \textbf{Restructuration syntaxique (Fusion / Division / Permutation) :} Modifier l'ordre des propositions, fusionner deux phrases en une, ou diviser une phrase trop dense.
    \begin{exemplebox}[Fusion]
    \emph{Source :} « Le BTS est un diplôme national. \textbf{Il} permet l'insertion. \textbf{Il} est reconnu par l'État. »
    \emph{Reformulation :} « Diplôme national reconnu par l'État, le BTS favorise l'insertion. » (1 phrase, 2 participes passés adjectivés).
    \end{exemplebox}
\end{enumerate}
\end{propriete}

\courssub{La gestion des connecteurs : fluidité et logique}

En contraction, les connecteurs longs (\emph{« par ailleurs, il convient de noter que »}) sont remplacés par des connecteurs brefs (\emph{« de plus, or, car, donc, mais »}) ou par la ponctuation (deux-points, point-virgule, parenthèses).

\begin{center}
\begin{tabular}{|l|l|l|}\hline
\rowcolor{popblueL}
\textbf{Connecteur source (lourd)} & \textbf{Connecteur cible (bref)} & \textbf{Relation} \\ \hline
Il est important de souligner que / Force est de constater que & \textbf{Fait que / Certes} (ou suppression) & Insistance / Évidence \\ \hline
En ce qui concerne / Quant à & \textbf{Sur / Pour} + Nom / \textbf{Quant à} & Thématisation \\ \hline
Cela dit / Cela étant / Toutefois / Néanmoins & \textbf{Mais / Or / Pourtant} & Opposition / Concession \\ \hline
Par conséquent / C'est pourquoi / Aussi & \textbf{Donc / Aussi / Dès lors} & Conséquence \\ \hline
En effet / Car / Parce que / Puisque & \textbf{Car / Car / Puisque} (début phrase) / \textbf{:} (deux-points) & Cause / Justification \\ \hline
De plus / En outre / Par ailleurs & \textbf{Et / De plus / Aussi} & Addition \\ \hline
\end{tabular}
\end{center}

\begin{attention}[Respect de la contrainte de mots : le comptage rigoureux]
\begin{methode}[Compter et ajuster]
\begin{enumerate}
    \item Rédigez un premier jet sans compter (viser ~120\% de la cible).
    \item Comptez \textbf{chaque mot} (mots-outils inclus : \emph{le, de, et, à, dans, pour...}). Hyphénés = 1 mot (\emph{porte-parole}), chiffres = 1 mot (\emph{2025}), sigles = 1 mot (\emph{BTS}).
    \item \textbf{Trop long ?} Coupez les adjectifs qualificatifs non techniques, remplacez les relatives par des participes, fusionnez les phrases.
    \item \textbf{Trop court ?} Vérifiez qu'aucun argument n'est manquant. Développez la thèse ou un argument clé avec un détail technique/chiffre gardé.
\end{enumerate}
\end{methode}
\end{attention}

\coursec{Application au \emph{Vieux Nègre et la Médaille} : Lectures méthodiques 5, 6 et Exposés}

\courssub{Contexte de l'œuvre (Rappel synthétique)}

\emph{Le Vieux Nègre et la Médaille} (1956) de \textbf{Ferdinand Oyono} (Cameroun) est une nouvelle satirique (récit à thèse) dénonçant l'aliénation coloniale et la collaboration indigène. Le protagoniste, \textbf{Meka}, vieux chef traditionnel, reçoit la « Médaille de la Reconnaissance française » pour sa fidélité. L'ironie dramatique structure le texte : Meka croit à l'honneur, le lecteur sait la duperie.

\courssub{Lecture méthodique n°5 : Structure argumentative implicite du récit}

Bien que ce soit une fiction, le texte porte une \textbf{thèse argumentative implicite} : « La collaboration coloniale est une aliénation qui vide l'Africain de sa substance humaine et historique. »

\begin{exoresolu}[Analyse structurelle du récit (Extrait : Cérémonie de remise de médaille)]
\textbf{Texte visé :} Chapitre final (Arrivée de l'Administrateur, discours, remise de médaille, réaction de Meka, chute).

\textbf{1. Thème :} La cérémonie officielle de décoration d'un chef africain par l'administration coloniale.
\textbf{2. Thèse implicite (Auteur) :} La médaille symbolise la servitude volontaire et le mépris colonial masqué en honneur.
\textbf{3. Arguments (distribués dans la narration) :}
\begin{itemize}
    \item \textbf{Arg 1 (Mise en scène du pouvoir) :} L'Administrateur incarne l'autorité lointaine, méprisante, pressée (regarde sa montre). \emph{Marqueurs :} Verbes d'action rapide, champ lexical du temps/horloge, distance physique.
    \item \textbf{Arg 2 (Langage vide / Performatif creux) :} Le discours de l'Administrateur est stéréotypé, déconnecté de la réalité de Meka. \emph{Marqueurs :} Performatifs rituels (\emph{« Je vous remets... », « La France vous remercie »}), lieux communs (\emph{« loyal serviteur »}).
    \item \textbf{Arg 3 (Aliénation du héros) :} Meka intériorise la valeur coloniale, oublie sa propre histoire (guerres, résistance oubliée). \emph{Marqueurs :} Monologue intérieur ému, fierté mal placée, oubli des siens (Kelara, ses fils).
    \item \textbf{Arg 4 (Ironie dramatique / Chute) :} La médaille est un objet clinquant, sans valeur réelle (pas de pension, pas de terre), Meka meurt peu après / reste « vieux nègre ».
\end{itemize}
\textbf{4. Schéma argumentatif narratif :}
\begin{center}
\begin{tikzpicture}[
    mindmap, grow cyclic, concept color=popdark,
    level 1/.style={level distance=4cm, sibling angle=120, concept color=popblue!80},
    level 2/.style={level distance=3cm, sibling angle=50, concept color=poporange!80},
    text=white, font=\small\bfseries
]
\node[root concept, text width=3cm, align=center] {THÈSE IMPLICITE\\Collaboration = Aliénation}
    child { node[concept, text width=2.5cm, align=center] {ARG 1\\Mise en scène\\pouvoir colonial}
        child { node[concept, text width=2cm] {Montre, distance, hâte} }
    }
    child { node[concept, text width=2.5cm, align=center] {ARG 2\\Discours\\performatif vide}
        child { node[concept, text width=2cm] {Stéréotypes, « loyal »} }
    }
    child { node[concept, text width=2.5cm, align=center] {ARG 3\\Conscience\\aliénée de Meka}
        child { node[concept, text width=2cm] {Fierté, oubli histoire} }
    }
    child { node[concept, text width=2.5cm, align=center] {ARG 4\\Chute ironique\\Objet vain}
        child { node[concept, text width=2cm] {Médaille clinquante, mort} }
    };
\end{tikzpicture}
\end{center}
\end{exoresolu}

\courssub{Lecture méthodique n°6 : Figures de pensée (Ironie, Paradoxe) et Variétés lexicales}

Oyono utilise l'ironie comme \textbf{arme critique principale} et le paradoxe pour révéler l'absurde colonial.

\begin{exemplebox}[Ironie dramatique et situationnelle]
\textbf{Passage :} Meka revêt son « habit du dimanche » (veste usée, pantalon court), se pare de la médaille. L'Administrateur : « \emph{Vous êtes un bel exemple de fidélité.} »
\begin{itemize}
    \item \textbf{Ironie situationnelle :} Meka pense être honoré ; le lecteur sait qu'il est exhibé comme trophée.
    \item \textbf{Antiphrase (Ironie verbale) :} « \emph{Bel exemple} » = triste marionnette. « \emph{Fidélité} » = soumission.
    \item \textbf{Contraction :} « La cérémonie tourne à la mascarade : l'Administrateur, pressé, épingle une médaille clinquante sur la veste usée de Meka, le qualifiant ironiquement de « bel exemple » de soumission. »
\end{itemize}
\end{exemplebox}

\begin{exemplebox}[Paradoxe colonial]
\textbf{Passage :} Meka a combattu pour les Allemands (contre les Français), puis pour les Français (contre les Allemands). Il reçoit une médaille \emph{française} pour sa « fidélité ».
\begin{itemize}
    \item \textbf{Paradoxe :} La « fidélité » récompensée est en réalité une \textbf{versatilité imposée par la force}. Le traître d'hier est le héros d'aujourd'hui.
    \item \textbf{Contraction :} « Le paradoxe de la « fidélité » récompensée masque une collaboration forcée, Meka ayant servi tour à tour des puissances ennemies. »
\end{itemize}
\end{exemplebox}

\begin{exemplebox}[Variétés lexicales : Français standard vs Français camerounais / Oralité]
\begin{itemize}
    \item \textbf{Français administratif (Performatifs) :} « \emph{Je vous décerne}, \emph{La République vous honore} ». $\rightarrow$ Garder en citation brève ou nominaliser (\emph{« la décoration », « l'hommage officiel »}).
    \item \textbf{Oralité / Français populaire camerounais (Meka, Kelara) :} « \emph{Mon vieux}, \emph{le blanc}, \emph{la médaille ça brille}, \emph{il a mangé l'argent} ».
    \item \textbf{Traitement en contraction :} Standardiser l'oralité (\emph{« le blanc » $\rightarrow$ « l'Administrateur/le colon » ; « il a mangé l'argent » $\rightarrow$ « il a détourné des fonds »}) sauf si l'analyse porte sur le style (exposé).
\end{itemize}
\end{exemplebox}

\courssub{Exposé n°1 : « La satire de la collaboration dans \emph{Le Vieux Nègre et la Médaille} » (Plan type 10-15 min)}

\begin{methode}[Structure de l'exposé oral / écrit]
\begin{enumerate}
    \item \textbf{Accroche :} Citation choc (ex. : « \emph{La médaille, c'est du toc qui brille pour les yeux des nègres.} » - Kelara) + Présentation œuvre/auteur/contexte (1956, Cameroun, pré-indépendance).
    \item \textbf{Problématique :} Comment Oyono utilise-t-il la forme narrative (ironie, paradoxe) pour démonter le mythe de la « mission civilisatrice » et de la « collaboration fructueuse » ?
    \item \textbf{Plan annoncé :}
    \begin{enumerate}
        \item \textbf{Une mise en scène du pouvoir colonial :} L'Administrateur, figure du mépris bureaucratique (temps, distance, performatifs vides).
        \item \textbf{La construction de l'aliénation :} Meka, du statut de sujet à l'objet de décor (intériorisation, oubli identitaire).
        \item \textbf{Le triomphe de l'ironie :} La chute (mort, inutilité de la médaille) comme verdict final de l'Histoire.
    \end{enumerate}
    \item \textbf{Développement :} Appuis textuels précis (citations courtes, numéros de page/édition). Analyse lexicale (champ lexical : horloge/métal vs chair/sang/terre), syntaxique (phrases nominales pour l'admin, phrases verbales/vivantes pour Kelara).
    \item \textbf{Conclusion :} Bilan : La médaille = symbole de l'aliénation. Ouverture : Actualité de la critique (néocolonialisme, décorations honorifiques, mémoire historique).
\end{enumerate}
\end{methode}

\courssub{Exposé n°2 : « Langage, pouvoir et résistance : étude stylistique » (Plan type)}

\begin{methode}[Focus Langue / Style]
\begin{enumerate}
    \item \textbf{Introduction :} Langue comme enjeu de pouvoir. Thèse : Oyono détourne la langue du colonisateur (français standard, performatifs) pour en faire une arme de résistance (ironie, oralité, proverbes).
    \item \textbf{I. La langue du pouvoir : performative, figée, déshumanisante.}
    \begin{itemize}
        \item Performatifs institutionnels (\emph{« Je déclare, Je remets »}) $\rightarrow$ Magie verbale qui ne change rien au réel.
        \item Nominalisation abstraite (\emph{« civilisation, fidélité, mérite »}) $\rightarrow$ Masque le concret (vol, travail forcé, impôts).
        \item Syntaxe rigide, impersonnelle.
    \end{itemize}
    \item \textbf{II. La langue de la résistance : vivante, imagée, ancrée.}
    \begin{itemize}
        \item \textbf{Kelara (la femme) :} Voix de la lucidité. Proverbes (\emph{« Le singe vieillit mais ne devient pas gorille »}), ironie directe, vocabulaire concret (\emph{« toc, ferraille, ventre »}).
        \item \textbf{Les fils (Kelara, les nationalistes) :} Silences lourds, regards, langue du corps (crachat, refus de serrer la main).
        \item \textbf{Meka (ambiguïté) :} Bascule de l'oralité fière (début) au français approximatif de la soumission (fin).
    \end{itemize}
    \item \textbf{III. L'ironie narrative comme méta-langage :} Le narrateur (focalisation zéro ironique) commente, juge, décalage permanent entre \emph{dire} (officiel) et \emph{voir} (réel).
    \item \textbf{Conclusion :} Oyono « camerounise » le français (créolisation syntaxique, emprunts culturels) pour dire l'indicible colonial. La contraction de ce texte doit préserver cette tension stylistique.
\end{enumerate}
\end{methode}

\coursec{Entraînement synthèse : Contraction d'un extrait du \emph{Vieux Nègre et la Médaille}}

\begin{exercice}[Sujet type Probatoire - Contraction 100 mots $\pm$ 10\%]
\textbf{Texte source (Extrait adapté, ~380 mots) :}
« L'Administrateur arriva, pressé, sa montre à la main. \emph{C'est avec un plaisir non dissimulé} qu'il lut le décret. « \emph{Meka, au nom de la République française, je vous remets cette médaille.} \emph{Vous êtes un loyal serviteur.} » Meka, les yeux brillants, tendit la poitrine. \emph{Il revoyait les portages, les coups de chicotte, les impôts arrachés à la faim de son peuple.} Mais il pensait : « \emph{Enfin, la France me reconnaît.} » Kelara, au fond de la cour, cracha : « \emph{Du toc ! Du toc pour les yeux des nègres !} » L'Administrateur, déjà reparti, lança par la portière : « \emph{Bonne fête, Chef ! } » Meka resta seul, la médaille froide sur le cœur, le ventre vide. »

\textbf{Consignes :}
\begin{enumerate}
    \item Identifiez : Thème, Thèse implicite (Auteur), 3 Arguments narratifs, Ironie/Paradoxe, Variétés lexicales.
    \item Tracez le schéma argumentatif (Arbre).
    \item Rédigez la contraction (90-110 mots).
\end{enumerate}
\end{exercice}


\begin{corrige}[Exercice - Contraction \emph{Vieux Nègre} - Éléments de correction]
\textbf{1. Analyse structurelle et stylistique}
\begin{itemize}
    \item \textbf{Thème :} Cérémonie de décoration coloniale / Collaboration.
    \item \textbf{Thèse implicite (Oyono) :} La médaille coloniale est une mascarade ironique qui scelle l'aliénation du collaborateur (Meka) et le mépris du colonisateur (Administrateur), sous le regard lucide de la résistance (Kelara).
    \item \textbf{Arg 1 (Pouvoir) :} Administrateur pressé, bureaucrate, performatifs rituels (\emph{« je vous remets »}), départ immédiat. $\rightarrow$ Mépris réel.
    \item \textbf{Arg 2 (Aliénation) :} Meka oublie souffrances (portages, chicotte, impôts) pour fierté illusoire (\emph{« La France me reconnaît »}). $\rightarrow$ Intériorisation du dominateur.
    \item \textbf{Arg 3 (Résistance/Lucidité) :} Kelara : antiphrase (\emph{« toc »}), vérité crue (\emph{« yeux des nègres »}). $\rightarrow$ Contre-discours oral, populaire.
    \item \textbf{Figures :} Ironie dramatique (Meka honoré / lecteur : humilié), Paradoxe (récompense pour souffrances subies), Antiphrase (\emph{« loyal serviteur »}).
    \item \textbf{Lexique :} Admin (standard, performatif) vs Kelara (oral, concret, métaphorique).
\end{itemize}

\textbf{2. Schéma argumentatif}
\begin{center}
\begin{tikzpicture}[mindmap, concept color=popdark, grow cyclic,
    level 1/.style={level distance=3.5cm, sibling angle=120, concept color=popblue!80},
    level 2/.style={level distance=2.5cm, sibling angle=50, concept color=popgreen!80},
    text=white, font=\tiny\bfseries]
\node[root concept, text width=2.5cm, align=center] {THÈSE\\Médaille = Aliénation\\Ironique}
    child { node[concept, text width=2cm, align=center] {ARG 1\\Admin : Mépris\\Bureaucratique}
        child { node[concept, text width=1.5cm] {Montre, Décret, Fuite} }
    }
    child { node[concept, text width=2cm, align=center] {ARG 2\\Meka : Fierté\\Aveuglée}
        child { node[concept, text width=1.5cm] {Oubli chicotte/impôts} }
    }
    child { node[concept, text width=2cm, align=center] {ARG 3\\Kelara : Lucidité\\Résistante}
        child { node[concept, text width=1.5cm] {« Toc », Crachat, Vérité} }
    };
\end{tikzpicture}
\end{center}

\textbf{3. Contraction proposée (98 mots)}
« Pressé, l'Administrateur épingle la médaille sur Meka par un performatif rituel : « loyal serviteur ». Ironie cruelle : Meka, oubliant portages, chicotte et impôts extorqués, croit à la reconnaissance française. Kelara, seule lucide, crache sur ce « toc pour yeux de nègres ». L'Administrateur fuit, lançant un « bonne fête » méprisant. Meka reste seul, médaille froide sur cœur vide, symbole de sa servitude consentie. La satire d'Oyono, par l'ironie dramatique et le choc des langages (bureaucratique vs oralité populaire), dénonce la collaboration comme aliénation mortifère, vidant l'Africain de son histoire et de sa dignité. »
\end{corrige}

\begin{aretenir}[Check-list finale avant l'épreuve]
\begin{itemize}
    \item [ ] \textbf{Analyse} : Thème, Thèse, Args, Exemples, Figures (Ironie/Paradoxe), Variétés lexicales, Connecteurs, Structure phrase.
    \item [ ] \textbf{Schéma} : Arbre dessiné, validé, hiérarchisé.
    \item [ ] \textbf{Rédaction} : Reformulation personnelle (synonymie, transposition, fusion), vocabulaire technique gardé, connecteurs brefs.
    \item [ ] \textbf{Contrôle} : Fidélité (pas de contresens, pas d'oubli d'arg principal), Nombre de mots ($\pm$ 10\%), Orthographe/Grammaire, Ponctuation.
\end{itemize}
\end{aretenir}

\coursec{Techniques de réduction et reformulation des idées essentielles}

Après avoir appris à \emph{dégager la structure argumentative du texte} (section précédente), nous disposons désormais du « squelette » de l'œuvre : nous savons où se trouve la thèse, quels sont les arguments principaux, les exemples et les connecteurs logiques. L'étape suivante, et non la moindre, consiste à \emph{passer de l'analyse à la production} : il faut maintenant rédiger le résumé (ou contraction) en respectant une contrainte de longueur stricte. C'est tout l'objet de cette section : maîtriser les \cle{techniques de réduction} et l'art de la \cle{reformulation} pour produire un texte court, fidèle, cohérent et correct.

\courssub{Les trois techniques fondamentales de réduction}

Réduire un texte ne signifie pas « couper au hasard ». C'est un travail de \emph{sélection} et de \emph{transformation} guidé par la structure argumentative repérée précédemment. Trois opérations techniques, souvent combinées, permettent de gagner de la place sans perdre le sens.

\begin{definition}[La contraction de texte]
La \cle{contraction de texte} (ou résumé) est une réduction fidèle d'un texte source qui en conserve les idées essentielles, la thèse et l'enchaînement logique, dans une longueur imposée (généralement le quart du texte original, avec une tolérance de $\pm 10\,\%$). Elle exclut tout commentaire personnel et toute citation littérale prolongée.
\end{definition}

\begin{center}
\begin{tabularx}{\linewidth}{|l|X|X|}\hline
\rowcolor{popblueL} \textbf{Technique} & \textbf{Principe et mécanisme} & \textbf{Exemple appliqué} \\ \hline
\textbf{1. Suppression (de l'accessoire)} & Éliminer les éléments non indispensables à la thèse : adjectifs qualificatifs, compléments de détail, exemples illustratifs, répétitions, incises, phrases de transition lourdes. & \emph{Source :} « Le \textbf{vieux} chauffeur, \textbf{fatigué par des décennies de labeur}, arrêta son taxi \textbf{jaune et cabossé} au bord de la route. » (20 mots) \par \emph{Cible :} « Le chauffeur arrêta son taxi. » (6 mots) \\ \hline
\textbf{2. Condensation (fusion)} & Remplacer un groupe de mots (syntagme) par un mot unique ou une expression plus dense : nominalisation, pronominalisation, synthèse de termes coordonnés. & \emph{Source :} « Il \textbf{prit son courage à deux mains} et \textbf{se lança dans l'aventure}. » (12 mots) \par \emph{Cible :} « Il \textbf{s'engagea résolument}. » (3 mots) \\ \hline
\textbf{3. Généralisation (abstraction)} & Remplacer une énumération de termes particuliers (hyponymes) par leur hyperonyme (terme générique), ou un fait précis par une idée générale. & \emph{Source :} « Les \textbf{motos-taxis, les taxis, les bus et les cars} encombrent la voirie. » (11 mots) \par \emph{Cible :} « Les \textbf{transports urbains} encombrent la voirie. » (6 mots) \\ \hline
\rowcolor{popgreenL} \textbf{Règle d'or} & \multicolumn{2}{l|}{Combiner les trois techniques \emph{après} avoir isolé les idées-forces. Ne jamais supprimer un argument principal ni un connecteur logique majeur.} \\ \hline
\end{tabularx}
\end{center}

\begin{exemplebox}[Exemple chiffré : la généralisation]
Dans un texte traitant de la mobilité à Douala, on lit : « Les \emph{motos-taxis}, les \emph{taxis} et les \emph{bus} assurent le transport des populations. » (13 mots).
\par \textbf{Généralisation :} « Les \emph{transports urbains} assurent la mobilité des populations. » (8 mots).
\par \textbf{Gain :} 5 mots économisés (environ $38\,\%$ de réduction) pour une information préservée.
\end{exemplebox}

\begin{attention}[Piège : la suppression abusive]
Supprimer un \emph{argument principal} ou un \emph{connecteur logique} (donc, cependant, par conséquent) brise la cohérence du résumé. Exemple : si l'auteur écrit « Il a plu, \textbf{donc} le match est annulé », supprimer « donc » produit « Il a plu. Le match est annulé. » : le lien de causalité disparaît. \textbf{Toujours conserver les opérateurs argumentatifs.}
\end{attention}

\courssub{La reformulation : changer la forme, garder le sens}

La contraction n'est pas une suite de « copier-coller » tronqués. Le correcteur attend une \emph{réécriture personnelle}. La reformulation repose sur deux leviers majeurs : le levier lexical et le levier syntaxique.

\begin{propriete}[Les deux leviers de la reformulation]
\textbf{1. Changer les mots (niveau lexical)}
\begin{itemize}
    \item \textbf{Synonymie / paraphrase :} « Il est \emph{décédé} » $\to$ « Il est \emph{mort} » ; « Il a \emph{mis un terme à ses jours} » $\to$ « Il s'est \emph{donné la mort} ».
    \item \textbf{Nominalisation (verbe $\to$ nom) :} « Le gouvernement \emph{a décidé} de baisser les prix » $\to$ « La \emph{décision} du gouvernement de baisser les prix... »
    \item \textbf{Verbalisation (nom $\to$ verbe) :} « Son \emph{arrivée} a surpris tout le monde » $\to$ « Il \emph{est arrivé}, surprenant tout le monde. »
\end{itemize}
\textbf{2. Changer la construction (niveau syntaxique)}
\begin{itemize}
    \item \textbf{Voix active $\leftrightarrow$ voix passive :} « L'auteur \emph{démontre} la thèse » $\leftrightarrow$ « La thèse \emph{est démontrée} par l'auteur. »
    \item \textbf{Subordonnée $\to$ groupe nominal ou proposition participiale :} « \emph{Parce qu'il pleuvait}, il est rentré » $\to$ « \emph{À cause de la pluie}, il est rentré » ou « La pluie \emph{l'ayant forcé}, il est rentré. »
    \item \textbf{Relative $\to$ adjectif ou apposition :} « Un homme \emph{qui a du courage} » $\to$ « Un homme \emph{courageux} » ; « Meka, \emph{qui est un vieux villageois}, reçoit une médaille » $\to$ « Meka, \emph{vieux villageois}, reçoit une médaille. »
\end{itemize}
\end{propriete}

\begin{exoresolu}[Réduction et reformulation d'un paragraphe (98 mots $\to$ environ 40 mots)]
\textbf{Texte source (extrait fictif, style essai) :}
« \textit{Il est indéniable que la déforestation massive, qui touche aujourd'hui l'ensemble du bassin du Congo, entraîne des conséquences dramatiques et irréversibles sur la biodiversité locale. En effet, les grands singes, tels que les gorilles et les chimpanzés, voient leur habitat naturel se réduire comme peau de chagrin. De plus, les populations riveraines, qui dépendent traditionnellement de la forêt pour leur subsistance, se trouvent démunies face à l'avancée des exploitations forestières industrielles. Par conséquent, la protection de cet écosystème unique apparaît comme une urgence absolue pour les gouvernements de la sous-région.} » (98 mots)

\textbf{Étape 1 : structure argumentative (rappel).}
Thèse : la déforestation du bassin du Congo menace la biodiversité et les populations ; sa protection est urgente.
Arguments : 1. perte d'habitat des grands singes ; 2. précarisation des populations riveraines. Connecteur de conclusion : \emph{Par conséquent}.

\textbf{Étape 2 : application des trois techniques et reformulation.}
\begin{itemize}
    \item \emph{Suppression :} « indéniable », « massive », « aujourd'hui », « dramatiques et irréversibles », « comme peau de chagrin », « traditionnellement », « industrielles », « unique », « absolue ».
    \item \emph{Condensation :} « exploitations forestières industrielles » $\to$ « exploitation forestière » ; « apparaît comme une urgence absolue » $\to$ « s'impose urgemment ».
    \item \emph{Généralisation :} « grands singes, tels que les gorilles et les chimpanzés » $\to$ « grands singes » ; « populations riveraines qui dépendent de la forêt » $\to$ « populations forestières ».
    \item \emph{Reformulation syntaxique :} « Il est indéniable que... entraîne des conséquences sur » $\to$ « ...menace » ; « voient leur habitat se réduire » $\to$ « perdent leur habitat » ; « se trouvent démunies face à » $\to$ « sont privées de leurs ressources par ».
\end{itemize}

\textbf{Résumé proposé (39 mots) :}
\par \textit{« La déforestation du bassin du Congo menace les grands singes en détruisant leur habitat et prive les populations forestières de leurs ressources. Face à l'exploitation forestière, la protection de cet écosystème s'impose urgemment aux gouvernements sous-régionaux. »}

\textbf{Vérification :} fidélité (thèse + 2 arguments + conclusion) ; reformulation (vocabulaire et syntaxe transformés) ; cohérence (connecteurs : \emph{et}, \emph{Face à}) ; longueur (39 mots, proche de la cible de 40) ; langue correcte.
\tcblower
\textbf{Analyse des gains :}
\begin{itemize}
    \item « Il est indéniable que la déforestation massive... entraîne des conséquences dramatiques et irréversibles sur » (17 mots) $\to$ « La déforestation... menace » (5 mots). \textbf{Gain : 12 mots} (suppression + condensation + voix active).
    \item « grands singes, tels que les gorilles et les chimpanzés » (9 mots) $\to$ « grands singes » (2 mots). \textbf{Gain : 7 mots} (généralisation).
    \item « populations riveraines, qui dépendent traditionnellement de la forêt pour leur subsistance, se trouvent démunies face à l'avancée des exploitations forestières industrielles » (22 mots) $\to$ « prive les populations forestières de leurs ressources. Face à l'exploitation forestière » (12 mots). \textbf{Gain : 10 mots} (suppression + condensation + reformulation).
\end{itemize}
\end{exoresolu}

\courssub{Le calcul du nombre de mots : la contrainte quantitative}

Au Probatoire comme au Baccalauréat technique, la consigne impose une longueur précise (par exemple : « Résumez ce texte en 125 mots environ » ou « Réduisez le texte au quart de sa longueur »). Le non-respect de la fourchette est pénalisant.

\begin{methode}[Compter et viser la fourchette ($\pm 10\,\%$)]
\begin{enumerate}
    \item \textbf{Compter le texte source ($N_{\text{source}}$).} Compter \emph{tous} les mots, mots-outils inclus (\emph{le, de, et, mais, dans...}). Convention : un mot est une suite de caractères séparée par un espace. Les nombres (2025) comptent pour 1 mot ; les mots composés à traits d'union (moto-taxi, porte-monnaie) comptent pour 1 mot.
    \item \textbf{Calculer la cible ($N_{\text{cible}}$).} Si la consigne dit « au quart » : $N_{\text{cible}} = N_{\text{source}} / 4$. Si la consigne dit « 125 mots » : $N_{\text{cible}} = 125$.
    \item \textbf{Calculer la fourchette de tolérance.} Marge $= N_{\text{cible}} \times 0{,}10$. Fourchette $= [N_{\text{cible}} - \text{marge}\,;\,N_{\text{cible}} + \text{marge}]$.
    \item \textbf{Rédiger un premier jet} sans compter mot à mot, en visant visuellement la cible.
    \item \textbf{Compter le jet ($N_{\text{jet}}$) et ajuster :}
    \begin{itemize}
        \item si $N_{\text{jet}}$ dépasse le maximum : resserrer (supprimer des adjectifs, condenser, généraliser davantage) ;
        \item si $N_{\text{jet}}$ est sous le minimum : développer légèrement (rétablir un connecteur, préciser un agent, ajouter une précision signifiante).
    \end{itemize}
    \item \textbf{Recompter la version finale} et noter le nombre de mots en marge de la copie.
\end{enumerate}
\end{methode}

\begin{savaistu}[La règle du quart au Probatoire et au Bac technique]
La tradition scolaire fixe la longueur du résumé à \textbf{un quart (25\,\%) du texte source}, avec une \textbf{marge de tolérance de 10\,\%}.
\par Exemple : texte de 500 mots $\to$ cible de 125 mots. Marge : $125 \times 0{,}10 = 12{,}5$ $\to$ fourchette acceptée : \textbf{113 à 138 mots} (arrondis).
\par \textbf{Attention :} certains sujets imposent un nombre fixe (par exemple « en 120 mots »). La tolérance de $\pm 10\,\%$ s'applique alors à ce nombre imposé (108 à 132 mots). Lisez \emph{attentivement} l'énoncé.
\end{savaistu}

\begin{popfigure}
\begin{center}
\begin{tikzpicture}
\begin{axis}[
    ybar,
    bar width=2.2cm,
    width=0.75\linewidth,
    height=7.5cm,
    ymin=0, ymax=580,
    ylabel={Nombre de mots},
    symbolic x coords={Texte source, R\'esum\'e (cible)},
    xtick=data,
    ytick={0,125,250,375,500},
    nodes near coords,
    every node near coord/.append style={font=\bfseries, color=popdark},
    title style={font=\bfseries\small, color=popdark},
    title={Rapport quantitatif : 500 mots $\to$ 125 mots (25\,\%)},
    enlarge x limits=0.5,
]
\addplot[fill=popblue, draw=popdark] coordinates {(Texte source,500)};
\addplot[fill=poporange, draw=popdark] coordinates {(R\'esum\'e (cible),125)};
\end{axis}
\end{tikzpicture}
\end{center}
\legende{Graphique en barres illustrant la règle du quart : un texte source de 500 mots (barre bleue) donne un résumé cible de 125 mots (barre orange), soit $25\,\%$ de la longueur initiale, avec une fourchette de tolérance de 113 à 138 mots.}
\end{popfigure}

\courssub{Les interdits du résumé : ce qui fait perdre des points}

\begin{aretenir}[Les quatre interdits majeurs (sanctionnés au Probatoire et au Bac)]
\begin{enumerate}
    \item \textbf{Le \cle{patchwork} (recopiage) :} enchaîner des phrases ou des fragments de phrases du texte source sans reformulation personnelle. Même tronquées, ce ne sont pas \emph{vos} phrases.
    \item \textbf{L'apport subjectif :} ajouter son opinion (« \emph{À mon avis}, l'auteur a raison »), un jugement de valeur (« \emph{Heureusement}, la loi a changé ») ou une connaissance extérieure absente du texte.
    \item \textbf{La trahison de la thèse :} inverser le sens, attribuer à l'auteur le contraire de sa pensée, ou omettre l'argument principal.
    \item \textbf{Le hors-longueur :} dépasser la borne haute ($N_{\text{cible}} + 10\,\%$) ou rester sous la borne basse ($N_{\text{cible}} - 10\,\%$). Un résumé trop court est incomplet ; trop long, il n'en est plus un.
\end{enumerate}
\end{aretenir}

\begin{attention}[Erreur fréquente : le « patchwork » masqué]
L'élève croit reformuler en changeant deux ou trois mots par phrase : « \emph{La déforestation massive entraîne des conséquences dramatiques} » $\to$ « \emph{La déforestation importante cause des conséquences graves} ». \textbf{C'est du patchwork.} Le correcteur repère la structure syntaxique identique (sujet + verbe + compléments). \textbf{Vraie reformulation :} changer la construction, par exemple « \emph{Menacée par la déforestation, la biodiversité subit un déclin grave} » (passage à une tournure participiale et à un nouveau sujet).
\end{attention}

\courssub{La rédaction finale : cohérence, style impersonnel, relecture}

Une fois les idées sélectionnées, réduites et reformulées au brouillon, reste l'écriture de la version définitive.

\begin{propriete}[Les marqueurs d'un résumé réussi]
\begin{itemize}
    \item \textbf{Style impersonnel :} l'auteur du résumé s'efface. On utilise : « \emph{L'auteur montre que...} », « \emph{Le texte démontre que...} », « \emph{Il ressort que...} ». \textbf{À proscrire :} « \emph{Je pense que...} », « \emph{D'après moi...} », « \emph{Dans ce texte, on voit que...} » (trop oral).
    \item \textbf{Connecteurs logiques :} ils reproduisent l'enchaînement argumentatif : \emph{d'abord, ensuite, en effet, par ailleurs, cependant, par conséquent, enfin}. Éviter la répétition de « et ».
    \item \textbf{Temps verbaux :} en général le \emph{présent} (de vérité générale) pour rendre compte des idées : « L'auteur \emph{affirme}, \emph{explique}, \emph{conclut} ». Le passé composé reste possible si le texte rapporte des faits datés.
    \item \textbf{Paragraphe unique :} sauf consigne contraire, le résumé forme \emph{un seul bloc} de texte, sans sauts de ligne artificiels.
\end{itemize}
\end{propriete}

\begin{methode}[Rédaction finale et relecture : la check-list en cinq points]
\begin{enumerate}
    \item \textbf{Fidélité :} ai-je conservé la thèse exacte et tous les arguments principaux, dans l'ordre, sans rien ajouter ?
    \item \textbf{Reformulation :} aucune phrase n'est-elle recopiée ? Le vocabulaire et la syntaxe sont-ils miens ? (Test : cacher le texte source et relire mon résumé : est-il fluide et autonome ?)
    \item \textbf{Cohérence :} les connecteurs assurent-ils la liaison ? Le texte se comprend-il sans le source ?
    \item \textbf{Longueur :} le compte final est-il dans la fourchette de $\pm 10\,\%$ ? Ai-je noté le nombre de mots en marge ?
    \item \textbf{Correction :} accords (sujet-verbe, groupe nominal), orthographe, ponctuation, majuscules. \textbf{Relire à l'envers} (de la dernière phrase vers la première) pour chasser les fautes d'inattention.
\end{enumerate}
\end{methode}

\courssub{Entraînement gradué : du paragraphe au texte long}

La progression pédagogique exige deux niveaux d'exercices, conformes aux épreuves du Probatoire et du Baccalauréat technique.

\begin{exercice}[Niveau 1 : paragraphe court (environ 100 mots $\to$ 30 mots $\pm 10\,\%$, soit 27 à 33 mots)]
\textbf{Texte :} « \textit{L'urbanisation galopante des métropoles africaines, comme Yaoundé ou Douala, engendre une pression foncière sans précédent. Les quartiers périphériques, jadis zones agricoles, sont grignotés par l'habitat spontané. Cette extension anarchique prive les citadins d'espaces verts et complique l'assainissement. Les autorités municipales, dépassées par l'ampleur du phénomène, peinent à doter ces zones en voiries et en réseaux d'eau potable. Pourtant, des plans d'urbanisme directeurs existent, mais leur application bute sur le manque de moyens financiers et la complexité des procédures d'expropriation. In fine, c'est la qualité de vie des populations les plus vulnérables qui se dégrade inexorablement.} » (102 mots)

\textbf{Consigne :} résumez ce paragraphe en 30 mots environ (fourchette : 27 à 33 mots), en appliquant la méthode en cinq étapes.
\end{exercice}

\begin{corrige}[Exercice niveau 1]
\textbf{Structure argumentative :} thèse : l'urbanisation rapide des métropoles africaines dégrade les conditions de vie. Argument 1 : l'habitat spontané gagne les périphéries (perte d'espaces verts, assainissement difficile). Argument 2 : les autorités sont dépassées ; les plans d'urbanisme ne sont pas appliqués (freins financiers et juridiques). Conclusion : dégradation de la qualité de vie des plus vulnérables.
\par \textbf{Réduction et reformulation :}
\begin{itemize}
    \item Suppression : « galopante », « jadis zones agricoles », « sans précédent », « anarchique », « inexorablement ».
    \item Condensation : « voiries et réseaux d'eau potable » $\to$ « infrastructures de base » ; « manque de moyens financiers et complexité des procédures d'expropriation » $\to$ « freins financiers et juridiques ».
    \item Généralisation : « Yaoundé ou Douala » $\to$ (déjà couvert par « métropoles africaines ») ; « espaces verts et assainissement » $\to$ intégré dans « infrastructures de base ».
\end{itemize}
\textbf{Résumé proposé (31 mots) :}
\par \textit{« L'urbanisation des métropoles africaines génère une pression foncière transformant les périphéries en habitat spontané, dépourvu d'infrastructures de base. Malgré des plans directeurs, des freins financiers et juridiques entravent leur application, dégradant la vie des populations vulnérables. »}
\par \textbf{Compte :} 31 mots, dans la fourchette $[27\,;\,33]$. \textbf{Validé.} Vérification : thèse conservée, deux arguments présents, connecteurs (\emph{malgré}), aucune phrase recopiée.
\end{corrige}

\begin{exercice}[Niveau 2 : texte long (500 mots $\to$ 125 mots $\pm 10\,\%$, soit 113 à 138 mots)]
Exercice similaire à l'épreuve type du Probatoire ou du Baccalauréat technique. Utiliser un extrait argumentatif du \emph{Vieux Nègre et la Médaille} de Ferdinand Oyono (par exemple le discours du commandant ou la réflexion de Meka sur sa médaille) ou un éditorial de presse.
\par \textbf{Consigne :} appliquez la méthode complète : lecture analytique $\to$ structure argumentative $\to$ brouillon (réduction et reformulation) $\to$ premier jet $\to$ comptage $\to$ ajustement $\to$ copie propre $\to$ relecture. Durée conseillée : 45 à 50 minutes.
\end{exercice}

\courssub{Amélioration du résumé : grille d'auto-évaluation}

Pour progresser, l'élève doit s'approprier les critères du correcteur. Voici une grille à utiliser \emph{systématiquement} sur chaque brouillon.

\begin{center}
\begin{tabularx}{\linewidth}{|l|X|X|c|}\hline
\rowcolor{popblueL} \textbf{Critère} & \textbf{Indicateurs de réussite} & \textbf{Action corrective si insuffisant} & \textbf{Points} \\ \hline
\textbf{1. Fidélité (thèse et arguments)} & Thèse exacte ? Arguments principaux tous présents, dans l'ordre ? Aucun ajout extérieur ? & Relire le source ; cocher chaque idée-force ; réintégrer l'oubli ; supprimer l'intrus. & /5 \\ \hline
\textbf{2. Reformulation personnelle} & Aucune phrase recopiée ? Vocabulaire varié ? Syntaxe transformée ? Style impersonnel ? & Réécrire les phrases « patchwork » ; nominaliser ; changer de voix ou de sujet. & /5 \\ \hline
\textbf{3. Cohérence (liaison)} & Connecteurs logiques adaptés ? Enchaînement fluide ? Texte compréhensible sans le source ? & Insérer \emph{donc, cependant, par ailleurs, en outre} ; vérifier les reprises anaphoriques (\emph{il, ce, cela}). & /4 \\ \hline
\textbf{4. Longueur (contrainte)} & Nombre de mots dans la fourchette de $\pm 10\,\%$ ? Noté en marge ? & Trop long : condenser, généraliser, supprimer l'accessoire. Trop court : rétablir un connecteur ou une précision. & /3 \\ \hline
\textbf{5. Correction (langue)} & Accords, orthographe, ponctuation, majuscules ? & Relire à l'envers ; vérifier chaque accord ; consulter le dictionnaire en cas de doute. & /3 \\ \hline
\end{tabularx}
\end{center}

\begin{experience}[Atelier de révision par les pairs]
\textbf{Matériel :} trois résumés anonymes du même texte (niveau 2), la grille d'évaluation ci-dessus, des surligneurs (vert = fidélité, orange = reformulation, rouge = erreur ou hors-longueur).
\par \textbf{Protocole :}
\begin{enumerate}
    \item Constituer des groupes de trois. Chaque élève reçoit deux résumés produits par d'autres groupes.
    \item Lecture silencieuse et remplissage de la grille (10 minutes).
    \item Discussion en trio : « Pourquoi as-tu jugé la fidélité insuffisante ? Quelle phrase est du patchwork ? » (10 minutes).
    \item Retour aux auteurs : lecture des grilles reçues, puis réécriture ciblée (10 minutes).
\end{enumerate}
\textbf{Observation :} les élèves détectent plus facilement le patchwork et les trahisons de sens sur la copie d'autrui que sur la leur.
\par \textbf{Interprétation :} le regard d'autrui développe l'esprit critique nécessaire à l'auto-évaluation. Cet atelier est à renouveler avant chaque devoir type Bac.
\end{experience}

\begin{aretenir}[Synthèse : la contraction, un exercice de rigueur intellectuelle]
La contraction de texte n'est pas un simple « raccourcissement » mais une \textbf{restructuration intelligente}. Elle exige :
\begin{enumerate}
    \item une \textbf{lecture analytique} (structure argumentative) -- vue à la section précédente ;
    \item une \textbf{maîtrise technique} : suppression, condensation, généralisation ;
    \item une \textbf{compétence linguistique} : reformulation lexicale et syntaxique, style impersonnel ;
    \item une \textbf{discipline quantitative} : comptage des mots, fourchette de $\pm 10\,\%$ ;
    \item une \textbf{éthique de la fidélité} : ni opinion personnelle, ni trahison du sens, ni patchwork.
\end{enumerate}
C'est cet ensemble qui est évalué au Probatoire et au Baccalauréat technique. L'entraînement régulier (un résumé par semaine) sur des textes variés -- essais, articles de presse, extraits du \emph{Vieux Nègre et la Médaille} -- est la voie la plus sûre vers la réussite.
\end{aretenir}

\coursec{Application au Vieux nègre et la médaille : lectures méthodiques 5 et 6, exposés}

Après avoir maîtrisé les \cle{techniques de réduction} (suppression, généralisation, substitution) et la \cle{reformulation} (changement de point de vue, nominalisation, passage à l'impersonnel), nous mettons désormais ces outils au service de l'analyse d'une œuvre au programme : \emph{Le Vieux nègre et la médaille} de Ferdinand Oyono. Cette application concrète permet de vérifier que la contraction n'est pas un exercice purement mécanique, mais qu'elle suppose une compréhension fine des \cle{procédés d'écriture} (ironie, énonciation performative, variétés linguistiques) et de la \cle{structure argumentative} implicite du récit. C'est en analysant \emph{comment} le texte dit ce qu'il dit que l'on peut en extraire l'essentiel pour le reformuler en 80 mots.

\courssub{Lecture méthodique n°5 : La scène de la remise de la médaille — Énoncés performatifs et ironie du discours officiel}

\begin{definition}[La lecture méthodique]
La \cle{lecture méthodique} est une analyse guidée d'un extrait qui articule trois niveaux : la \textbf{compréhension littérale} (que se passe-t-il ?), l'étude des \textbf{procédés d'écriture} (comment est-ce dit ? : syntaxe, lexique, figures, énonciation) et l'\textbf{interprétation} (que veut dire le texte ? quelle vision du monde propose-t-il ?). Elle prépare l'exposé oral et l'écrit d'examen.
\end{definition}

L'extrait étudié (début du chapitre 3, environ 400 mots) décrit la cérémonie officielle où l'administrateur colonial remet la médaille à Meka. Le texte oppose deux discours : celui de l'administrateur (solennel, vide) et celui du narrateur (ironique, lucide).

\begin{methode}[Repérer les énoncés performatifs dans un texte narratif]
\begin{enumerate}
\item \textbf{Isoler les séquences dialoguées} ou rapportées où un personnage investi d'une autorité institutionnelle parle (ici : l'administrateur).
\item \textbf{Identifier le verbe performatif} (verbe d'action : « je décerne », « je nomme », « je déclare ») à la 1\textsuperscript{re} personne du présent de l'indicatif.
\item \textbf{Vérifier les conditions de félicité} (Austin) : le locuteur a-t-il le pouvoir ? Le contexte est-il rituel ? La formule est-elle conventionnelle ?
\item \textbf{Analyser l'effet produit} : l'énoncé \emph{fait} ce qu'il \emph{dit} (Meka devient médaillé), mais le narrateur en montre la vacuité.
\end{enumerate}
\end{methode}

Dans l'extrait, l'administrateur prononce : \emph{« Au nom du Président de la République, je vous décerne la médaille de la Légion d'honneur. »} C'est un \cle{énoncé performatif} pur : il transforme socialement Meka en « chevalier ». Pourtant, le narrateur commente : \emph{« Meka ne comprenait pas un mot de ce charabia. »} L'\cle{ironie} naît du décalage entre la solennité du rituel (le français administratif, \cle{variété standard} normative) et l'incompréhension du récipiendaire (qui parle le \cle{français populaire camerounais} ou le \cle{bassa}). Le discours officiel est \emph{constatif} en apparence (il décrit un mérite), mais \emph{performatif} en réalité : il \emph{crée} une allégeance.

\begin{attention}[Erreur fréquente : raconter au lieu d'analyser]
Ne pas écrire : « L'administrateur donne la médaille à Meka qui est content. »
Écrire : « L'énoncé performatif "je vous décerne" instaure une relation de domination symbolique ; l'ironie narrative (focalisation interne sur l'incompréhension de Meka) déconstruit la portée honorifique du rituel. »
\end{attention}

\begin{popfigure}
\begin{tikzpicture}[
    timeline/.style={very thick, popblue},
    event/.style={circle, fill=popblue, inner sep=2pt},
    labelbox/.style={rectangle, rounded corners, fill=popblueL, draw=popblue, inner sep=4pt, font=\small, align=center, text width=3.5cm},
    arrow/.style={-Stealth, thick, poporange}
]
% Ligne de temps
\draw[timeline] (0,0) -- (15,0);
% Événements
\node[event] at (1,0) {};
\node[labelbox, anchor=south, align=center] at (1,0.8){\textbf{Attente}\\Meka attend la médaille,\\ espoir naïf};
\draw[arrow] (1,0.05) -- ++(0,0.5);
\node[event] at (4,0) {};
\node[labelbox, anchor=south, align=center] at (4,0.8){\textbf{Médaille}\\Cérémonie officielle,\\ performatif \& ironie};
\draw[arrow] (4,0.05) -- ++(0,0.5);
\node[event] at (7,0) {};
\node[labelbox, anchor=south, align=center] at (7,0.8){\textbf{Fête}\\Banquet au village,\\ fierté apparente};
\draw[arrow] (7,0.05) -- ++(0,0.5);
\node[event] at (10,0) {};
\node[labelbox, anchor=south, align=center] at (10,0.8){\textbf{Arrestation}\\Arrestation brutale,\\ humiliation publique};
\draw[arrow] (10,0.05) -- ++(0,0.5);
\node[event] at (13,0) {};
\node[labelbox, anchor=south, align=center] at (13,0.8){\textbf{Désillusion}\\Prison, médaille\\ confisquée, lucidité};
\draw[arrow] (13,0.05) -- ++(0,0.5);
% Flèche nord / temps
\node[font=\small, popmuted] at (15.5,0) {$\rightarrow$ Temps narratif};
\end{tikzpicture}
\legende{Frise chronologique des étapes du récit de Meka : de l'attente naïve à la désillusion finale, en passant par l'apogée ironique de la remise de médaille.}
\end{popfigure}

\courssub{Lecture méthodique n°6 : La chute de Meka — Dégagement de la structure argumentative implicite}

L'extrait de la fin (chapitre 6) relate l'arrestation de Meka, accusé de vol par le commandant qui lui avait remis la médaille. La structure argumentative du récit n'est pas explicite (pas de thèse/antithèse formelle), elle est \emph{implicite}, construite par la \cle{structure narrative} elle-même.

\begin{propriete}[Structure argumentative implicite du récit oyonesque]
Le récit suit un schéma \textbf{ironique en trois temps} :
\begin{enumerate}
\item \textbf{Thèse apparente (l'illusion) :} La médaille = reconnaissance, intégration, honneur. (Discours colonial).
\item \textbf{Antithèse réelle (la preuve par les faits) :} L'arrestation = trahison, violence, spoliation. (Réalité coloniale).
\item \textbf{Synthèse amère (la lucidité) :} La médaille = symbole de l'aliénation, « chaîne en or ». (Vision de l'auteur).
\end{enumerate}
Cette structure mime le parcours de la conscience de Meka : de l'aliénation à la prise de conscience.
\end{propriete}

Le lexique opère un \cle{basculement sémantique} : « honneur », « mérite », « France » (champ lexical de la grandeur) cèdent la place à « voleur », « prison », « menottes », « sale nègre » (champ lexical de l'avilissement). L'ironie devient \cle{paradoxe} : la récompense est l'instrument de la chute. Meka \emph{porte} la médaille au moment où on lui met les menottes ; l'objet même de sa fierté devient la preuve de sa culpabilité aux yeux du colon.

\begin{exemplebox}[Analyse d'un paradoxe textuel]
\textbf{Texte :} « Il portait la médaille au cou comme une croix. On lui mit les menottes aux poignets. »
\textbf{Analyse :} Parallélisme syntaxique (portait / mit) et opposition sémantique (médaille / menottes, cou / poignets). La médaille, symbole chrétien de la rédemption (« comme une croix »), devient le pendant ironique des menottes, symbole de l'oppression. Le récit argumente sans le dire : l'ordre colonial « récompense » pour mieux enchaîner.
\end{exemplebox}

\begin{exoresolu}[Résumé d'un extrait de 400 mots en 100 mots (cible : 80 mots)]
\textbf{Extrait :} Scène de l'arrestation (fin du roman). Meka, médaillé, est accusé de vol par le commandant. Il est menotté, humilié devant les siens, la médaille arrachée.
\textbf{Consigne :} Résumer en conservant la désillusion de Meka et l'ironie du narrateur (4-5 phrases).
\tcblower
\textbf{Proposition de résumé (82 mots) :}
« Après la cérémonie fastueuse où l'administrateur lui décerne la Légion d'honneur, Meka est brutalement arrêté par le même commandant pour un vol qu'il n'a pas commis. L'ironie narrative atteint son paroxysme : l'homme qui portait fièrement la médaille « comme une croix » se voit menotter sous les yeux de son village. La récompense coloniale, symbole d'intégration, se révèle être le piège qui livre le « vieux nègre » à l'humiliation suprême. Oyono démontre que l'honneur colonial n'est qu'une chaîne dorée. »
\textbf{Justification :} Suppression des détails anecdotiques (noms des témoins, description du temps), généralisation (« cérémonie fastueuse »), reformulation au discours indirect libre (« Meka est brutalement arrêté »), conservation du paradoxe central (médaille/croix vs menottes) et de l'intention auctoriale (dernière phrase).
\end{exoresolu}

\begin{exemplebox}[Exemple chiffré : Contraction ciblée]
\begin{center}
\begin{tabularx}{\linewidth}{|l|X|X|}\hline
\textbf{Étape} & \textbf{Texte original (extrait)} & \textbf{Texte contracté (résumé)} \\ \hline
Suppression & Description du temps qu'il fait, noms des villageois présents. & — \\ \hline
Généralisation & « Le commandant, un grand maigre au visage rougeaud... » & « Le commandant » \\ \hline
Substitution & « Il mit les menottes à ses poignets, il le poussa, il l'insulta. » & « Il l'arrête et l'humilie. » \\ \hline
Nominalisation & « Meka ne comprenait pas pourquoi on l'arrêtait. » & « L'incompréhension de Meka » \\ \hline
Changement de point de vue & Discours direct : « Voleur ! » & Discours rapporté : « L'accusation de vol » \\ \hline
\end{tabularx}
\end{center}
\legende{Application des techniques de réduction sur l'extrait de l'arrestation.}
\end{exemplebox}

\courssub{Analyse du lexique colonial et des variétés de français dans les extraits}

Le roman met en scène une \cle{contact de langues} et de variétés. Trois niveaux cohabitent :
\begin{itemize}
\item \textbf{Le français standard administratif} (administrateur, commandant) : syntaxe complexe, vocabulaire abstrait (\"décerner\", \"insigne\", \"république\"), énoncés performatifs. Variété de \cle{pouvoir} et de \cle{distance}.
\item \textbf{Le français populaire camerounais / « petit nègre »} (Meka, villageois) : syntaxe simplifiée, emprunts au bassa (\"ndolo\", \"mami\"), proverbes traduits. Variété de \cle{proximité} et de \cle{résistance} identitaire.
\item \textbf{Le français littéraire du narrateur} (Oyono) : ironie, métaphores, intertextualité biblique (\"comme une croix\"). Variété \cle{esthétique} et \cle{critique}.
\end{itemize}

\begin{methode}[Analyser les variétés lexicales pour le résumé]
\begin{enumerate}
\item Repérer les \textbf{marqueurs lexicaux} propres à chaque variété (ex: \"insigne\" vs \"médaille\" vs \"breloque\").
\item Noter les \textbf{gloses} (traductions intégrées) : \"ndolo (amour)\".
\item Dans le résumé, \textbf{neutraliser la variété orale} (passer au standard) mais \textbf{conserver les termes clés} culturels s'ils sont porteurs de sens (ex: \"médaille\" reste \"médaille\", pas \"distinction\").
\end{enumerate}
\end{methode}

\courssub{Exposé n°1 : « L'ironie comme arme de dénonciation chez Oyono » (Travail de groupe)}

\begin{methode}[Pour un exposé réussi : structure et oralité]
\begin{enumerate}
\item \textbf{Accroche (30 s) :} Citation choc + problématique. Ex: « "Il portait la médaille comme une croix." Comment l'ironie transforme-t-elle l'objet en arme ? »
\item \textbf{Annonce du plan (2 parties ou 3) :} 
    \begin{itemize}
    \item I. L'ironie situationnelle : le décalage cérémonie/réalité.
    \item II. L'ironie verbale : le discours officiel vidé de son sens.
    \item III. L'ironie dramatique : le lecteur sait (Meka volé) avant le personnage.
    \end{itemize}
\item \textbf{Développement (6 min) :} Une idée = un paragraphe = une citation analysée (procédé + effet).
\item \textbf{Conclusion (30 s) :} Réponse à la problématique + ouverture (autres œuvres : \emph{Une vie de boy}, \emph{Le Pauvre Christ de Bomba}).
\item \textbf{Regard public / voix portée / temps de parole équitable.}
\end{enumerate}
\end{methode}

\courssub{Exposé n°2 : « La médaille : symbole de l'illusion coloniale » (Travail de groupe, appui sur citations)}

Cet exposé exige un travail sur le \cle{symbole}. La médaille n'est pas un simple accessoire ; elle structure le roman (titre, début, fin).

\begin{popfigure}
\begin{tikzpicture}[
    nodeStyle/.style={rectangle, rounded corners, draw=popdark, thick, inner sep=8pt, font=\small, align=center, text width=3.8cm},
    arrowStyle/.style={-Stealth, thick, poporange},
    doubleArrow/.style={<->, thick, poppurple, dashed},
    centerNode/.style={circle, draw=popblue, thick, fill=popblueL, inner sep=10pt, font=\bfseries\small, text width=2.5cm, align=center}
]
% Nœud central
\node[centerNode] (medaille) at (0,0) {LA MÉDAILLE};
% Sens apparent (gauche)
\node[nodeStyle, fill=popgreenL, draw=popgreen, align=center] (app) at (-5, 2.5){
    \textbf{Sens apparent}\\(Discours colonial)\\
    \begin{itemize}\item Récompense du mérite
    \item Intégration à la « Grande France »
    \item Honneur, fierté, statut social
    \item Objet sacré (\"comme une croix\")
    \end{itemize}
};
% Sens réel (droite)
\node[nodeStyle, fill=poppinkL, draw=poppink, align=center] (reel) at (5, 2.5){
    \textbf{Sens réel}\\(Réalité vécue / Ironie)\\
    \begin{itemize}\item Instrument de domination
    \item Aliénation / « Chaîne en or »
    \item Piège : attire la jalousie, la répression
    \item Objet confisqué (spoliation finale)
    \end{itemize}
};
% Basculement (bas)
\node[nodeStyle, fill=poporangeL, draw=poporange, align=center] (basculement) at (0, -3){
    \textbf{Basculement narratif}\\
    L'ironie opère le passage de l'apparent au réel.\\
    Meka \emph{porte} le symbole $\rightarrow$ Meka \emph{subit} le symbole.\\
    \emph{« La médaille lui brûlait le cou. »}
};
% Flèches
\draw[arrowStyle] (medaille) -- node[above, sloped, font=\tiny] {représente} (app);
\draw[arrowStyle] (medaille) -- node[above, sloped, font=\tiny] {masque} (reel);
\draw[doubleArrow] (app) -- node[above, font=\tiny] {Ironie / Paradoxe} (reel);
\draw[arrowStyle] (app) -- (basculement);
\draw[arrowStyle] (reel) -- (basculement);
\end{tikzpicture}
\legende{Schéma de la médaille comme symbole : opposition structurale entre le sens apparent (discours officiel) et le sens réel (dénonciation romanesque), médiatisé par l'ironie narrative.}
\end{popfigure}

\begin{savaistu}[Prix Fénéon et programme camerounais]
Publié en \textbf{1956}, \emph{Le Vieux nègre et la médaille} a obtenu le \textbf{prix Fénéon} (révélant de jeunes auteurs de langue française). C'est le premier roman de Ferdinand Oyono (1929–2010), diplomate et écrivain camerounais. L'œuvre figure au programme officiel du Cameroun depuis les années 1960 car elle incarne la \cle{littérature de la décolonisation} : elle utilise la langue du colonisateur (le français) pour en démonter les mécanismes (ironie, énonciation, lexique). Étudier ce texte en Terminale technique, c'est analyser comment la littérature africaine d'expression française a détourné l'outil linguistique pour dire l'histoire des vaincus.
\end{savaistu}

\courssub{Synthèse : Rédaction en classe d'un résumé de 80 mots, compte-rendu collectif et remédiation ciblée}

\begin{exercice}[Épreuve synthèse : Résumé contracté]
\textbf{Texte support :} Extrait de la scène de l'arrestation (fin du chapitre 6, env. 380 mots).
\textbf{Consigne :} Rédiger un résumé de \textbf{80 mots (± 10\%)} qui rende compte de la chute de Meka, de l'ironie du sort (médaille/menottes) et de la dénonciation implicite du système colonial.
\textbf{Contraintes :} 4 à 5 phrases. Vocabulaire standard. Pas de citation. Respect de la chronologie logique (cérémonie $\rightarrow$ arrestation $\rightarrow$ sens).
\end{exercice}

\begin{corrige}[Exercice n° Synthèse — Proposition de correction commentée]
\textbf{Résumé proposé (79 mots) :}
« Lors de la cérémonie officielle, l'administrateur décerne à Meka la Légion d'honneur, le propulsant au rang de notable. Mais cette gloire est de courte durée : le même commandant l'accuse de vol et le fait arrêter sous les yeux de son village. L'ironie cruelle veut que Meka, qui portait la médaille « comme une croix », se retrouve menotté, l'insigne confisqué. Par ce renversement, Oyono dévoile la médaille comme un leurre : l'honneur colonial n'est qu'une chaîne dorée rivée au cou des dominés. »

\textbf{Analyse de la correction (Remédiation) :}
\begin{itemize}
\item \textbf{Structure :} 4 phrases. 1. Contexte/Cérémonie. 2. Rupture/Arrestation. 3. Paradoxe visuel (médaille/croix vs menottes). 4. Interprétation/Thèse.
\item \textbf{Techniques utilisées :} Généralisation (\"notable\"), Suppression (détails du vol, dialogues), Substitution (\"propulsant\" pour \"remise solennelle\"), Nominalisation (\"ce renversement\"), Reformulation énonciative (passage du \"je\" narratif au \"Oyono dévoile\").
\item \textbf{Points de vigilance (Remédiation ciblée) :}
    \begin{itemize}
    \item \textbf{Attention :} Ne pas écrire \"Meka est volé\" (ambigu) mais \"Meka est accusé de vol\".
    \item \textbf{Attention :} Ne pas oublier l'ironie (\"gloire de courte durée\", \"leurre\"). Un résumé sans ironie rate le sens du texte.
    \item \textbf{Comptage :} Vérifier le nombre de mots (ici 79). Si 95 mots $\rightarrow$ resserrer les périphrases.
    \end{itemize}
\end{itemize}
\end{corrige}

\begin{aretenir}[Bilan de la séquence Contraction / Le Vieux nègre et la médaille]
\begin{itemize}
\item La \cle{contraction} exige une \cle{lecture méthodique} préalable : comprendre l'ironie, repérer les performatifs, analyser le lexique.
\item La \cle{structure argumentative} d'un récit peut être implicite (schéma ironique : illusion $\rightarrow$ rupture $\rightarrow$ lucidité).
\item Le \cle{résumé} (80 mots) est une \emph{reconstruction} : il sélectionne, hiérarchise, reformule au style indirect, neutralise l'oralité mais garde la charge critique (le paradoxe final).
\item L'\cle{exposé} oral structure la pensée : problématique $\rightarrow$ plan annoncé $\rightarrow$ citations analysées $\rightarrow$ conclusion argumentée.
\end{itemize}
\end{aretenir}

\newpage

\coursec{Activité d'intégration}

\courssub{Objectifs de l'activité}
Mobiliser l'ensemble des techniques de contraction de texte étudiées dans la leçon (analyse de la structure argumentative, élimination des redondances, reformulation en français standard, respect de la contrainte de longueur) pour produire un résumé administratif professionnel, fidèle et cohérent, destiné à une autorité hiérarchique (le sous-préfet).

\begin{aretenir}[Rappel : ce qu'on attend d'un résumé réussi]
Un résumé est un texte \cle{bref}, \cle{fidèle} (aucune idée ajoutée, aucun contresens), \cle{reformulé} (phrases et vocabulaire personnels) et \cle{cohérent} (enchaînement logique assuré par des connecteurs). Il respecte la contrainte de longueur imposée.
\end{aretenir}

\courssub{Situation}
Vous êtes chefs de quartier adjoints dans l'arrondissement de Yaoundé II. Suite aux fortes pluies de la semaine dernière, le sous-préfet vous demande un rapport de synthèse sur la situation du marché de Mokolo pour une réunion de crise prévue dans deux heures. Il ne dispose que de peu de temps pour lire votre note ; elle ne doit pas dépasser \textbf{150 mots} (tolérance $\pm 10\%$, soit entre 135 et 165 mots).

Vous disposez de l'article de presse ci-dessous, publié ce matin dans \emph{La Voix du Koa}, comme unique source d'information brute. Cet article est un compte rendu journalistique d'environ 600 mots.

\begin{center}
\textbf{\textsc{Le marché de Mokolo face aux pluies : le cri de détresse des commerçants}}
\end{center}

\emph{Par notre correspondant, Yaoundé, le 15 octobre 2023.}

Depuis le début de la grande saison des pluies, le marché de Mokolo, poumon économique de la capitale, vit au ralenti. La thèse est sans appel : \textbf{les inondations récurrentes paralysent l'activité commerciale et engendrent des pertes financières colossales pour les opérateurs économiques}. Si rien n'est fait d'urgence, c'est tout un pan de l'économie informelle qui risque de s'effondrer.

Le premier argument, c'est la dégradation avancée des infrastructures. \og{}L'eau ne coule plus, elle stagne\fg{}, soupire Madame Ngo Biboum, revendeuse de légumes au \og{}Carrefour Mvog-Ada\fg{} depuis vingt ans. Selon le délégué du marché, près de \textbf{300 étals} sur les 1\,200 que compte le site principal sont actuellement impraticables. Les deux artères principales, l'allée \og{}A\fg{} et l'allée \og{}B\fg{}, sont coupées par des flaques d'eau boueuse atteignant parfois 40 centimètres de hauteur. Les caniveaux, bouchés par des tonnes de déchets plastiques et de sable, ne jouent plus leur rôle. \og{}C'est une honte pour la ville\fg{}, lance, ironique, un jeune pousse-pousse : \og{}On a des piscines gratuites, mais pas de clients\fg{}. Ce paradoxe --- des piscines sans nageurs --- illustre l'absurdité de la situation.

Deuxièmement, l'impact économique est immédiat et brutal. Les denrées périssables (tomates, poisson, feuilles) pourrissent sur place. \og{}J'ai jeté pour 25\,000 FCFA de tomates rien que ce matin\fg{}, témoigne Monsieur Ousmanou, grossiste. La fédération des commerçants estime la baisse du chiffre d'affaires hebdomadaire à \textbf{15\,\%} en moyenne, soit une perte sèche de \textbf{50 millions de FCFA par semaine} pour l'ensemble du marché. Les clients, rebutés par la boue et l'odeur nauséabonde, désertent les lieux pour aller vers les supermarchés ou les petits marchés de quartier plus accessibles. \og{}Le client a horreur de la boue\fg{}, constate amèrement une vendeuse de pagnes. C'est un énoncé constatif qui sonne comme une sentence.

Troisièmement, les risques sanitaires et sécuritaires montent en flèche. L'eau stagnante devient un nid à moustiques (paludisme) et à bactéries (choléra, fièvre typhoïde). Les tas d'ordures non collectés s'accumulent aux abords des étals. Pire, des câbles électriques dénudés trempent dans l'eau par endroits, créant un risque d'électrocution réel. \og{}On travaille la peur au ventre\fg{}, confie un couturier. Le délégué du marché a adressé un \textbf{énoncé performatif} fort à la mairie : \og{}Je déclare l'état d'urgence sanitaire dans cette zone\fg{}. Cet énoncé, par sa force illocutoire, engage la responsabilité des autorités.

Face à ce tableau noir, les solutions existent mais tardent. Le curage des caniveaux, la réhabilitation du réseau d'évacuation des eaux pluviales et l'installation de pompes de relevage sont les mesures techniques préconisées par les experts. À court terme, un relogement temporaire des commerçants des zones les plus touchées vers l'esplanade de la mairie d'arrondissement est à l'étude. \og{}Nous attendons les actes, pas les promesses\fg{}, conclut Madame Ngo Biboum. Le temps presse : la prochaine averse est annoncée pour ce soir.

\begin{attention}[Repères de langue dans le texte]
\begin{itemize}
    \item \textbf{Énoncé constatif} : il décrit un fait, il peut être vrai ou faux. Exemple : \og{}Le client a horreur de la boue\fg{}.
    \item \textbf{Énoncé performatif} : il accomplit un acte par le seul fait d'être prononcé. Exemple : \og{}Je déclare l'état d'urgence sanitaire\fg{}.
    \item \textbf{Ironie} : dire le contraire de ce qu'on pense pour se moquer. Exemple : \og{}On a des piscines gratuites\fg{}.
    \item \textbf{Paradoxe} : réunion de deux idées contradictoires. Exemple : \og{}des piscines sans nageurs\fg{}.
\end{itemize}
Ces éléments relèvent du commentaire journalistique : ils seront supprimés ou neutralisés dans le résumé administratif.
\end{attention}

\courssub{Consignes}
Travail en groupes de 4 élèves. Chaque groupe rend une seule production collective.

\begin{enumerate}
    \item \textbf{Analyse structurelle} : Représentez la structure argumentative de l'article sous forme d'un \textbf{schéma arborescent} (thèse, arguments, exemples/illustrations, conclusion/appel). Identifiez explicitement les énoncés constatifs et performatifs repérés, ainsi que l'ironie et le paradoxe.
    \item \textbf{Sélection / Réduction} : Sur une photocopie de l'article, \textbf{barrez au crayon} tous les éléments supprimables (exemples détaillés, citations longues, commentaires journalistiques, répétitions, adjectifs qualificatifs non essentiels) en justifiant brièvement chaque suppression en marge (ex. : \og{}exemple illustratif\fg{}, \og{}répétition\fg{}, \og{}subjectivité\fg{}).
    \item \textbf{Rédaction du résumé} : Rédigez le \textbf{rapport pour le sous-préfet} (150 mots $\pm 10\%$).
        \begin{itemize}
            \item Langue : français standard administratif (pas de registre familier, pas de citations directes).
            \item Contenu : thèse reformulée, 3 arguments synthétisés avec chiffres clés, conclusion opérationnelle.
            \item Structure : un paragraphe unique ou deux paragraphes au maximum, connecteurs logiques obligatoires.
        \end{itemize}
    \item \textbf{Évaluation par les pairs} : Échangez votre production avec un autre groupe. Évaluez-la sur 20 points à l'aide de la grille ci-dessous. Inscrivez la note et trois commentaires argumentés. Remettez la copie évaluée au professeur.
\end{enumerate}

\begin{methode}[De l'article au rapport : la démarche en 5 étapes]
\begin{enumerate}
    \item \textbf{Lire} deux fois l'article et souligner la thèse et les mots d'articulation (\og{}premier argument\fg{}, \og{}deuxièmement\fg{}, \og{}troisièmement\fg{}, \og{}face à ce tableau noir\fg{}).
    \item \textbf{Dégager} la structure : 1 thèse, 3 arguments, 1 conclusion avec propositions.
    \item \textbf{Supprimer} : citations directes, témoignages nommés, figures de pensée (ironie, paradoxe), détails pittoresques, répétitions.
    \item \textbf{Conserver} : la thèse, les 3 arguments, les chiffres clés (300 étals sur 1\,200 ; 15\,\% ; 50 millions de FCFA par semaine), les solutions proposées.
    \item \textbf{Reformuler} en français standard, à l'infinitif ou à la troisième personne, avec des connecteurs (d'abord, ensuite, enfin, par conséquent), puis \textbf{compter les mots}.
\end{enumerate}
\end{methode}

\begin{center}
\textbf{GRILLE D'ÉVALUATION (20 pts)}
\end{center}
\begin{center}
\begin{tabularx}{\linewidth}{|l|X|c|}\hline
\rowcolor{popblueL}\textbf{Critère} & \textbf{Indicateurs} & \textbf{Points} \\ \hline
Fidélité au texte source & Idées principales respectées, pas de contresens, chiffres exacts conservés & 5 \\ \hline
Reformulation & Vocabulaire personnel, phrases nouvelles, pas de copier-coller, style administratif & 5 \\ \hline
Cohérence / Structure & Enchaînement logique, connecteurs, structure argumentative lisible & 4 \\ \hline
Respect de la contrainte & Longueur 135--165 mots, format rapport (objet, destinataire, date) & 3 \\ \hline
Qualité de la langue & Orthographe, grammaire, syntaxe, vocabulaire précis, impersonnalité & 3 \\ \hline
\end{tabularx}
\end{center}

\courssub{Ressources à mobiliser}
\begin{itemize}
    \item \textbf{Structure argumentative} : Thèse $\rightarrow$ Arguments (Argument 1 : infrastructures ; Argument 2 : économie ; Argument 3 : sanitaire et sécurité) $\rightarrow$ Illustrations (chiffres, citations) $\rightarrow$ Conclusion et propositions.
    \item \textbf{Techniques de réduction} : \cle{suppression} (exemples, détails, citations, répétitions, métadiscours journalistique) ; \cle{généralisation} (regrouper \og{}tomates, poisson, feuilles\fg{} $\rightarrow$ \og{}denrées périssables\fg{}) ; \cle{intégration} (fusionner plusieurs phrases en une seule).
    \item \textbf{Reformulation} : changement de point de vue (journalistique $\rightarrow$ administratif), nominalisation (\og{}l'eau stagne\fg{} $\rightarrow$ \og{}la stagnation des eaux\fg{} ; \og{}les clients désertent le marché\fg{} $\rightarrow$ \og{}la désertion du marché par la clientèle\fg{}), remplacement des citations directes par des propositions rapportées au discours indirect ou par de simples constats.
\end{itemize}

\begin{exemplebox}[Modèle de rapport attendu (150 mots)]
\textbf{À l'attention de Monsieur le Sous-préfet de Yaoundé II --- Yaoundé, le 15 octobre 2023.}

\textbf{Objet :} situation du marché de Mokolo après les fortes pluies.

Les inondations récurrentes paralysent l'activité du marché de Mokolo et provoquent d'importantes pertes financières. D'abord, les infrastructures sont très dégradées : environ 300 étals sur 1\,200 sont impraticables et les caniveaux bouchés ne permettent plus l'évacuation des eaux. Ensuite, l'impact économique est sévère : les denrées périssables sont perdues et le chiffre d'affaires hebdomadaire a baissé d'environ 15\,\%, soit une perte estimée à 50 millions de FCFA par semaine ; la clientèle se détourne vers d'autres points de vente. Enfin, les risques sanitaires et sécuritaires s'aggravent : eaux stagnantes, ordures accumulées et câbles électriques dénudés menacent la santé publique. Par conséquent, il est demandé le curage urgent des caniveaux, la réhabilitation du réseau d'évacuation et le relogement temporaire des commerçants les plus touchés.
\end{exemplebox}

\begin{attention}[Erreurs fréquentes à éviter]
\begin{itemize}
    \item Recopier des phrases entières de l'article (copier-coller sanctionné par le critère \og{}reformulation\fg{}).
    \item Conserver les citations directes et les noms des témoins (Mme Ngo Biboum, M. Ousmanou) : un rapport administratif est impersonnel.
    \item Ajouter des idées personnelles absentes du texte (contresens de fidélité).
    \item Oublier de compter les mots : hors de la fourchette 135--165, des points sont perdus.
    \item Employer un registre familier (\og{}c'est la galère\fg{}) au lieu du français standard.
\end{itemize}
\end{attention}

\begin{exercice}[Pour aller plus loin]
Reprenez le rapport de votre groupe et produisez une \textbf{version de 80 mots} destinée à être lue à la radio locale. Quels éléments supprimez-vous en priorité ? Justifiez vos choix en deux phrases.
\end{exercice}

\begin{corrige}[Exercice « Pour aller plus loin »]
On supprime d'abord les chiffres secondaires (1\,200 étals), puis les détails des solutions techniques (pompes de relevage), en ne gardant que : la thèse (inondations, marché paralysé), les trois arguments réduits à leur noyau (infrastructures dégradées, pertes de 50 millions de FCFA par semaine, risques sanitaires) et l'appel à l'action urgente. Justification : à la radio, l'auditeur ne peut pas revenir en arrière ; seuls les faits essentiels et l'appel doivent être retenus, dans des phrases courtes.
\end{corrige}

\newpage

\coursec{Méthodes et exercices résolus}

\coursec{Contraction de texte : reformuler les idées essentielles en vue d'un résumé}

\courssub{La structure de la phrase : outil de la contraction}

\begin{methode}[Analyser la structure syntaxique pour réduire le texte]
    \begin{enumerate}
        \item \textbf{Identifier le noyau de la phrase} : repérer le verbe (ou le prédicat) et son sujet. C'est l'information indispensable.
        \item \textbf{Repérer les expansions détachables} : circonstants de temps, de lieu, de manière ; compléments du nom longs ; propositions relatives explicatives ; incises. Ce sont les cibles prioritaires de la suppression.
        \item \textbf{Transformer les subordonnées en groupes nominaux} : 
            \begin{itemize}
                \item Proposition relative $\rightarrow$ adjectif ou participe passé (ex: \emph{le livre que j'ai lu} $\rightarrow$ \emph{le livre lu}).
                \item Proposition conjonctive (cause, but, conséquence) $\rightarrow$ préposition + infinitif ou nom (ex: \emph{parce qu'il pleut} $\rightarrow$ \emph{à cause de la pluie}).
            \end{itemize}
        \item \textbf{Regrouper les coordonnées} : fusionner plusieurs phrases simples ou coordonnées en une phrase complexe unique si le lien logique le permet.
        \item \textbf{Vérifier la cohérence syntaxique} : accord, ponctuation, ordre des mots après coupes.
    \end{enumerate}
\end{methode}

\begin{exoresolu}[Réduction syntaxique d'un paragraphe argumentatif]
\textbf{Énoncé :} Réduisez le paragraphe suivant à 40 mots environ (± 10\%), en conservant les idées forces et la progression logique. Comptez les mots.
\begin{quote}
« L'éducation nationale, qui est un pilier fondamental de toute société démocratique, doit impérativement s'adapter aux mutations technologiques rapides. En effet, les élèves, qui sont nés à l'ère du numérique, ne peuvent plus se contenter des méthodes pédagogiques traditionnelles, basées sur la transmission passive du savoir. Par conséquent, les enseignants sont contraints de repenser leurs pratiques pour intégrer les outils numériques, ce qui demande une formation continue soutenue. »
\end{quote}
\tcblower
\textbf{Solution détaillée :}
\begin{enumerate}
    \item \textbf{Analyse structurelle :}
    \begin{itemize}
        \item P1 : \emph{L'éducation nationale [qui est un pilier...] doit s'adapter...} $\rightarrow$ Sujet : L'éducation nationale. Verbe : doit s'adapter. Relative explicative supprimable.
        \item P2 : \emph{En effet, les élèves [qui sont nés...] ne peuvent plus se contenter...} $\rightarrow$ Cause. Relative explicative supprimable. \emph{Basées sur...} = participe passé (déjà réduit).
        \item P3 : \emph{Par conséquent, les enseignants sont contraints de repenser... [ce qui demande...]} $\rightarrow$ Conséquence. Relative de conséquence $\rightarrow$ préposition + nom.
    \end{itemize}
    \item \textbf{Application des transformations :}
    \begin{itemize}
        \item \emph{L'éducation nationale, pilier de toute société démocratique, doit s'adapter aux mutations technologiques.} (Suppression relative, apposition).
        \item \emph{Les élèves, nés à l'ère numérique, rejettent les méthodes traditionnelles passives.} (Relative $\rightarrow$ participe ; \emph{ne peuvent plus se contenter de} $\rightarrow$ \emph{rejettent} ; \emph{basées sur...} gardé).
        \item \emph{D'où la nécessité pour les enseignants de repenser leurs pratiques via l'intégration d'outils numériques, exigeant une formation continue.} (Conséquence fusionnée ; \emph{ce qui demande} $\rightarrow$ \emph{exigeant}).
    \end{itemize}
    \item \textbf{Assemblage et comptage :}
    \emph{« L'éducation nationale, pilier de toute société démocratique, doit s'adapter aux mutations technologiques. Les élèves, nés à l'ère numérique, rejettent les méthodes traditionnelles passives. D'où la nécessité pour les enseignants de repenser leurs pratiques via l'intégration d'outils numériques, exigeant une formation continue. »}
    \item \textbf{Comptage :} 42 mots (dans la fourchette 36--44). \textbf{Conforme.}
\end{enumerate}
\textbf{Conclusion :} La maîtrise de la syntaxe (apposition, participe, nominalisation) permet de condenser sans trahir la logique argumentative.
\end{exoresolu}

\courssub{Énoncés constatifs et performatifs}

\begin{methode}[Distinguier constatif et performatif pour filtrer l'essentiel]
    \begin{enumerate}
        \item \textbf{Repérer le verbe performatif} (explicite : \emph{je promets, j'ordonne, je déclare, je nomme} ; ou implicite/indirect).
        \item \textbf{Tester la vérité} : Un constatif est \emph{vrai} ou \emph{faux}. Un performatif n'est ni vrai ni faux, il est \emph{heureux} (réussi) ou \emph{inheureux} (raté : conditions non remplies).
        \item \textbf{Évaluer l'apport informationnel} :
            \begin{itemize}
                \item Constatif $\rightarrow$ Apporte une connaissance sur le monde $\rightarrow$ \textbf{À résumer} (souvent idée force).
                \item Performatif $\rightarrow$ Réalise une action (engagement, nomination, ordre) $\rightarrow$ Souvent \textbf{secondaire dans un résumé} (sauf si l'acte \emph{est} le sujet du texte : ex. discours de mariage, loi).
            \end{itemize}
        \item \textbf{Reformuler l'acte de langage} si nécessaire : \emph{« Je vous déclare mari et femme »} $\rightarrow$ \emph{Le maire les unit.} (Passage du performatif au constatif rapporté).
    \end{enumerate}
\end{methode}

\begin{exoresolu}[Filtrage constatif/performatif dans un texte administratif]
\textbf{Énoncé :} Résumez le texte suivant en 30 mots max. Identifiez au préalable les énoncés performatifs et justifiez leur traitement (conservation/suppression/reformulation).
\begin{quote}
« Le Ministre de l'Éducation de Base, vu la Constitution, déclare ouverts les concours d'entrée en 6ème pour la session 2025. Il invite les candidats à se présenter munis de leur convocation. Nous souhaitons pleine réussite à tous les postulants. »
\end{quote}
\tcblower
\textbf{Solution détaillée :}
\begin{enumerate}
    \item \textbf{Analyse des énoncés :}
    \begin{itemize}
        \item \emph{« déclare ouverts les concours... »} : Verbe \emph{déclarer} à la 1ère pers. indicatif présent $\rightarrow$ \textbf{Performatif explicite} (acte d'ouverture officiel). Condition : qualité du locuteur (Ministre) remplie.
        \item \emph{« invite les candidats... »} : Verbe \emph{inviter} $\rightarrow$ \textbf{Performatif} (acte directif / invitation).
        \item \emph{« Nous souhaitons pleine réussite... »} : Verbe \emph{souhaiter} $\rightarrow$ \textbf{Performatif} (acte expressif / vœu). Aucune valeur informative factuelle.
    \end{itemize}
    \item \textbf{Stratégie de résumé :} Le texte \emph{est} un acte administratif. L'information constative (le fait nouveau) est : \emph{ouverture des concours 2025, condition de convocation}. Les performatifs sont le \emph{moyen} de l'annonce, pas l'info nouvelle pour le lecteur (sauf l'autorité émettrice).
    \item \textbf{Rédaction :}
    \emph{« Le Ministre de l'Éducation de Base ouvre les concours d'entrée en 6ème session 2025 ; les candidats doivent présenter leur convocation. »}
    \item \textbf{Comptage :} 22 mots. \textbf{Conforme.}
    \item \textbf{Justification :} Les performatifs \emph{déclare, invite, souhaitons} ont été convertis en constatifs rapportés (\emph{ouvre, doivent présenter}) ou supprimés (vœu). L'autorité (Ministre) est conservée comme source.
\end{enumerate}
\textbf{Conclusion :} Dans un résumé, on rapporte l'effet du performatif (le fait accompli) plutôt que l'énonciation de l'acte lui-même, sauf analyse de discours.
\end{exoresolu}

\courssub{Ironie et paradoxe : figures de pensée}

\begin{methode}[Décoder l'ironie et le paradoxe pour ne pas les aplatir]
    \begin{enumerate}
        \item \textbf{Détecter l'ironie} : Décalage entre le sens littéral (positif) et le sens visé (négatif, critique). Indices : hyperboles, guillemets, ton emphatique, contexte critique.
        \item \textbf{Détecter le paradoxe} : Affirmation apparemment contradictoire mais révélant une vérité profonde (ex: \emph{« Je ne sais qu'une chose, c'est que je ne sais rien »}).
        \item \textbf{Règle d'or du résumé} : \textbf{Ne jamais garder le sens littéral seul.} Il faut restituer le \textbf{sens visé (ironique) ou la vérité profonde (paradoxe)}.
        \item \textbf{Techniques de reformulation} :
            \begin{itemize}
                \item Ironie $\rightarrow$ Mots-clés : \emph{critique, dénonce, moque, feint de louer, tourne en dérision}.
                \item Paradoxe $\rightarrow$ Structure : \emph{Contrairement aux apparences, [vérité profonde] / Si [A] semble vrai, c'est [B] qui l'est en réalité.}
            \end{itemize}
        \item \textbf{Vérifier} : Le résumé lu seul doit être compréhensible sans le texte original.
    \end{enumerate}
\end{methode}

\begin{exoresolu}[Rédaction d'un résumé intégrant une ironie soutenue]
\textbf{Énoncé :} Résumez le texte suivant en 35 mots ± 10\%. Le texte contient une ironie soutenue.
\begin{quote}
« Quelle admirable invention que la réunionite aigüe ! Nos dirigeants excellent dans l'art de convoquer des assemblées plénières pour décider... de la date de la prochaine réunion. Pendant ce temps, les dossiers s'empilent, la poussière les recouvre, et l'efficacité, cette grande oubliée, prend des vacances prolongées. Bravo pour cette gestions si « productive » ! »
\end{quote}
\tcblower
\textbf{Solution détaillée :}
\begin{enumerate}
    \item \textbf{Décodage :}
    \begin{itemize}
        \item \emph{« admirable invention », « excellent », « Bravo », « productive »} (guillemets) = \textbf{Ironie}. Sens visé : \emph{perte de temps, inefficacité, bureaucratie stérile}.
        \item \emph{« décider de la date de la prochaine réunion »} = \textbf{Paradoxe apparent} (réunion pour réunion) révélant la vacuité.
    \end{itemize}
    \item \textbf{Idées forces (sens visé) :}
    \begin{enumerate}
        \item Critique de la multiplication stérile des réunions (\emph{réunionite}).
        \item Conséquence : négligence du travail réel (dossiers en attente).
        \item Constat : inefficacité gestionnaire.
    \end{enumerate}
    \item \textbf{Rédaction (reformulation du sens visé) :}
    \emph{« L'auteur dénonce la "réunionite" : les dirigeants enchaînent des réunions improductives qui ne décident que de se revoir, laissant les dossiers réels en plan et l'efficacité au placard. »}
    \item \textbf{Comptage :} 33 mots (fourchette 31,5--38,5). \textbf{Conforme.}
    \item \textbf{Vérification :} L'ironie est explicitée (\emph{dénonce, improductives, laissant... en plan}). Le paradoxe est résolu (\emph{ne décident que de se revoir}).
\end{enumerate}
\textbf{Conclusion :} Résumer un texte ironique impose de \emph{traduire} le second degré en énoncé constatif clair et critique.
\end{exoresolu}

\courssub{Les variétés lexicales du français}

\begin{methode}[Uniformiser le registre lexical pour la reformulation]
    \begin{enumerate}
        \item \textbf{Identifier la variété source} : Français standard, soutenu, familier, argotique, régional (français camerounais, ivoirien, etc.), technique/jargon.
        \item \textbf{Déterminer la variété cible} : \textbf{Français standard} (langue de l'écrit scolaire, du résumé d'examen). Neutre, précis, sans marqueurs familiers/régionaux.
        \item \textbf{Opérer les substitutions lexicales} (Dictionnaire mental) :
            \begin{itemize}
                \item Familier $\rightarrow$ Standard : \emph{boulot $\rightarrow$ travail, bouffer $\rightarrow$ manger, fric $\rightarrow$ argent, gamin $\rightarrow$ enfant}.
                \item Régional (Cameroun) $\rightarrow$ Standard : \emph{appeler (au sens de téléphoner) $\rightarrow$ téléphoner/contacter ; le bureau $\rightarrow$ le poste de travail ; se débrouiller $\rightarrow$ faire face / gérer ; un taxi $\rightarrow$ un véhicule de transport en commun (si ambigu)}.
                \item Technique/Jargon $\rightarrow$ Standard (ou terme générique) : \emph{bug $\rightarrow$ dysfonctionnement/anomalie ; process $\rightarrow$ procédure/processus}.
            \end{itemize}
        \item \textbf{Attention aux faux amis / anglicismes} : \emph{Opportunité $\neq$ occasion (sens anglais) ; Réaliser $\neq$ comprendre ; Actuellement $\neq$ en ce moment (sens anglais).}
        \item \textbf{Privilégier l'hyponymie (généralisation)} : \emph{Voiture, moto, camion $\rightarrow$ véhicules ; Citron, orange, mangue $\rightarrow$ agrumes/fruits}.
    \end{enumerate}
\end{methode}

\begin{exoresolu}[Reformulation standardisée d'un extrait oral/camerounais]
\textbf{Énoncé :} Reformulez en français standard (résumé 25 mots max) l'extrait suivant, relevé dans une enquête de terrain.
\begin{quote}
« Mon papa, il a appelé le boss pour lui dire que le chantier, ça coince grave. Les gars n'ont pas le matos, du coup c'est le bazar total. Il faut qu'il se débrouille pour envoyer les sous vite fait. »
\end{quote}
\tcblower
\textbf{Solution détaillée :}
\begin{enumerate}
    \item \textbf{Repérage des variétés (Familier / Régional Cameroun / Oral) :}
    \begin{center}
    \begin{tabularx}{\linewidth}{|l|X|X|}\hline
    \textbf{Original} & \textbf{Variété} & \textbf{Standard (Cible)} \\ \hline
    Mon papa & Familier / Oral & Mon père \\ \hline
    A appelé (au sens téléphoner) & Régional (Cameroun) & A téléphoné / contacté \\ \hline
    Le boss & Familier / Arg. métier & Le chef de chantier / l'entrepreneur / le responsable \\ \hline
    Ça coince grave & Familier / Intensif. oral & Est bloqué / rencontre de graves difficultés \\ \hline
    Les gars & Familier & Les ouvriers / les travailleurs \\ \hline
    Le matos & Arg. métier (matériel) & Le matériel / les équipements \\ \hline
    C'est le bazar total & Familier / Imagé & C'est le désordre total / la situation est chaotique \\ \hline
    Se débrouiller & Régional (Cameroun) / Standard étendu & Prendre les dispositions / Agir / Résoudre \\ \hline
    Envoyer les sous & Familier (sous = argent) & Verser les fonds / Payer / Débloquer le financement \\ \hline
    Vite fait & Oral / Familier & Rapidement / Dans les meilleurs délais \\ \hline
    \end{tabularx}
    \end{center}
    \item \textbf{Synthèse des idées forces :} Père alerte responsable $\rightarrow$ chantier bloqué (manque matériel) $\rightarrow$ urgence paiement.
    \item \textbf{Rédaction standard :}
    \emph{« Mon père a contacté le responsable : le chantier est bloqué par manque de matériel, causant un désordre total. Il lui demande de débloquer les fonds d'urgence. »}
    \item \textbf{Comptage :} 25 mots. \textbf{Parfait.}
\end{enumerate}
\textbf{Conclusion :} La reformulation en français standard est une exigence de l'épreuve ; elle gomme l'oralité et les particularismes pour garantir l'intelligibilité universelle.
\end{exoresolu}

\courssub{Dégager la structure argumentative du texte}

\begin{methode}[Cartographier l'argumentation pour structurer le résumé]
    \begin{enumerate}
        \item \textbf{Segmenter le texte} : Découper en paragraphes logiques (unité de thème/argument). Numéroter les segments.
        \item \textbf{Identifier le type de structure} :
            \begin{itemize}
                \item \textbf{Déductive} : Thèse $\rightarrow$ Arguments $\rightarrow$ Conclusion.
                \item \textbf{Inductive} : Faits/Exemples $\rightarrow$ Généralisation/Thèse.
                \item \textbf{Dialectique} : Thèse $\rightarrow$ Antithèse $\rightarrow$ Synthèse.
                \item \textbf{Analytique} : Problème $\rightarrow$ Causes $\rightarrow$ Conséquences $\rightarrow$ Solutions.
            \end{itemize}
        \item \textbf{Repérer les connecteurs logiques} (marqueurs de relation) : \emph{donc, car, mais, cependant, en effet, par conséquent, d'une part... d'autre part}. Ils dessinent le squelette.
        \item \textbf{Construire le « squelette » du résumé} : Une phrase (ou groupe) par segment argumentatif majeur, reliée par les connecteurs logiques conservés ou rétablis.
        \item \textbf{Élaguer les illustrations} : Exemples, anecdotes, citations, statistiques détaillées $\rightarrow$ Remplacer par \emph{« par exemple »} ou généraliser (\emph{« de nombreux cas »}).
    \end{enumerate}
\end{methode}

\begin{exoresolu}[Analyse structurelle et résumé d'un texte dialectique]
\textbf{Énoncé :} Dégagez la structure argumentative du texte ci-dessous, puis rédigez un résumé de 50 mots ± 10\%.
\begin{quote}
« L'intelligence artificielle (IA) fascine par ses performances. \textbf{D'une part}, elle révolutionne la médecine : diagnostic précoce des cancers, chirurgie assistée par robot. \textbf{D'autre part}, elle inquiète : suppression massive d'emplois, biais algorithmiques discriminants, menace sur la vie privée. \textbf{Cependant}, interdire l'IA serait une erreur historique. \textbf{Il faut} l'encadrer par une éthique rigoureuse et une régulation juridique internationale, seule garantie d'un progrès partagé. »
\end{quote}
\tcblower
\textbf{Solution détaillée :}
\begin{enumerate}
    \item \textbf{Segmentation \& Structure (Dialectique) :}
    \begin{itemize}
        \item [\textcircled{1}] \textbf{Thèse / Exorde} : L'IA fascine (performances).
        \item [\textcircled{2}] \textbf{Argument 1 (Pour)} : Révolution médicale (exemples précis).
        \item [\textcircled{3}] \textbf{Argument 2 (Contre / Antithèse)} : Risques sociaux/éthiques (emplois, biais, vie privée). Marqueur : \emph{D'une part / D'autre part}.
        \item [\textcircled{4}] \textbf{Synthèse / Position de l'auteur} : Refus de l'interdiction ($\leftarrow$ \emph{Cependant}). Proposition : Encadrement éthique/juridique international. Marqueur : \emph{Il faut}.
    \end{itemize}
    \item \textbf{Élague :} Exemples médicaux (cancers, robot) $\rightarrow$ \emph{révolutionne la médecine}. Liste risques $\rightarrow$ \emph{menaces sociales et éthiques}.
    \item \textbf{Rédaction suivant le squelette :}
    \emph{« Si l'IA fascine par ses performances et révolutionne la médecine, elle menace l'emploi, l'équité et la vie privée. L'interdire serait une erreur : l'auteur préconise un encadrement éthique et juridique international pour garantir un progrès partagé. »}
    \item \textbf{Comptage :} 41 mots. \textbf{Attention :} Fourchette 45--55. Trop court.
    \item \textbf{Ajustement (réintégrer densité) :}
    \emph{« L'IA fascine par ses performances, révolutionnant notamment la médecine. Mais elle suscite de vives inquiétudes : suppressions d'emplois, biais discriminatoires, atteinte à la vie privée. Refusant son interdiction, l'auteur réclame un encadrement éthique et juridique international, seule voie vers un progrès partagé. »}
    \item \textbf{Comptage final :} 48 mots. \textbf{Conforme.}
\end{enumerate}
\textbf{Conclusion :} La structure dialectique (Pour/Contre $\rightarrow$ Synthèse) dicte l'architecture du résumé : il faut rendre visible le mouvement de la pensée.
\end{exoresolu}

\courssub{Techniques de réduction et reformulation des idées essentielles}

\begin{methode}[Appliquer les 4 opérations fondamentales de la contraction]
    \begin{enumerate}
        \item \textbf{La suppression (Élagage) :} Supprimer les redondances, pléonasmes, exemples illustratifs, détails anecdotiques, adjectifs qualificatifs non essentiels, incises.
        \item \textbf{La généralisation (Nomination / Hyponymie) :} Remplacer une énumération ou une explication par un hyperonyme (terme plus général).
            \begin{itemize}
                \item Énumération $\rightarrow$ Hyperonyme : \emph{Chiens, chats, oiseaux $\rightarrow$ animaux domestiques}.
                \item Définition/Explication $\rightarrow$ Terme technique/précis : \emph{« L'organe qui pompe le sang » $\rightarrow$ Le cœur}.
            \end{itemize}
        \item \textbf{La nomination (Nominalisation) :} Transformer une proposition (verbe + sujet) en groupe nominal (GN).
            \begin{itemize}
                \item \emph{« Les prix augmentent » $\rightarrow$ L'augmentation des prix / La hausse des prix}.
                \item \emph{« Il a décidé de partir » $\rightarrow$ Sa décision de partir / Son départ}.
            \end{itemize}
        \item \textbf{L'intégration (Fusion) :} Regrouper plusieurs phrases simples en une phrase complexe (subordination, coordination, participe, gérondif).
            \begin{itemize}
                \item \emph{« Il pleut. Le match est annulé. » $\rightarrow$ « La pluie entraînant l'annulation du match... » / « À cause de la pluie, le match est annulé. »}
            \end{itemize}
        \item \textbf{Contrôle final} : Comptage mots, cohérence logique, fidélité au sens, correction syntaxique.
    \end{enumerate}
\end{methode}

\begin{exoresolu}[Application combinée des 4 techniques sur un texte dense]
\textbf{Énoncé :} Contractez le texte suivant à 30 mots exacts (tolérance 0). Appliquez explicitement les 4 techniques.
\begin{quote}
« Le changement climatique est une réalité scientifique avérée. Les températures moyennes augmentent régulièrement depuis le début de l'ère industrielle. Cette augmentation entraîne la fonte des glaciers et la montée des eaux. Par conséquent, les îles basses et les côtes sont menacées de submersion. Les populations qui y vivent devront être déplacées. Cela coûtera très cher aux États. »
\end{quote}
\tcblower
\textbf{Solution détaillée :}
\begin{enumerate}
    \item \textbf{Analyse par technique :}
    \begin{center}
    \begin{tabularx}{\linewidth}{|l|X|l|}\hline
    \textbf{Segment original} & \textbf{Technique(s) appliquée(s)} & \textbf{Résultat} \\ \hline
    \emph{Le changement climatique est une réalité scientifique avérée.} & Suppression (pléonasme \emph{réalité/avérée}, \emph{scientifique} implicite) & \emph{Le changement climatique est avéré.} \\ \hline
    \emph{Les températures moyennes augmentent régulièrement depuis le début de l'ère industrielle.} & Généralisation (\emph{régulièrement, depuis...} $\rightarrow$ \emph{s'accroissent}) + Nominalisation & \emph{L'augmentation continue des températures depuis l'ère industrielle...} \\ \hline
    \emph{Cette augmentation entraîne la fonte des glaciers et la montée des eaux.} & Intégration (cause/conséquence) + Nominalisation & \emph{...entraîne fonte glaciaire et montée des eaux...} \\ \hline
    \emph{Par conséquent, les îles basses et les côtes sont menacées de submersion.} & Généralisation (\emph{îles basses et côtes} $\rightarrow$ \emph{territoires littoraux}) & \emph{...menaçant de submersion les territoires littoraux.} \\ \hline
    \emph{Les populations qui y vivent devront être déplacées.} & Nominalisation (\emph{devront être déplacées} $\rightarrow$ \emph{déplacement des populations}) & \emph{...et le déplacement des populations.} \\ \hline
    \emph{Cela coûtera très cher aux États.} & Intégration (conséquence finale) + Vocabulaire précis (\emph{coûtera cher} $\rightarrow$ \emph{coût exorbitant}) & \emph{...au coût exorbitant pour les États.} \\ \hline
    \end{tabularx}
    \end{center}
    \item \textbf{Assemblage fluide :}
    \emph{« Le changement climatique, avéré, voit les températures s'accroître depuis l'ère industrielle, entraînant fonte glaciaire et montée des eaux, menaçant de submersion les territoires littoraux et imposant le déplacement des populations, au coût exorbitant pour les États. »}
    \item \textbf{Comptage rigoureux :}
    1-Le 2-changement 3-climatique, 4-avéré, 5-voit 6-les 7-températures 8-s'accroître 9-depuis 10-l'ère 11-industrielle, 12-entraînant 13-fonte 14-glaciaire 15-et 16-montée 17-des 18-eaux, 19-menaçant 20-de 21-submersion 22-les 23-territoires 24-littoraux 25-et 26-imposant 27-le 28-déplacement 29-des 30-populations, 31-au 32-coût 33-exorbitant 34-pour 35-les 36-États. = \textbf{36 mots}. \textbf{Hors tolérance (30 exact).}
    \item \textbf{Réduction drastique (re-nominalisation/intégration) :}
    \emph{« Avéré, le réchauffement climatique accélère depuis l'ère industrielle, causant fonte glaciaire, montée des eaux, submersion des côtes, déplacement des populations et coûts exorbitants pour les États. »}
    \item \textbf{Comptage final :}
    1-Avéré, 2-le 3-réchauffement 4-climatique 5-accélère 6-depuis 7-l'ère 8-industrielle, 9-causant 10-fonte 11-glaciaire, 12-montée 13-des 14-eaux, 15-submersion 16-des 17-côtes, 18-déplacement 19-des 20-populations 21-et 22-coûts 23-exorbitants 24-pour 25-les 26-États. = \textbf{26 mots}. \textbf{Encore trop court.}
    \item \textbf{Ajustement précis pour 30 mots :}
    \emph{« Le réchauffement climatique, avéré, s'accélère depuis l'ère industrielle, causant fonte glaciaire, montée des eaux, submersion côtière, déplacement forcé des populations et coûts exorbitants pour les États. »}
    \item \textbf{Comptage validé :} 30 mots exactement.
\end{enumerate}
\textbf{Conclusion :} L'atteinte d'une contrainte stricte (30 mots exacts) impose une maîtrise chirurgicale des 4 techniques, surtout la nomination et l'intégration.
\end{exoresolu}

\courssub{Application au \emph{Vieux nègre et la médaille} : lectures méthodiques 5 et 6, exposés}

\begin{methode}[Synthèse méthodique pour l'exposé et le résumé d'œuvre intégrale]
    \begin{enumerate}
        \item \textbf{Lecture méthodique n°5 (Structure narrative/argumentative) :}
            \begin{itemize}
                \item Découper le roman en \textbf{séquences narratives} (Situation initiale, Élément déclencheur, Péripéties, Climax, Dénouement, Situation finale).
                \item Identifier la \textbf{structure argumentative implicite} : La médaille comme symbole $\rightarrow$ Thèse coloniale (honneur) vs Antithèse vécue (humiliation/ironie) $\rightarrow$ Synthèse tragique (mort/révolte muette).
            \end{itemize}
        \item \textbf{Lecture méthodique n°6 (Figures de style / Variétés langagières) :}
            \begin{itemize}
                \item Repérer l'\textbf{ironie} (titre, discours officiel vs réalité), le \textbf{paradoxe} (médaillé mais misérable), l'\textbf{antiphrase}.
                \item Noter le \textbf{français oralisé / créolisé} du Vieux Nègre (variété régionale) vs le \textbf{français administratif standard} du colonisateur. $\rightarrow$ Enjeu de reformulation en standard pour le résumé.
            \end{itemize}
        \item \textbf{Préparation de l'Exposé n°1 (Analyse) :} Structurer en 3 parties : 1. Contexte/Sujet, 2. Analyse thématique (colonialisme, mémoire, ironie), 3. Analyse stylistique (langues, figures, structure).
        \item \textbf{Préparation de l'Exposé n°2 (Contraction/Résumé) :} Appliquer la \textbf{Méthode 6 (4 techniques)} sur la base du squelette argumentatif (Méthode 5). Produire un résumé \emph{standard} de 10\% du texte cible (souvent un extrait clé au Bac).
        \item \textbf{Compte-rendu / Remédiation} : S'auto-évaluer sur la grille : \emph{Fidélité / Concision / Correction / Structure / Reformulation lexicale}.
    \end{enumerate}
\end{methode}

\begin{exoresolu}[Résumé d'un extrait clé du \emph{Vieux nègre et la médaille} (Simulation Bac)]
\textbf{Énoncé :} Résumez en 40 mots ± 10\% l'extrait suivant (fin du roman, mort du Vieux Nègre).
\begin{quote}
« Le Vieux Nègre, lui, ne bronchait plus. Il était là, assis par terre, le dos contre le mur de la case, la médaille au cou, brillant au soleil couchant comme un miroir. Il ne respirait plus. La médaille, elle, vivait encore. Elle brillait de tous ses feux, indifférente, éternelle, insultante. Les fourmis commençaient déjà à grimper le long de la chaîne pour aller explorer le visage du mort. C'était fini. La grande farce coloniale avait eu son dénouement. »
\end{quote}
\tcblower
\textbf{Solution détaillée :}
\begin{enumerate}
    \item \textbf{Analyse rapide (Méthodes 2, 3, 4, 5) :}
    \begin{itemize}
        \item \textbf{Constatifs :} Mort du Vieux Nègre, médaille intacte/brillante, fourmis sur le corps, crépuscule.
        \item \textbf{Ironie/Paradoxe (Fig. pensée) :} \emph{« La médaille vivait encore », « insultante », « grande farce coloniale »}. Sens visé : La décoration honore le mort mais symbolise l'absurdité et la cruauté du système colonial qui survit à l'individu.
        \item \textbf{Variétés :} Standard littéraire (pas de difficulté majeure).
        \item \textbf{Structure :} Tableau final (mort/brillance) $\rightarrow$ Sens général (farce coloniale).
    \end{itemize}
    \item \textbf{Application techniques (Méthode 6) :}
    \begin{itemize}
        \item \emph{Suppression :} \emph{assis par terre, dos contre mur, soleil couchant, comme un miroir, ne respirait plus, tous ses feux, commençaient déjà à grimper, explorer le visage.}
        \item \emph{Généralisation :} \emph{Fourmis... chaîne... visage} $\rightarrow$ \emph{Fourmis envahissent le corps.}
        \item \emph{Nominalisation :} \emph{Il ne respirait plus} $\rightarrow$ \emph{Sa mort} ; \emph{La médaille brillait} $\rightarrow$ \emph{La brillance de la médaille}.
        \item \emph{Intégration :} Fusion mort / contraste médaille / sens final.
    \end{itemize}
    \item \textbf{Rédaction :}
    \emph{« Le Vieux Nègre meurt, sa médaille coloniale brillant toujours, indifférente et insultante. Les fourmis envahissent son corps. Ce contraste final dévoile l'absurdité cruelle de la "farce coloniale" qui survit à ses victimes. »}
    \item \textbf{Comptage :} 36 mots. (Fourchette 36--44). \textbf{Conforme.}
    \item \textbf{Vérification items Bac :}
    \begin{itemize}
        \item Idées forces : Mort, Médaille vivante (ironie), Fourmis (détail réaliste/symbolique), Jugement final (Farce). \checkmark
        \item Reformulation standard / ironie explicitée (\emph{indifférente et insultante, absurdité cruelle}). \checkmark
        \item Structure logique (Chronologique $\rightarrow$ Sens). \checkmark
    \end{itemize}
\end{enumerate}
\textbf{Conclusion :} Le résumé d'œuvre littéraire au Bac technique exige de \emph{nominaliser la narration} et d'\emph{expliciter l'implicite ironique} sans le trahir.
\end{exoresolu}

\begin{attention}[Les 5 erreurs qui coûtent des points au Bac]
    \begin{enumerate}
        \item \textbf{Garder le sens littéral de l'ironie/paradoxe} : Écrire \emph{« Il trouve la réunion admirable »} au lieu de \emph{« Il dénonce l'inutilité de la réunion »}. $\rightarrow$ \textbf{0 point sur l'idée force.}
        \item \textbf{Conserver le registre familier/régional} : Laisser \emph{« mon papa, le boss, le matos, ça coince grave »} dans la copie. $\rightarrow$ \textbf{Pénalité correction / inadaptation à la consigne (français standard).}
        \item \textbf{Oublier la structure argumentative (effet « liste de courses »)} : Aligner les idées sans connecteurs logiques (\emph{donc, mais, cependant, par conséquent}). $\rightarrow$ \textbf{Perte de points « Structure/Coherence » (souvent 2 à 3 pts).}
        \item \textbf{Dépasser la limite de mots (hors tolérance)} : Ne pas compter ou tronquer brutalement la fin (souvent la thèse/synthèse). $\rightarrow$ \textbf{Pénalité stricte (souvent 1 pt par tranche de 10\% excédent, ou note plafonnée).}
        \item \textbf{Confondre résumé et commentaire / avis personnel} : Ajouter \emph{« Je trouve ce texte très touchant », « L'auteur a raison »}. $\rightarrow$ \textbf{Hors sujet partiel, sanctionné heavily.}
    \end{enumerate}
\end{attention}

\newpage

\coursec{Exercices}

\courssub{Je vérifie mes connaissances}

\begin{exercice}[QCM : structure de la phrase et énoncés]
Parmi les propositions suivantes, une seule est exacte pour chaque question. Recopiez le numéro et la lettre correspondante.
\begin{enumerate}
    \item Dans la phrase « Le directeur a signé le décret », le groupe verbal (GV) est :
    \begin{itemize}
        \item[A.] a signé
        \item[B.] a signé le décret
        \item[C.] le décret
        \item[D.] Le directeur a signé
    \end{itemize}
    \item Un énoncé performatif se caractérise par le fait qu'il :
    \begin{itemize}
        \item[A.] décrit une réalité observable
        \item[B.] accomplit l'action qu'il énonce
        \item[C.] exprime un doute sur la véracité des faits
        \item[D.] utilise nécessairement l'ironie
    \end{itemize}
    \item La nominalisation consiste à :
    \begin{itemize}
        \item[A.] transformer un verbe en adverbe
        \item[B.] transformer un verbe ou une phrase en nom ou en groupe nominal
        \item[C.] remplacer un nom par un pronom
        \item[D.] mettre une phrase à la voix passive
    \end{itemize}
    \item Laquelle de ces phrases relève de l'ironie ?
    \begin{itemize}
        \item[A.] « Quelle belle journée ! » dit sous un soleil radieux.
        \item[B.] « Quelle belle journée ! » dit sous une pluie torrentielle.
        \item[C.] Il pleut beaucoup aujourd'hui.
        \item[D.] Est-ce qu'il pleut ?
    \end{itemize}
\end{enumerate}
\end{exercice}

\begin{exercice}[Vrai / Faux justifié : figures de pensée et variétés du français]
Indiquez si les affirmations suivantes sont vraies ou fausses. Justifiez brièvement chaque réponse.
\begin{enumerate}
    \item Le paradoxe est une affirmation en apparence contradictoire mais qui révèle une vérité profonde.
    \item Le français standard est la seule variété lexicale admise dans une contraction de texte au Baccalauréat.
    \item Dans \emph{Le Vieux nègre et la médaille}, l'ironie sert uniquement à faire rire le lecteur.
    \item Un énoncé constatif peut être vrai ou faux ; un énoncé performatif ne s'évalue pas en termes de vérité mais de réussite (félicité ou infélicité).
\end{enumerate}
\end{exercice}

\begin{exercice}[Définitions et notions clés]
Donnez une définition précise pour chacun des termes suivants, en lien avec la contraction de texte :
\begin{enumerate}
    \item Structure argumentative (thèse, arguments, connecteurs).
    \item Reformulation (par rapport à la citation ou à la paraphrase).
    \item Énoncé performatif explicite (donnez un exemple).
    \item Camerounisme (donnez deux exemples courants).
\end{enumerate}
\end{exercice}

\begin{exercice}[Identification directe d'énoncés]
Pour chacune des phrases ci-dessous, précisez s'il s'agit d'un énoncé \textbf{constatif} ou \textbf{performatif}. Si c'est un performatif, indiquez le verbe performatif.
\begin{enumerate}
    \item « Je déclare la séance ouverte. »
    \item « Le taux de chômage des jeunes atteint 15 \% dans la région. »
    \item « Je vous promets de rendre le livre demain. »
    \item « Le Vieux nègre attend sa médaille depuis des années. »
    \item « Nous condamnons fermement ces actes de violence. »
\end{enumerate}
\end{exercice}

\courssub{Je m'entraîne}

\begin{exercice}[Repérage dans \emph{Le Vieux nègre et la médaille} (extraits)]
Voici deux extraits courts inspirés de l'œuvre de Ferdinand Oyono. Pour chaque extrait :
\begin{enumerate}
    \item Soulignez \textbf{un énoncé performatif} et justifiez son identification.
    \item Repérez \textbf{un passage ironique} (ou relevant du paradoxe) et expliquez la cible de l'ironie.
\end{enumerate}
\textbf{Extrait 1 :} « Le commandant fit signe au vieux nègre de s'approcher : "Tu as bien servi la France, lui dit-il. Voici ta médaille." Meka tendit la main, tremblante d'émotion. Il touchait au but. Toute une vie de fidélité couronnée par ce petit morceau de métal. »

\textbf{Extrait 2 :} « Le prêtre, dans son sermon, avait longuement parlé de la charité chrétienne. Les Européens, au premier rang, approuvaient de la tête. Mais quand la quête passa, les porte-monnaie restèrent fermés. "Pauvre noir, pensait le prêtre, tu n'as pas encore compris que la charité, ici, ne traverse pas la barrière de la couleur." »
\end{exercice}

\begin{exercice}[Condensation par nominalisation]
Transformez chacune des phrases complexes suivantes en une phrase simple (ou une phrase à structure nominale dense) sans perte de sens essentiel, en utilisant la nominalisation et la suppression du superflu.
\begin{enumerate}
    \item « Le gouvernement a décidé de créer de nouveaux centres de formation professionnelle parce que le chômage des jeunes augmente. »
    \item « Bien que les entreprises demandent des diplômes, beaucoup de jeunes n'ont pas les moyens de payer les études. »
    \item « Les jeunes qui ont suivi une formation technique trouvent plus facilement du travail que ceux qui ont fait des études générales. »
    \item « Il est nécessaire que l'État investisse dans l'enseignement technique pour que l'économie se développe. »
    \item « Le Vieux nègre a attendu la médaille pendant longtemps et il a été déçu quand il a vu qu'elle ne changeait rien à sa misère. »
\end{enumerate}
\end{exercice}

\begin{exercice}[Registre de langue : du familier au standard]
Réécrivez le paragraphe suivant en \textbf{français standard} (langue de l'administration et de l'écrit scolaire), en remplaçant les camerounismes, anglicismes et tournures familières par leurs équivalents standards.
« Mon \textbf{grand}, le \textbf{boss} m'a dit que le \textbf{dossier} est \textbf{au courant}. Il faut qu'on \textbf{se débrouille} pour \textbf{caler} les \textbf{papiers} avant la \textbf{deadline}. Si on \textbf{gagne} pas le marché, on va \textbf{manger} la \textbf{galère}. C'est pas \textbf{la joie} en ce moment, les \textbf{sous} sont \textbf{raides}. »
\end{exercice}

\begin{exercice}[Analyse de la structure argumentative]
Lisez le court texte argumentatif ci-dessous. Puis :
\begin{enumerate}
    \item Identifiez la \textbf{thèse} défendue.
    \item Relevez les \textbf{trois arguments principaux} (en les numérotant).
    \item Indiquez les \textbf{connecteurs logiques} (ou opérateurs argumentatifs) qui relient ces arguments à la thèse.
\end{enumerate}
\textbf{Texte :} « L'insertion professionnelle des jeunes diplômés au Cameroun reste un défi majeur. \textbf{D'abord}, le système éducatif forme davantage aux métiers de bureau qu'aux besoins réels du marché de l'emploi. \textbf{Ensuite}, le secteur privé, principal pourvoyeur d'emplois, peine à absorber la masse croissante de sortants universitaires faute de croissance économique suffisante. \textbf{Enfin}, l'absence de mécanismes efficaces d'orientation et de stage empêche l'adéquation formation-emploi. \textbf{Par conséquent}, une réforme profonde de l'enseignement technique et professionnel s'impose comme une urgence nationale. »
\end{exercice}

\courssub{Je me prépare au Bac}

\begin{exercice}[Contraction de texte type Baccalauréat Technique — Sujet : Insertion professionnelle des jeunes]
\textbf{Consigne :} Vous ferez une contraction du texte ci-dessous en \textbf{120 mots (\textpm 10 \%)}, soit entre 108 et 132 mots. Votre résumé devra reprendre \textbf{toutes les idées essentielles} dans l'ordre du texte, en \textbf{reformulant} (pas de copier-coller), sans opinion personnelle, dans un français standard correct.

\textbf{Texte source (environ 480 mots) :}

L'avenir professionnel de la jeunesse camerounaise constitue aujourd'hui l'une des préoccupations centrales des pouvoirs publics et de la société civile. Chaque année, les universités et les grandes écoles déversent sur le marché du travail des dizaines de milliers de diplômés dont une proportion alarmante ne trouve pas d'emploi correspondant à leur qualification. Ce phénomène, loin d'être conjoncturel, révèle des dysfonctionnements structurels profonds qui appellent des réponses systémiques et non plus seulement palliatives.

En premier lieu, l'inadéquation entre l'offre de formation et la demande du marché du travail apparaît comme la cause première de ce chômage des diplômés. Le système éducatif, hérité d'un modèle colonial tourné vers l'administration, continue de privilégier les filières générales et théoriques au détriment de l'enseignement technique, professionnel et technologique. Les programmes, souvent obsolètes, n'intègrent pas assez les compétences numériques, entrepreneuriales et pratiques exigées par les entreprises modernes. Par ailleurs, les équipements des établissements techniques restent vétustes, limitant l'acquisition de gestes professionnels réels. Les jeunes sortent donc avec des diplômes qui ne garantissent pas leur employabilité.

En second lieu, la structure même de l'économie camerounaise freine l'absorption de cette main-d'œuvre qualifiée. Le secteur formel privé, qui devrait être le premier employeur, reste étroit, dominé par quelques grandes entreprises et une multitude de très petites entreprises (TPE) aux capacités d'embauche limitées. Le secteur public, jadis pourvoyeur d'emplois massifs, a gelé les recrutements sous l'effet des politiques d'ajustement structurel et de la maîtrise de la masse salariale. L'économie informelle, qui emploie près de 90 \% de la population active, offre des conditions de travail précaires, sans protection sociale ni perspective de carrière, reléguant les diplômés dans un sous-emploi dévalorisant.

Face à ce constat, plusieurs pistes de solution s'imposent avec urgence. La réforme curriculaire doit aligner les contenus de formation sur les besoins réels identifiés par une concertation régulière avec le monde professionnel : stages obligatoires, alternance, projets tutorés. L'État doit inciter fiscalement les entreprises qui accueillent des stagiaires ou recrutent de jeunes diplômés. Surtout, la promotion de l'entrepreneuriat des jeunes, par l'accès facilité au financement (fonds jeunes, microcrédit), l'accompagnement technique et la simplification administrative, doit devenir un pilier de la politique de l'emploi. Il ne s'agit plus seulement de former pour l'emploi, mais de former à la création d'emplois.

En définitive, résoudre la crise de l'insertion professionnelle exige une rupture avec les logiques d'attente. Elle commande une refonte courageuse de l'école, une dynamisation de l'appareil productif et une confiance réelle en la capacité d'initiative de la jeunesse. C'est à ce prix que le dividende démographique tant espéré deviendra une réalité pour le Cameroun.
\end{exercice}

\begin{methode}[Rappel : les étapes de la contraction]
\begin{enumerate}
    \item \textbf{Lire} le texte deux fois et dégager la thèse et le plan (ici : constat, deux causes, solutions, conclusion).
    \item \textbf{Surligner} les idées essentielles : une idée par paragraphe au minimum.
    \item \textbf{Reformuler} avec ses propres mots : nominalisations, synonymes, suppression des exemples et des répétitions.
    \item \textbf{Rédiger} en français standard, avec des connecteurs, sans avis personnel ni citation.
    \item \textbf{Compter} les mots et vérifier la fourchette imposée (108 à 132 mots).
\end{enumerate}
\end{methode}

\begin{exercice}[Question de synthèse type Probatoire — \emph{Le Vieux nègre et la médaille}]
\textbf{Sujet :} « Montrez, à partir de l'extrait étudié en classe (la remise de la médaille par le commandant), comment l'ironie de Ferdinand Oyono dénonce l'illusion de la reconnaissance coloniale. »
\textbf{Consignes :}
\begin{enumerate}
    \item Présentez le contexte de la scène et la thèse implicite de l'auteur.
    \item Analysez deux procédés ironiques précis (lexique, ton, situation, paradoxe) et leur effet de sens.
    \item Concluez sur la portée critique de cette scène vis-à-vis du système colonial.
\end{enumerate}
\textbf{Contraintes :} réponse organisée en trois paragraphes, arguments appuyés sur le texte, environ \textbf{15 lignes} au total, français standard soigné.
\end{exercice}

\begin{exercice}[Exercice d'amélioration : correction d'un résumé d'élève]
Voici une proposition de résumé du texte sur l'insertion professionnelle (exercice précédent) rédigée par un élève. Ce résumé comporte \textbf{cinq erreurs majeures} (recopie, dépassement de longueur, opinion personnelle, omission d'idée essentielle, maladresse de reformulation).
\textbf{Travail demandé :}
\begin{enumerate}
    \item Identifiez et numérotez les \textbf{cinq erreurs} en les localisant dans le texte de l'élève (citez le passage fautif).
    \item Pour chaque erreur, précisez la \textbf{règle ou la méthode} violée.
    \item Proposez une \textbf{version corrigée} de ce résumé (respectant la consigne de 120 mots \textpm 10 \%).
\end{enumerate}

\textbf{Résumé de l'élève (165 mots) :}
« L'avenir des jeunes au Cameroun est un grand problème. Le texte dit que l'école ne forme pas bien pour le travail. Les programmes sont vieux et il n'y a pas d'ordinateurs dans les lycées techniques. Moi je pense que c'est la faute du gouvernement qui ne donne pas d'argent. L'économie ne marche pas : le privé est petit et le public n'embauche plus. L'informel c'est 90 \% mais c'est la galère. Pour finir, l'auteur propose de changer les programmes, de faire des stages, de donner de l'argent aux jeunes pour créer des entreprises. C'est urgent. Si on ne fait rien, le dividende démographique sera raté. Il faut vraiment que l'école change et que les patrons embauchent. C'est une question de survie pour le pays. »
\end{exercice}



\newpage

\coursec{Corrigés des exercices}

\begin{corrige}[Exercice 1 — QCM : structure de la phrase et énoncés]
\begin{enumerate}
    \item \textbf{B.} Le groupe verbal (GV) comprend le verbe et tous les compléments qui lui sont rattachés : « a signé le décret ». « Le directeur » constitue le groupe sujet (GS) ; « le décret » seul n'est que le COD.
    \item \textbf{B.} L'énoncé performatif \emph{accomplit} l'action qu'il énonce : dire « je promets », c'est promettre. La proposition A décrit le constatif.
    \item \textbf{B.} La nominalisation transforme un verbe (« décider ») ou une proposition en nom (« la décision ») ou en groupe nominal ; elle est un outil majeur de condensation en contraction de texte.
    \item \textbf{B.} L'ironie repose sur l'antiphrase : dire le contraire de ce que l'on pense. L'écart entre l'éloge (« belle journée ») et la réalité (la pluie torrentielle) crée l'effet ironique. En A, l'énoncé est simplement constatif et sincère.
\end{enumerate}
\textbf{Points de vigilance :} ne pas confondre GV et verbe seul ; ne pas croire que tout énoncé exclamatif est ironique : l'ironie exige un écart perceptible entre le dit et le pensé.
\end{corrige}

\begin{corrige}[Exercice 2 — Vrai / Faux justifié]
\begin{enumerate}
    \item \textbf{Vrai.} Le paradoxe heurte l'opinion commune par une contradiction apparente, mais invite à découvrir une vérité plus profonde (exemple : « il faut reculer pour mieux sauter »).
    \item \textbf{Vrai.} La contraction est un exercice scolaire et administratif : elle exige le français standard, sans camerounismes, anglicismes ni tournures familières.
    \item \textbf{Faux.} Chez Ferdinand Oyono, l'ironie est une arme critique : elle dénonce l'hypocrisie, l'injustice et la dérision du système colonial ; elle ne vise pas seulement le rire.
    \item \textbf{Vrai.} « Il pleut » est vrai ou faux ; « je déclare la séance ouverte » n'est ni vrai ni faux : il réussit (félicité) si les conditions sont remplies, sinon il échoue (infélicité).
\end{enumerate}
\end{corrige}

\begin{corrige}[Exercice 3 — Définitions et notions clés]
\begin{enumerate}
    \item \textbf{Structure argumentative :} organisation d'un texte qui défend une prise de position. Elle comprend la \cle{thèse} (l'idée défendue), les \cle{arguments} (raisons qui la justifient), les exemples qui les illustrent, et les \cle{connecteurs logiques} (d'abord, ensuite, donc, cependant\ldots) qui articulent l'ensemble. Dégager cette structure est la première étape de la contraction.
    \item \textbf{Reformulation :} reprise des idées du texte avec ses propres mots (synonymes, nominalisations, changements de construction), sans rien ajouter ni retrancher d'essentiel. Elle se distingue de la \emph{citation} (reprise mot à mot, interdite dans un résumé) et de la \emph{paraphrase} (qui reprend tout sans réduire) : la reformulation dans la contraction est à la fois fidèle et condensée.
    \item \textbf{Énoncé performatif explicite :} énoncé qui accomplit l'action qu'il dit, grâce à un verbe performatif à la première personne du présent. Exemple : « Je vous promets de venir » (verbe : \emph{promettre}) ; « Je déclare la séance ouverte » (verbe : \emph{déclarer}).
    \item \textbf{Camerounisme :} mot ou tournure propre au français parlé au Cameroun, à éviter dans l'écrit scolaire. Exemples : « être au courant » pour « être en règle / informé » ; « caler un papier » pour « préparer un document » ; « le ndolé est au menu » (mot local intégré au français).
\end{enumerate}
\end{corrige}

\begin{corrige}[Exercice 4 — Identification directe d'énoncés]
\begin{enumerate}
    \item \textbf{Performatif} — verbe performatif : \emph{déclarer} (déclarer la séance ouverte, c'est l'ouvrir).
    \item \textbf{Constatif} — énoncé descriptif, vérifiable, vrai ou faux.
    \item \textbf{Performatif} — verbe performatif : \emph{promettre}.
    \item \textbf{Constatif} — description d'un fait narratif.
    \item \textbf{Performatif} — verbe performatif : \emph{condamner} (l'énonciation accomplit l'acte de condamnation).
\end{enumerate}
\textbf{Point de vigilance :} le performatif exige souvent la première personne et le présent ; une phrase au passé (« j'ai promis ») redevient un constatif, car elle rapporte l'acte au lieu de l'accomplir.
\end{corrige}

\begin{corrige}[Exercice 5 — Repérage dans \emph{Le Vieux nègre et la médaille}]
\textbf{Extrait 1.} Performatif : « Voici ta médaille » — la parole, accompagnée du geste du commandant, accomplit l'acte de remise. Ironie : « Toute une vie de fidélité couronnée par ce petit morceau de métal » — l'écart entre l'immensité du sacrifice (une vie entière) et l'insignifiance de la récompense (un « petit morceau de métal ») dénonce la dérision de la reconnaissance coloniale.

\textbf{Extrait 2.} Performatif : le sermon lui-même (parler de charité, c'est accomplir un acte de prédication) ; on retiendra surtout le paradoxe : les Européens approuvent la charité en paroles mais ferment leurs porte-monnaie. L'ironie vise l'hypocrisie des colons et du prêtre : la « charité chrétienne » proclamée s'arrête à « la barrière de la couleur », ce qui contredit le message universel du christianisme.
\end{corrige}

\begin{corrige}[Exercice 6 — Condensation par nominalisation (propositions de réponses)]
\begin{enumerate}
    \item Face à la hausse du chômage des jeunes, le gouvernement décide la création de centres de formation professionnelle.
    \item Malgré les diplômes exigés par les entreprises, le coût des études reste inaccessible à beaucoup de jeunes.
    \item La formation technique garantit une meilleure insertion professionnelle que les études générales.
    \item L'investissement de l'État dans l'enseignement technique est indispensable au développement économique.
    \item Après une longue attente, Meka découvre avec amertume l'inutilité de sa médaille face à sa misère.
\end{enumerate}
\textbf{Points de vigilance :} la nominalisation (« a décidé de créer » $\rightarrow$ « décide la création ») ne doit supprimer ni l'agent ni la cause ; vérifier que la phrase condensée garde le sens essentiel et reste grammaticalement correcte.
\end{corrige}

\begin{corrige}[Exercice 7 — Registre de langue : proposition de réécriture]
« Mon frère (ou : mon ami), le directeur m'a informé que le dossier est en règle. Nous devons nous organiser pour préparer les documents administratifs avant la date limite. Si nous n'obtenons pas le marché, nous connaîtrons de grandes difficultés. La situation est difficile en ce moment : les finances sont épuisées. »

\textbf{Correspondances à retenir :} \emph{boss} $\rightarrow$ directeur / chef ; \emph{au courant} $\rightarrow$ en règle, à jour ; \emph{se débrouiller} $\rightarrow$ s'organiser ; \emph{caler} $\rightarrow$ préparer, régler ; \emph{deadline} $\rightarrow$ date limite, échéance ; \emph{manger la galère} $\rightarrow$ connaître de grandes difficultés ; \emph{les sous sont raides} $\rightarrow$ les fonds sont épuisés, nous manquons d'argent.
\end{corrige}

\begin{corrige}[Exercice 8 — Analyse de la structure argumentative]
\begin{enumerate}
    \item \textbf{Thèse :} une réforme profonde de l'enseignement technique et professionnel s'impose comme une urgence nationale (dernière phrase, introduite par « par conséquent »).
    \item \textbf{Arguments :} (1) inadéquation entre la formation scolaire et les besoins réels du marché de l'emploi ; (2) faiblesse d'absorption du secteur privé, faute de croissance suffisante ; (3) absence de mécanismes efficaces d'orientation et de stage.
    \item \textbf{Connecteurs :} « d'abord », « ensuite », « enfin » (énumération des arguments) ; « par conséquent » (conclusion logique reliant les arguments à la thèse).
\end{enumerate}
\textbf{Point de vigilance :} la thèse n'est pas toujours au début ; ici le texte procède par induction (du constat vers la conclusion). Repérer le connecteur conclusif (« par conséquent », « donc », « c'est pourquoi ») aide à la localiser.
\end{corrige}

\begin{corrige}[Exercice 9 — Contraction type Baccalauréat Technique : proposition de résumé (environ 120 mots)]
« L'insertion professionnelle des jeunes diplômés camerounais est devenue une préoccupation nationale majeure, révélant des dysfonctionnements structurels. D'abord, la formation dispensée, trop théorique et inadaptée aux besoins des entreprises, ne garantit pas l'employabilité des sortants. Ensuite, l'économie peine à absorber cette main-d'œuvre : le secteur privé formel reste étroit, le recrutement public est gelé et le secteur informel, largement dominant, n'offre que des emplois précaires. Pour y remédier, l'auteur préconise une réforme des programmes en concertation avec le monde professionnel, des incitations fiscales à l'embauche et surtout la promotion de l'entrepreneuriat des jeunes par le financement et l'accompagnement. Seule cette rupture permettra au Cameroun de transformer sa jeunesse en véritable atout démographique. » (123 mots)

\textbf{Barème indicatif (sur 10) :}
\begin{center}
\begin{tabularx}{\linewidth}{|X|c|}
\hline
\rowcolor{popblueL} \textbf{Critère} & \textbf{Points} \\
\hline
Respect du nombre de mots (108 à 132) & 1 \\
\hline
Présence de toutes les idées essentielles (constat, deux causes, solutions, conclusion) & 3 \\
\hline
Reformulation (aucune phrase recopiée) & 2 \\
\hline
Fidélité et objectivité (aucune opinion personnelle) & 2 \\
\hline
Qualité de la langue (français standard, connecteurs) & 2 \\
\hline
\end{tabularx}
\end{center}

\textbf{Points de vigilance :} compter les mots avant de rendre ; ne pas recopier les exemples chiffrés (les 90 \% peuvent être résumés par « largement dominant ») ; conserver l'ordre des idées du texte ; aucune phrase d'ouverture personnelle du type « ce texte parle de\ldots ».
\end{corrige}

\begin{corrige}[Exercice 10 — Question de synthèse type Probatoire : éléments de réponse]
\textbf{Paragraphe 1 — Contexte et thèse implicite.} Dans la scène de la remise de la médaille, Meka, vieux tirailleur fidèle, reçoit enfin du commandant la récompense promise pour ses services à la France. Oyono y dénonce implicitement l'illusion de la reconnaissance coloniale : l'administration flatte la loyauté des colonisés sans jamais les traiter en égaux.

\textbf{Paragraphe 2 — Deux procédés ironiques.} D'abord, le \emph{lexique dévalorisant} : la médaille est réduite à un « petit morceau de métal », alors qu'elle couronne « toute une vie de fidélité » ; l'écart entre le sacrifice et la récompense crée l'antiphrase. Ensuite, le \emph{paradoxe de la situation} : Meka, « tremblant d'émotion », croit toucher au but, alors que cette consécration ne changera rien à sa condition ; le lecteur, lui, perçoit la supercherie, ce qui instaure une complicité ironique contre le personnage du commandant et son paternalisme (« Tu as bien servi la France »).

\textbf{Paragraphe 3 — Portée critique.} La scène révèle que la médaille n'est qu'un instrument de domination symbolique : elle récompense la soumission et entretient l'illusion de la gratitude coloniale. Par l'ironie, Oyono invite le lecteur à démasquer l'hypocrisie du système et à mesurer l'aliénation de ceux qui ont cru à ses promesses.

\textbf{Barème indicatif (sur 10) :} contexte et thèse (2 pts) ; analyse de deux procédés précis appuyés sur le texte (4 pts) ; conclusion sur la portée critique (2 pts) ; organisation en trois paragraphes et qualité de la langue (2 pts).

\textbf{Points de vigilance :} citer brièvement le texte pour appuyer chaque analyse ; ne pas résumer toute l'œuvre ; distinguer clairement ironie (antiphrase) et paradoxe (contradiction apparente).
\end{corrige}

\begin{corrige}[Exercice 11 — Exercice d'amélioration : correction du résumé d'élève]
\textbf{Les cinq erreurs :}
\begin{enumerate}
    \item \textbf{Dépassement de longueur :} 165 mots au lieu de 108 à 132. Règle violée : respect du nombre de mots imposé (\textpm 10 \%).
    \item \textbf{Opinion personnelle :} « Moi je pense que c'est la faute du gouvernement ». Règle violée : le résumé doit rester objectif et fidèle à l'auteur, sans avis du résumant.
    \item \textbf{Recopie / reformulation insuffisante :} « le dividende démographique sera raté » reprend la fin du texte presque mot pour mot ; « changer les programmes, faire des stages » calque le texte. Règle violée : reformuler avec ses propres mots.
    \item \textbf{Omission et déformation d'idées essentielles :} l'incitation fiscale des entreprises et l'accompagnement de l'entrepreneuriat (microcrédit, simplification administrative) sont absents ; la vétusté des équipements est déformée en « pas d'ordinateurs ».
    \item \textbf{Registre familier et maladresses :} « c'est la galère », « les patrons », « l'économie ne marche pas ». Règle violée : usage obligatoire du français standard.
\end{enumerate}

\textbf{Version corrigée possible (environ 120 mots) :}

« L'insertion professionnelle des jeunes diplômés camerounais est devenue une préoccupation nationale majeure, révélant des dysfonctionnements structurels. D'abord, la formation dispensée, trop théorique et inadaptée aux besoins des entreprises, ne garantit pas l'employabilité des sortants. Ensuite, l'économie peine à absorber cette main-d'œuvre : le secteur privé formel reste étroit, le recrutement public est gelé et le secteur informel, largement dominant, n'offre que des emplois précaires. Pour y remédier, l'auteur préconise une réforme des programmes en concertation avec le monde professionnel, des incitations fiscales à l'embauche et surtout la promotion de l'entrepreneuriat des jeunes par le financement et l'accompagnement. Seule cette rupture permettra au Cameroun de transformer sa jeunesse en véritable atout démographique. »

\textbf{Point de vigilance :} corriger un résumé, c'est d'abord vérifier les cinq critères du cahier des charges : longueur, exhaustivité, reformulation, objectivité, langue standard.
\end{corrige}

\newpage

\coursec{Fiche bilan}

\begin{popfigure}
\begin{tikzpicture}[mindmap, grow cyclic, every node/.style={concept, text=white, font=\small},
  root concept/.append style={concept color=popblue, font=\bfseries\small, minimum size=3.1cm},
  level 1/.append style={level distance=3.4cm, sibling angle=51.4},
  level 1 concept/.append style={font=\scriptsize, minimum size=2.2cm}]
\node[root concept]{Contraction de texte : reformuler les idées essentielles}
  child[concept color=poporange]{ node {Structure de la phrase : sujet, verbe, compléments ; subordination} }
  child[concept color=popgreen]{ node {Énoncés constatifs / performatifs} }
  child[concept color=poppurple]{ node {Figures de pensée : ironie, paradoxe} }
  child[concept color=poppink]{ node {Variétés lexicales du français} }
  child[concept color=popgold]{ node {Structure argumentative : thèse, arguments, exemples} }
  child[concept color=popblue!70]{ node {Techniques de réduction : suppression, généralisation, fusion} }
  child[concept color=popmuted]{ node {Application : \emph{Le Vieux nègre et la médaille}} };
\end{tikzpicture}
\legende{Carte mentale de la leçon : les sept savoirs articulés autour de la contraction de texte.}
\end{popfigure}

\begin{bilan}
\textbf{1. Définitions clés}
\begin{itemize}
  \item \cle{Contraction de texte} : exercice qui consiste à réduire un texte à une fraction de sa longueur (souvent le quart) en conservant toutes les idées essentielles, sans commentaire personnel.
  \item \cle{Phrase simple} : une seule proposition (un verbe conjugué). \cle{Phrase complexe} : plusieurs propositions liées par juxtaposition, coordination ou subordination.
  \item \cle{Énoncé constatif} : décrit un fait, peut être vrai ou faux (\emph{Il pleut}). \cle{Énoncé performatif} : accomplit l'action en la disant (\emph{Je vous déclare mariés} ; verbes comme \emph{promettre, ordonner, jurer}).
  \item \cle{Ironie} : dire le contraire de ce qu'on pense pour railler. \cle{Paradoxe} : affirmation contraire à l'opinion commune mais qui peut être vraie.
  \item \cle{Variétés lexicales} : le français varie selon les régions (français du Cameroun, de Belgique, du Québec), les situations (registres familier, courant, soutenu) et les domaines (vocabulaires spécialisés).
\end{itemize}

\textbf{2. Structure argumentative à repérer}
\begin{center}
\begin{tabularx}{\linewidth}{|l|X|}\hline
\rowcolor{popblueL} \textbf{Élément} & \textbf{Repères} \\ \hline
Thèse (soutenue ou rejetée) & position de l'auteur, souvent en introduction ou conclusion \\ \hline
Arguments & raisons logiques marquées par \emph{car, en effet, d'abord, ensuite} \\ \hline
Exemples & illustrations marquées par \emph{ainsi, par exemple, c'est le cas de} \\ \hline
Connecteurs & hiérarchisent : \emph{mais, donc, cependant, enfin} \\ \hline
\end{tabularx}
\end{center}

\textbf{3. Techniques de réduction (méthode express)}
\begin{enumerate}
  \item \textbf{Supprimer} : exemples, répétitions, citations, détails descriptifs, parenthèses.
  \item \textbf{Généraliser} : remplacer une énumération par un terme générique (\emph{bananes, ignames, manioc} $\rightarrow$ \emph{les produits vivriers}).
  \item \textbf{Fusionner} : condenser plusieurs phrases en une seule grâce à la subordination ou à un nom résumant une phrase (\emph{il refusa la médaille} $\rightarrow$ \emph{son refus}).
  \item \textbf{Reformuler} : avec ses propres mots, au présent de vérité générale, sans copier les tournures de l'auteur.
\end{enumerate}

\textbf{4. Application au \emph{Vieux nègre et la médaille} (F. Oyono)} : l'ironie de Meka (la médaille, symbole de la « gratitude » coloniale, devient source d'humiliation) nourrit la thèse critique du roman ; le résumé doit en dégager les étapes sans raconter tous les épisodes.
\end{bilan}

\begin{aretenir}[Checklist avant le Bac]
\begin{itemize}
  \item[$\square$] Je sais définir la contraction de texte et la distinguer du résumé-commentaire.
  \item[$\square$] Je respecte la consigne de longueur (le quart du texte, marge de 10 \%).
  \item[$\square$] Je repère la thèse, les arguments et les exemples avant de réduire.
  \item[$\square$] Je souligne les connecteurs logiques pour suivre la progression du texte.
  \item[$\square$] Je supprime exemples, répétitions et citations sans perdre l'essentiel.
  \item[$\square$] Je généralise les énumérations par un terme générique.
  \item[$\square$] Je fusionne les phrases grâce à la subordination et à la nominalisation.
  \item[$\square$] Je reformule avec mes propres mots, sans copier-coller de longs passages.
  \item[$\square$] Je distingue un énoncé constatif d'un énoncé performatif.
  \item[$\square$] Je reconnais l'ironie et le paradoxe et j'explique leur effet.
  \item[$\square$] Je n'introduis aucune opinion personnelle dans mon résumé.
  \item[$\square$] Je rédige au présent, je relis et je compte mes mots.
\end{itemize}
\end{aretenir}

\findecours

\end{document}
